% La santé nutritionnelle : glycémie et pression artérielle — SVTEEHB Tle TI — Cours signé OnBuch+
% Programme officiel MINESEC. Compiler avec : tectonic N08-sante-nutritionnelle.tex
\newcommand{\DOCMATIERE}{SVTEEHB}
\newcommand{\DOCNIVEAU}{Tle TI}
\newcommand{\DOCMODULE}{Module II : L'Éducation à la Santé}
\newcommand{\DOCLECON}{N08}
\newcommand{\DOCTITRE}{La santé nutritionnelle : glycémie et pression artérielle}
% OnBuch+ — préambule « cours signés » ENRICHI (compilé avec tectonic / XeLaTeX)
% Système visuel « Pop » d'OnBuch+ : fond crème (jamais blanc), bordures encre
% épaisses, ombres « dures » décalées, titres Archivo Black en capitales.
% Version MATHÉMATIQUES : graphes (pgfplots/tikz), boîtes pédagogiques
% (définition, propriété, méthode, attention, le savais-tu, exercice résolu,
% exercice, TP, bilan), page de garde et signature OnBuch+ ancrée en bas.
%
% Macros attendues dans main.tex AVANT \input{preamble} :
%   \DOCMATIERE  \DOCNIVEAU  \DOCMODULE  \DOCLECON (numéro)  \DOCTITRE
\documentclass[11pt]{article}

\usepackage{fontspec}
\usepackage[french]{babel}
\usepackage[a4paper,margin=2cm,top=3.4cm,bottom=2.8cm,headheight=1.4cm,footskip=1.3cm]{geometry}
\usepackage{amsmath,amssymb,amsfonts}
\usepackage{mathtools} % \xmapsto, \coloneqq... souvent utilisés par les modèles
\usepackage{cancel} % simplifications barrées (bilans de forces)
% \abs{z} (module, valeur absolue) et \norm{u} (norme), utilisés par les modèles
\providecommand{\abs}{}\let\abs\relax\DeclarePairedDelimiter{\abs}{\lvert}{\rvert}
\providecommand{\norm}{}\let\norm\relax\DeclarePairedDelimiter{\norm}{\lVert}{\rVert}
\usepackage[table]{xcolor} % [table] : fournit aussi \rowcolors (alternance de lignes)
\usepackage{pagecolor}
\usepackage{tikz}
\usetikzlibrary{arrows.meta,positioning,calc,shapes.geometric,shapes.misc,shapes.arrows,decorations.pathmorphing,decorations.pathreplacing,decorations.markings,patterns,shadows,mindmap,trees,fit,backgrounds,matrix,intersections,angles,quotes}
\usepackage{pgfplots}
\usepackage[version=4]{mhchem} % équations nucléaires \ce{^{235}_{92}U}
\usepackage[european,siunitx]{circuitikz} % schémas électriques
\pgfplotsset{compat=1.18}
\usepgfplotslibrary{fillbetween,groupplots} % aires entre courbes (calcul intégral)
\usepackage{enumitem}
\usepackage{fancyhdr}
% [most] charge toutes les bibliothèques tcolorbox (listings, minted, raster,
% documentation...) et gonfle inutilement la mémoire TeX sur un document long
% (~300 boîtes) : on ne charge que ce qui est réellement utilisé.
\usepackage[skins,breakable]{tcolorbox}
\usepackage{graphicx}
\usepackage{colortbl}
\usepackage{booktabs}
\usepackage{tabularx}
\usepackage{multirow}
\usepackage{diagbox} % case d'en-tête diagonale des tableaux à double entrée
% Colonnes extensibles centrée / gauche / droite (C, L, R), souvent utilisées par les modèles
\newcolumntype{C}{>{\centering\arraybackslash}X}
\newcolumntype{L}{>{\raggedright\arraybackslash}X}
\newcolumntype{R}{>{\raggedleft\arraybackslash}X}
\usepackage{array}
\usepackage{multicol}
\usepackage{siunitx}
% parse-numbers=false : accepte aussi \SI{2,5 \times 10^{-3}}{...}
\sisetup{locale=FR,output-decimal-marker={,},group-minimum-digits=4,parse-numbers=false}
\DeclareSIUnit\ohm{\ensuremath{\symup{\Omega}}} % U+2126 absent de la police
\DeclareSIUnit\division{div} % réglages d'oscilloscope (ms/div, V/div)
\DeclareSIUnit\tr{tr} % tours (vitesse de rotation, ex. \SI{1500}{\tr\per\minute})
\DeclareSIUnit\molar{M} % concentration molaire (ex. \SI{3}{\molar}, \SI{1}{\micro\molar})
\DeclareSIUnit\an{an} % année (le modèle confond parfois avec \year, un compteur TeX) — ex. \SI{5}{\kilo\meter\per\an}
\DeclareSIUnit\FCFA{FCFA} % franc CFA (ex. \SI{500}{\FCFA\per\kilo\gram})
\DeclareSIUnit\degreeBrix{\textdegree Brix} % teneur en sucre d'un sirop/jus (réfractométrie)
\DeclareSIUnit\month{mois} % le modèle utilise parfois \month, un compteur TeX, comme unité de durée
\DeclareSIUnit\microliter{\micro\liter} % le modèle écrit parfois \microliter en un seul token
\DeclareSIUnit\million{millions} % dénombrement (ex. \SI{15}{\million\per\mL} spermatozoïdes)
\usepackage{qrcode}
% \degree : commande fréquemment utilisée par les modèles pour un angle ou une
% température en dehors de \SI{}{} (non fournie par les paquets chargés).
\providecommand{\degree}{^{\circ}}
% \qed (fin de démonstration), utilisé par les modèles sans amsthm
\providecommand{\qed}{\hfill$\square$}
% \barre : double barre des valeurs interdites dans un tableau de variations (tkz-tab)
\providecommand{\barre}{\ensuremath{\|}}
% environnement proof (amsthm non chargé) utilisé par les modèles
\ifdefined\proof\else\newenvironment{proof}[1][Démonstration]{\par\noindent\textit{#1.}\ }{\qed\par}\fi
\usepackage{lastpage}
\usepackage{hyperref}
\hypersetup{hidelinks}

% ============================================================
% Palette « Pop » OnBuch+ — mêmes hex que la classe P (plus_ui.dart)
% ============================================================
\definecolor{poporange}{HTML}{F59321}
\definecolor{popdark}{HTML}{E07A0C}
\definecolor{popcream}{HTML}{FFF6E8}
\definecolor{popink}{HTML}{1C1714}
\definecolor{poppurple}{HTML}{7A5AE0}
\definecolor{popgreen}{HTML}{2E9E6A}
\definecolor{poppink}{HTML}{E0457B}
\definecolor{popblue}{HTML}{2F7DE1}
\definecolor{popgold}{HTML}{C98A12}
\definecolor{popmuted}{HTML}{8A7F76}
\definecolor{popline}{HTML}{EADFD2}
\definecolor{popgray}{HTML}{8A7F76}
\colorlet{popgrey}{popgray}
% Alias tolérants (le modèle nomme parfois les couleurs différemment)
\colorlet{popwhite}{white}
\colorlet{popblack}{popink}
\colorlet{popred}{poppink}
\colorlet{popyellow}{popgold}
% Teintes claires (fonds de tableaux/encadrés) — le modèle nomme parfois
% les couleurs avec un suffixe L (ex. popblueL) sans les avoir définies.
\colorlet{poporangeL}{poporange!20}
\colorlet{popdarkL}{popdark!20}
\colorlet{poppurpleL}{poppurple!20}
\colorlet{popgreenL}{popgreen!20}
\colorlet{poppinkL}{poppink!20}
\colorlet{popblueL}{popblue!20}
\colorlet{popgoldL}{popgold!20}
\colorlet{popredL}{popred!20}
\colorlet{popgrayL}{popgray!20}
% Teintes claires pour fonds de tableaux / zones
\colorlet{popblueL}{popblue!10!white}
\colorlet{poporangeL}{poporange!14!white}
\colorlet{popgreenL}{popgreen!12!white}
\colorlet{poppurpleL}{poppurple!12!white}
\colorlet{poppinkL}{poppink!12!white}
\colorlet{popgoldL}{popgold!14!white}

\pagecolor{popcream}

% ============================================================
% Typographie
% ============================================================
\defaultfontfeatures{Ligatures=TeX}
\setmainfont{PlusJakartaSans-Medium.ttf}[Path=../../../fonts/,BoldFont=PlusJakartaSans-Bold.ttf]
% Maths en sans-serif (Fira Math) avec chiffres et lettres droites de Plus
% Jakarta, pour que les formules se fondent dans le texte.
\usepackage[math-style=ISO,bold-style=ISO]{unicode-math}
\setmathfont{FiraMath-Regular.otf}[Path=../../../fonts/]
\setmathfont{PlusJakartaSans-Medium.ttf}[Path=../../../fonts/,range={up/num,up/{Latin,latin},\mathup}]
\usepackage{icomma} % virgule décimale sans espace en mode math
% Caractères mathématiques Unicode absents des polices de texte (Plus Jakarta,
% Archivo Black) : rendus par leur équivalent mathématique, en texte comme en
% formule — sinon ils s'affichent en carré vide (ex. « Arithmétique dans ℤ »).
\usepackage{newunicodechar}
\newunicodechar{ℝ}{\ensuremath{\mathbb{R}}}
\newunicodechar{ℤ}{\ensuremath{\mathbb{Z}}}
\newunicodechar{ℂ}{\ensuremath{\mathbb{C}}}
\newunicodechar{ℕ}{\ensuremath{\mathbb{N}}}
\newunicodechar{ℚ}{\ensuremath{\mathbb{Q}}}
\newunicodechar{𝔻}{\ensuremath{\mathbb{D}}}
\newunicodechar{α}{\ensuremath{\alpha}}
\newunicodechar{β}{\ensuremath{\beta}}
\newunicodechar{γ}{\ensuremath{\gamma}}
\newunicodechar{δ}{\ensuremath{\delta}}
\newunicodechar{ε}{\ensuremath{\varepsilon}}
\newunicodechar{θ}{\ensuremath{\theta}}
\newunicodechar{λ}{\ensuremath{\lambda}}
\newunicodechar{μ}{\ensuremath{\mu}}
\newunicodechar{σ}{\ensuremath{\sigma}}
\newunicodechar{τ}{\ensuremath{\tau}}
\newunicodechar{φ}{\ensuremath{\varphi}}
\newunicodechar{ψ}{\ensuremath{\psi}}
\newunicodechar{ω}{\ensuremath{\omega}}
\newunicodechar{Σ}{\ensuremath{\Sigma}}
% majuscules produites par \MakeUppercase dans les titres (α-aminés -> Α)
\newunicodechar{Α}{\ensuremath{\alpha}}
\newunicodechar{Β}{\ensuremath{\beta}}
\newunicodechar{Γ}{\ensuremath{\Gamma}}
\newunicodechar{Δ}{\ensuremath{\Delta}}
\newunicodechar{Ε}{\ensuremath{\varepsilon}}
\newunicodechar{Θ}{\ensuremath{\Theta}}
\newunicodechar{Λ}{\ensuremath{\Lambda}}
\newunicodechar{Μ}{\ensuremath{\mu}}
\newunicodechar{Τ}{\ensuremath{\tau}}
\newunicodechar{Φ}{\ensuremath{\Phi}}
\newunicodechar{Ψ}{\ensuremath{\Psi}}
\newunicodechar{Ω}{\ensuremath{\Omega}}
\newunicodechar{↦}{\ensuremath{\mapsto}}
\newunicodechar{∈}{\ensuremath{\in}}
\newunicodechar{∉}{\ensuremath{\notin}}
\newunicodechar{∧}{\ensuremath{\wedge}}
\newunicodechar{⇔}{\ensuremath{\Leftrightarrow}}
\newunicodechar{⟺}{\ensuremath{\Longleftrightarrow}}
\newunicodechar{⇒}{\ensuremath{\Rightarrow}}
\newunicodechar{⟹}{\ensuremath{\Longrightarrow}}
\newunicodechar{∘}{\ensuremath{\circ}}
\newunicodechar{∪}{\ensuremath{\cup}}
\newunicodechar{∩}{\ensuremath{\cap}}
\newunicodechar{≡}{\ensuremath{\equiv}}
\newunicodechar{⊂}{\ensuremath{\subset}}
\newunicodechar{⊥}{\ensuremath{\perp}}
\newunicodechar{∃}{\ensuremath{\exists}}
\newunicodechar{∀}{\ensuremath{\forall}}
\newunicodechar{∛}{\ensuremath{\sqrt[3]{\,}}}
\newunicodechar{∜}{\ensuremath{\sqrt[4]{\,}}}
\newunicodechar{⃗}{\ensuremath{{}^{\to}}}
\newunicodechar{ⁿ}{\ifmmode^{n}\else\textsuperscript{\ensuremath{n}}\fi}
\newunicodechar{ˣ}{\ifmmode^{x}\else\textsuperscript{\ensuremath{x}}\fi}
\newunicodechar{ᵇ}{\ifmmode^{b}\else\textsuperscript{\ensuremath{b}}\fi}
\newunicodechar{ʳ}{\ifmmode^{r}\else\textsuperscript{\ensuremath{r}}\fi}
\newunicodechar{ᵘ}{\ifmmode^{u}\else\textsuperscript{\ensuremath{u}}\fi}
\newunicodechar{ʸ}{\ifmmode^{y}\else\textsuperscript{\ensuremath{y}}\fi}
\newunicodechar{ᵃ}{\ifmmode^{a}\else\textsuperscript{\ensuremath{a}}\fi}
\newunicodechar{ᵉ}{\ifmmode^{e}\else\textsuperscript{\ensuremath{e}}\fi}
\newunicodechar{ᵏ}{\ifmmode^{k}\else\textsuperscript{\ensuremath{k}}\fi}
\newunicodechar{ᵖ}{\ifmmode^{p}\else\textsuperscript{\ensuremath{p}}\fi}
\newunicodechar{ᴺ}{\ifmmode^{N}\else\textsuperscript{\ensuremath{N}}\fi}
\newunicodechar{ᶜ}{\ifmmode^{c}\else\textsuperscript{\ensuremath{c}}\fi}
\newunicodechar{ʰ}{\ifmmode^{h}\else\textsuperscript{\ensuremath{h}}\fi}
\newunicodechar{ᵠ}{\ifmmode^{q}\else\textsuperscript{\ensuremath{q}}\fi}
\newunicodechar{ₙ}{\ifmmode_{n}\else\textsubscript{\ensuremath{n}}\fi}
\newunicodechar{ₐ}{\ifmmode_{a}\else\textsubscript{\ensuremath{a}}\fi}
\newunicodechar{ᵢ}{\ifmmode_{i}\else\textsubscript{\ensuremath{i}}\fi}
\newunicodechar{ᵣ}{\ifmmode_{r}\else\textsubscript{\ensuremath{r}}\fi}
\newunicodechar{ₖ}{\ifmmode_{k}\else\textsubscript{\ensuremath{k}}\fi}
\newunicodechar{ᵦ}{\ifmmode_{\beta}\else\textsubscript{\ensuremath{\beta}}\fi}
\newunicodechar{ₓ}{\ifmmode_{x}\else\textsubscript{\ensuremath{x}}\fi}
\newfontfamily\popdisplay{ArchivoBlack-Regular.ttf}[Path=../../../fonts/]
\newfontfamily\poplogo{SpaceGrotesk-Bold.ttf}[Path=../../../fonts/]
\newfontfamily\popsemi{PlusJakartaSans-SemiBold.ttf}[Path=../../../fonts/]
\newfontfamily\popxbold{PlusJakartaSans-ExtraBold.ttf}[Path=../../../fonts/]
\newfontfamily\popxboldls{PlusJakartaSans-ExtraBold.ttf}[Path=../../../fonts/,LetterSpace=3.0]

\setlist[enumerate]{leftmargin=1.7em,itemsep=5pt,topsep=4pt}
\setlist[itemize]{leftmargin=1.7em,itemsep=4pt,topsep=4pt,label={\color{poporange}\textbullet}}
\setlength{\parindent}{0pt}
\setlength{\parskip}{5pt}
\renewcommand{\arraystretch}{1.25}

% pgfplots : style Pop par défaut
\pgfplotsset{
  every axis/.append style={axis line style={popink,line width=1pt},tick style={popink},
    grid style={popline,line width=0.5pt},label style={font=\small},tick label style={font=\footnotesize}},
  popcurve/.style={poporange,line width=1.8pt,no markers,smooth},
}
% Même style disponible dans un \draw TikZ simple (hors pgfplots)
\tikzset{popcurve/.style={poporange,line width=1.8pt}}
\tikzset{popgrid/.style={popline,very thin}}
\colorlet{popcurve}{poporange} % le nom du style est parfois utilisé comme couleur

% ============================================================
% Pill / squiggle
% ============================================================
\newcommand{\poppill}[3]{%
  \tikz[baseline=-0.55ex]\node[draw=popink,line width=1.2pt,rounded corners=2.6pt,%
    fill=#2,inner xsep=6.5pt,inner ysep=3pt]{%
    \popxboldls\fontsize{7}{7}\selectfont\color{#3}\MakeUppercase{#1}};%
}
\newcommand{\popsquiggle}[1]{%
  \tikz[baseline=-0.4ex]{
    \draw[#1,line width=2.6pt,line cap=round]
      plot [smooth] coordinates {(0,0.09) (0.34,0.01) (0.68,0.21) (1.02,0.01) (1.36,0.10)};
  }%
}
% Mot-clé surligné (marqueur orange sous le texte)
\newcommand{\cle}[1]{{\popxbold\color{popdark}#1}}

% ============================================================
% En-tête / pied
% ============================================================
\pagestyle{fancy}
\fancyhf{}
\renewcommand{\headrulewidth}{0pt}
\renewcommand{\footrulewidth}{0pt}
\lhead{\raisebox{-0.15ex}{\poplogo\fontsize{13}{13}\selectfont\color{popink}OnBuch\color{poporange}+}}
\rhead{\poppill{\DOCMATIERE\ \textbullet\ \DOCNIVEAU\ \textbullet\ Leçon \DOCLECON}{white}{popink}}
\lfoot{\popxbold\fontsize{7.6}{7.6}\selectfont\color{popmuted}\MakeUppercase{Cours signé OnBuch+}\ \textbullet\ {\popsemi\DOCTITRE}}
\rfoot{\tikz[baseline=-0.6ex]\node[rounded corners=8.5pt,fill=popink,minimum height=17pt,minimum width=17pt,inner xsep=5pt,inner sep=0]{\popxbold\fontsize{8}{8}\selectfont\color{popcream}\thepage\,/\,\pageref*{LastPage}};}

% ============================================================
% Titres
% ============================================================
\newcommand{\chaptitre}[2]{%
  \begin{tcolorbox}[enhanced,colback=poporange,colframe=popink,boxrule=2.6pt,arc=11pt,
      drop shadow=popink,left=15pt,right=15pt,top=12pt,bottom=11pt,before skip=3pt,after skip=18pt]
    \poppill{#2}{popink}{popcream}\par\vspace{9pt}
    {\raggedright\hyphenpenalty=10000\exhyphenpenalty=10000\popdisplay\fontsize{22}{24}\selectfont\color{popcream}\MakeUppercase{#1}\par}
  \end{tcolorbox}
}
% Section numérotée : pastille orange avec le numéro + titre Archivo
\newcounter{popsec}
\newcommand{\coursec}[1]{%
  \stepcounter{popsec}\setcounter{popsub}{0}%
  \par\vspace{16pt}\noindent
  \begin{minipage}{\linewidth}
  \tikz[baseline=-0.7ex]\node[draw=popink,line width=1.6pt,fill=poporange,rounded corners=4pt,minimum size=22pt,inner sep=0]{\popdisplay\fontsize{12}{12}\selectfont\color{popcream}\thepopsec};%
  \hspace{8pt}{\popdisplay\fontsize{15}{17}\selectfont\color{popink}\MakeUppercase{#1}}\par
  \vspace{4pt}\noindent{\color{poporange}\rule{36pt}{3.6pt}}
  \end{minipage}\par\nopagebreak\vspace{8pt}
}
\newcounter{popsub}
\newcommand{\courssub}[1]{%
  \stepcounter{popsub}%
  \par\vspace{9pt}\noindent{\popxbold\fontsize{11.5}{13}\selectfont\color{popink}{\color{poporange}\thepopsec.\thepopsub}\hspace{6pt}#1}\par\nopagebreak\vspace{4pt}
}

% ============================================================
% Boîtes pédagogiques — même recette : fond blanc, bordure épaisse colorée,
% ombre dure encre, badge plein. Titre optionnel [#1].
% ============================================================
\tcbset{popbase/.style={breakable,enhanced,colback=white,boxrule=2.3pt,arc=9pt,
  shadow={3pt}{-3pt}{0pt}{popink},left=14pt,right=14pt,top=11pt,bottom=11pt,before skip=11pt,after skip=15pt,
  fontupper=\popsemi\fontsize{10.6}{14.5}\selectfont\color{popink},
  fontlower=\popsemi\fontsize{10.6}{14.5}\selectfont\color{popink}}}
% Coche dessinée (le glyphe ✓ n'existe pas dans les polices du document)
\newcommand{\popcheck}[1]{\tikz[baseline=-0.5ex]\draw[#1,line width=1.6pt,line cap=round,line join=round](-0.13,0.0)--(-0.04,-0.09)--(0.13,0.1);}
\providecommand{\coche}{\popcheck{popgreen}}
\providecommand{\resultat}[1]{\par\smallskip\noindent\fcolorbox{popgreen}{popgreenL}{\parbox{\dimexpr\linewidth-2\fboxsep-2\fboxrule\relax}{\textbf{Résultat.} #1}}\par\smallskip}
\newcommand{\popboxhead}[3]{\poppill{#1}{#2}{white}\ifx\relax#3\relax\else\hspace{6pt}{\popxbold\fontsize{10.6}{12}\selectfont\color{popink}#3}\fi\par\vspace{7pt}}

\newtcolorbox{aretenir}[1][]{popbase,colframe=popgreen,before upper={\popboxhead{À retenir}{popgreen}{#1}}}
\newtcolorbox{exemplebox}[1][]{popbase,colframe=poporange,before upper={\popboxhead{Exemple}{poporange}{#1}}}
\newtcolorbox{definition}[1][]{popbase,colframe=popblue,before upper={\popboxhead{Définition}{popblue}{#1}}}
\newtcolorbox{propriete}[1][]{popbase,colframe=poppurple,before upper={\popboxhead{Propriété}{poppurple}{#1}}}
\newtcolorbox{methode}[1][]{popbase,colframe=popgold,before upper={\popboxhead{Méthode}{popgold}{#1}}}
\newtcolorbox{attention}[1][]{popbase,colframe=poppink,before upper={\popboxhead{Attention}{poppink}{#1}}}
\newtcolorbox{experience}[1][]{popbase,colframe=popblue,colback=popblueL,before upper={\popboxhead{Expérience}{popblue}{#1}}}
% « Le savais-tu ? » : carte violette pleine (contexte, culture, Cameroun)
\newtcolorbox{savaistu}[1][]{popbase,colback=poppurple,colframe=popink,
  fontupper=\popsemi\fontsize{10.4}{14.2}\selectfont\color{white},
  before upper={\poppill{Le savais-tu\,?}{white}{poppurple}\ifx\relax#1\relax\else\hspace{6pt}{\popxbold\color{white}#1}\fi\par\vspace{7pt}}}
% Exercice résolu : énoncé en haut, \tcblower puis la solution sur fond crème
\newcounter{popexo}\newcounter{popexr}
\newtcolorbox{exoresolu}[1][]{popbase,colframe=poporange,colbacklower=popcream,
  segmentation style={popink,line width=1.2pt,dashed},
  before upper={\stepcounter{popexr}\popboxhead{Exercice résolu \thepopexr}{poporange}{#1}},
  before lower={\poppill{Solution}{popink}{popcream}\par\vspace{6pt}}}
% Exercice (non résolu en place) — la correction est dans la partie Corrigés
\newtcolorbox{exercice}[1][]{popbase,colframe=popink,
  before upper={\stepcounter{popexo}\popboxhead{Exercice \thepopexo}{popink}{#1}}}
% Corrigé
\newtcolorbox{corrige}[1][]{popbase,colframe=popgreen,colback=popgreenL,
  before upper={\popboxhead{Corrigé}{popgreen}{#1}}}
% Objectifs / prérequis / bilan
\newtcolorbox{objectifs}{popbase,colframe=popink,colback=poporangeL,
  before upper={\popboxhead{Objectifs de la leçon}{popink}{}}}
\newtcolorbox{prerequis}{popbase,colframe=popink,colback=popblueL,
  before upper={\popboxhead{Prérequis}{popblue}{}}}
\newtcolorbox{bilan}{popbase,colframe=popink,colback=white,boxrule=2.8pt,
  before upper={\popboxhead{Fiche bilan}{poporange}{L'essentiel à connaître pour l'examen}}}
% Figure encadrée avec légende
\newtcolorbox{popfigure}[1][]{enhanced,breakable=false,colback=white,colframe=popline,boxrule=1.4pt,arc=8pt,
  left=10pt,right=10pt,top=10pt,bottom=8pt,before skip=10pt,after skip=12pt,halign=center,
  lower separated=false,halign lower=center,
  fontlower=\popsemi\fontsize{9}{11.5}\selectfont\color{popmuted},
  #1}
\newcommand{\legende}[1]{\par\vspace{6pt}{\popsemi\fontsize{9}{11.5}\selectfont\color{popmuted}#1}}

% ============================================================
% Signature OnBuch+
% ============================================================
\newcommand{\poplogoinline}[1]{{\poplogo\fontsize{#1}{#1}\selectfont\color{popink}OnBuch\color{poporange}+}}
\newcommand{\signatureonbuch}{%
  \begin{tcolorbox}[enhanced,colback=popink,colframe=popink,boxrule=2.2pt,arc=10pt,
      drop shadow=poporange,left=14pt,right=14pt,top=10pt,bottom=10pt,before skip=0pt,after skip=0pt]
    \tikz[baseline=-0.7ex]\node[circle,fill=popgreen,draw=popcream,line width=1.3pt,minimum size=22pt,inner sep=0]{\popcheck{white}};%
    \hspace{9pt}\begin{minipage}[c]{0.62\linewidth}
      {\popxbold\fontsize{12}{13}\selectfont\color{popcream}Signé\ \textbullet\ {\poplogo OnBuch\color{poporange}+}}\par\vspace{2pt}
      {\raggedright\popsemi\fontsize{8}{10.5}\selectfont\color{popline}\MakeUppercase{Équipe pédagogique OnBuch+}\\\MakeUppercase{Conforme au programme officiel MINESEC}\par}
    \end{minipage}\hfill
    {\poplogo\fontsize{22}{22}\selectfont\color{popcream}OnBuch\color{poporange}+}
  \end{tcolorbox}%
}

% ============================================================
% Page de garde
% #1 titre · #2 matière · #3 niveau · #4 module · #5 n° leçon · #6 durée ·
% #7 description / objectifs (texte ou itemize)
% ============================================================
\newcommand{\pagegarde}[7]{%
  \thispagestyle{empty}
  \newgeometry{a4paper,margin=1.7cm,top=1.6cm,bottom=1.4cm}%
  \begin{tikzpicture}[remember picture,overlay]
    \fill[poporange!18] ([xshift=-3.2cm,yshift=-3.6cm]current page.north east) circle (4.6cm);
    \fill[poppurple!14] ([xshift=2.4cm,yshift=7.6cm]current page.south west) circle (3.8cm);
  \end{tikzpicture}%
  {\poplogo\fontsize{30}{30}\selectfont\color{popink}OnBuch\color{poporange}+}\hfill
  \raisebox{0.6ex}{\poppill{Cours signé \textbullet\ Premium}{popink}{popcream}}\par
  \vspace{1pt}\popsquiggle{poppurple}\par
  \vspace{8pt}
  \begin{tcolorbox}[enhanced,colback=poporange,colframe=popink,boxrule=2.8pt,arc=13pt,
      drop shadow=popink,left=18pt,right=18pt,top=12pt,bottom=13pt,after skip=10pt]
    \poppill{#2 \textbullet\ #3}{popink}{popcream}\hspace{5pt}\poppill{#4}{white}{popink}\par\vspace{8pt}
    {\popdisplay\fontsize{14}{14}\selectfont\color{popink}LEÇON #5}\par\vspace{4pt}
    {\raggedright\hyphenpenalty=10000\exhyphenpenalty=10000\popdisplay\fontsize{25}{27}\selectfont\color{popcream}\MakeUppercase{#1}\par}
    \vspace{8pt}
    {\popxbold\fontsize{9.5}{9.5}\selectfont\color{popink}\MakeUppercase{Durée conseillée : #6}}
  \end{tcolorbox}
  \begin{tcolorbox}[enhanced,colback=white,colframe=popink,boxrule=2.2pt,arc=10pt,
      drop shadow=popink,left=16pt,right=16pt,top=10pt,bottom=10pt,after skip=8pt]
    \poppill{Dans cette leçon}{poppurple}{white}\par\vspace{5pt}
    \popsemi\fontsize{9.3}{12.6}\selectfont\color{popink}#7
  \end{tcolorbox}
  \noindent\begin{minipage}[t]{0.487\linewidth}
    \begin{tcolorbox}[enhanced,colback=popgreen,colframe=popink,boxrule=2.2pt,arc=10pt,
        drop shadow=popink,left=11pt,right=11pt,top=10pt,bottom=10pt,height=3.1cm,valign=center]
      \tcbox[colback=white,colframe=popink,boxrule=1.3pt,arc=5pt,boxsep=0pt,nobeforeafter,
        left=3pt,right=3pt,top=3pt,bottom=3pt,valign=center]{\qrcode[height=1.9cm]{https://play.google.com/store/apps/details?id=com.luvvix.onbuch&hl=fr}}%
      \hspace{8pt}\begin{minipage}[c]{0.5\linewidth}
        {\popdisplay\fontsize{11}{12.5}\selectfont\color{white}\MakeUppercase{Télécharge OnBuch}}\par\vspace{4pt}
        {\raggedright\popsemi\fontsize{8.4}{11}\selectfont\color{white}Gratuit \textbullet\ Google Play\\Tuteur IA, annales, concours, résultats\par}
      \end{minipage}
    \end{tcolorbox}
  \end{minipage}\hfill
  \begin{minipage}[t]{0.487\linewidth}
    \begin{tcolorbox}[enhanced,colback=poppurple,colframe=popink,boxrule=2.2pt,arc=10pt,
        drop shadow=popink,left=11pt,right=11pt,top=10pt,bottom=10pt,height=3.1cm,valign=center]
      \tikz[baseline=-0.7ex]\node[circle,fill=white,draw=popink,line width=1.3pt,minimum size=1.6cm,inner sep=0]{\popdisplay\fontsize{24}{24}\selectfont\color{poppurple}+};%
      \hspace{8pt}\begin{minipage}[c]{0.6\linewidth}
        {\popdisplay\fontsize{11}{12.5}\selectfont\color{white}\MakeUppercase{Passe à OnBuch+}}\par\vspace{4pt}
        {\raggedright\popsemi\fontsize{8.4}{11}\selectfont\color{white}Tous les cours signés, Tuteur IA illimité, fiches PDF — onglet OnBuch+ de l'app\par}
      \end{minipage}
    \end{tcolorbox}
  \end{minipage}
  \vspace{8pt}
  \popcontenu
  \vfill
  \signatureonbuch
  \restoregeometry
  \newpage
}

% Rangée « Ce cours contient » (page de garde)
\newcommand{\poptile}[3]{% #1 couleur #2 gros texte #3 légende
  \begin{tcolorbox}[enhanced,colback=#1,colframe=popink,boxrule=2pt,arc=9pt,shadow={2.5pt}{-2.5pt}{0pt}{popink},
      left=6pt,right=6pt,top=9pt,bottom=9pt,halign=center,valign=center,height=1.75cm,width=\linewidth,nobeforeafter]
    {\popdisplay\fontsize{12.5}{13}\selectfont\color{white}\MakeUppercase{#2}}\par\vspace{3pt}
    {\centering\popsemi\fontsize{7.8}{9.5}\selectfont\color{white}#3\par}
  \end{tcolorbox}}
\newcommand{\popcontenu}{%
  \par\noindent{\popxboldls\fontsize{7.5}{7.5}\selectfont\color{popmuted}CE COURS SIGNÉ CONTIENT}\par\vspace{7pt}
  \noindent\begin{minipage}[t]{0.235\linewidth}\poptile{popblue}{Cours}{complet, illustré, pas à pas}\end{minipage}\hfill
  \begin{minipage}[t]{0.235\linewidth}\poptile{poporange}{Méthodes}{et exercices résolus}\end{minipage}\hfill
  \begin{minipage}[t]{0.235\linewidth}\poptile{poppink}{Exercices}{type examen + corrigés}\end{minipage}\hfill
  \begin{minipage}[t]{0.235\linewidth}\poptile{popgreen}{Fiche bilan}{l'essentiel à retenir}\end{minipage}\par}

% Sommaire
\newtcolorbox{sommaire}{enhanced,colback=white,colframe=popink,boxrule=2.2pt,arc=10pt,
  drop shadow=popink,left=18pt,right=18pt,top=13pt,bottom=13pt,before skip=6pt,after skip=20pt,
  before upper={\poppill{Sommaire}{poppurple}{white}\par\vspace{9pt}\popsemi\fontsize{10.6}{14.5}\selectfont\color{popink}}}

% Signature de fin de cours — poussée tout en bas de la dernière page
\newcommand{\findecours}{%
  \par\vfill\nopagebreak
  \begin{center}\popsquiggle{poporange}\end{center}\vspace{4pt}
  \signatureonbuch
}
\AtBeginDocument{\def\llbracket{\mathopen{[\mkern-3.5mu[}}\def\rrbracket{\mathclose{]\mkern-3.5mu]}}} % intervalles d'entiers (glyphe ⟦ absent de la police)
% Glyphes absents de Fira Math : redéfinis à partir de glyphes disponibles
\AtBeginDocument{%
  \def\setminus{\mathbin{\backslash}}\let\smallsetminus\setminus
  \def\vdots{\mathord{\vbox{\baselineskip3.2pt\lineskiplimit0pt\kern4pt\hbox{.}\hbox{.}\hbox{.}}}}%
  \def\ddots{\mathinner{\mkern1mu\raise6.5pt\hbox{.}\mkern2mu\raise3.6pt\hbox{.}\mkern2mu\raise0.7pt\hbox{.}\mkern1mu}}%
  \def\triangle{\mathord{\scalebox{0.9}{$\symup{\Delta}$}}}%
  \def\nabla{\mathord{\raisebox{\depth}{\scalebox{1}[-1]{$\symup{\Delta}$}}}}%
  \def\checkmark{\popcheck{popdark}}%
  \def\star{\mathbin{\ast}}\def\bigstar{\mathord{\ast}}%
  \def\diamond{\mathbin{\scalebox{0.6}{$\Diamond$}}}%
  \def\bigcup{\mathop{\mathchoice{\raisebox{-0.3em}{\scalebox{1.7}{$\cup$}}}{\scalebox{1.25}{$\cup$}}{\cup}{\cup}}\displaylimits}%
  \def\bigcap{\mathop{\mathchoice{\raisebox{-0.3em}{\scalebox{1.7}{$\cap$}}}{\scalebox{1.25}{$\cap$}}{\cap}{\cap}}\displaylimits}%
  \def\lesseqqgtr{\mathrel{\lessgtr}}\def\bigoplus{\mathop{\oplus}\displaylimits}%
}
\providecommand{\ssi}{si et seulement si} % abréviation fréquente
\AtBeginDocument{\providecommand{\wideparen}{\overparen}} % arc de cercle

%% FIGFIX-BEGIN v5 — correctifs globaux des figures (docs/onbuch-generation/tools/figfix)
\usepackage{adjustbox}\usepackage{etoolbox}\tcbuselibrary{raster}
% toute figure TikZ/pgfplots plus large que la ligne est réduite à la largeur disponible
\BeforeBeginEnvironment{tikzpicture}{\begin{adjustbox}{max width=\linewidth}}
\AfterEndEnvironment{tikzpicture}{\end{adjustbox}}
% légendes pgfplots lisibles par-dessus les courbes
\pgfplotsset{every axis/.append style={legend style={fill=white,fill opacity=0.92,text opacity=1,draw=black!15,inner sep=3pt}}}
% étiquettes d'arêtes (syntaxe edge["..."]) avec fond blanc
\tikzset{every edge quotes/.append style={fill=white,inner sep=1.2pt}}
%% FIGFIX-END
\begin{document}

\pagegarde{La santé nutritionnelle : glycémie et pression artérielle}{SVTEEHB}{Tle TI}{Module II : L'Éducation à la Santé}{N08}{6 h}{Cette leçon étudie les mécanismes de régulation de la glycémie (rôles du pancréas et du foie, insuline et glucagon) et de la pression artérielle, établis à partir de faits expérimentaux. Elle en déduit les conséquences pathologiques des dérèglements permanents : diabètes, hypoglycémie, hypertension et hypotension artérielles, ainsi que les moyens de lutte adaptés au contexte camerounais.
\vspace{4pt}
\begin{multicols}{2}\small
\begin{itemize}
\item À la fin de cette leçon, je sais analyser et interpréter des courbes de variation de la glycémie avant et après un repas
\item identifier les facteurs de variation de la glycémie et de la pression artérielle
\item expliquer, à partir de faits expérimentaux, les rôles du pancréas et du foie dans la régulation de la glycémie
\item réaliser un schéma fonctionnel de la régulation hormonale de la glycémie
\item interpréter les résultats d'analyses médicales (glycémie, HGPO, glycosurie) dans le cas du diabète
\item distinguer les types de diabète, leurs causes et les moyens de lutte
\end{itemize}\end{multicols}}

\begin{sommaire}
\begin{multicols}{2}
\begin{enumerate}
\item La glycémie : définition, mesure et facteurs de variation
\item La régulation hormonale de la glycémie : rôles du pancréas et du foie
\item Le diabète et l'hypoglycémie : conséquences des dérèglements de la glycémie
\item La pression artérielle : définition et mesure
\item Variations de la pression artérielle, hypertension et hypotension
\item Activité d'intégration
\item Méthodes et exercices résolus
\item Exercices
\item Corrigés des exercices
\item Fiche bilan
\end{enumerate}
\end{multicols}
\end{sommaire}

\begin{objectifs}
\begin{itemize}
\item À la fin de cette leçon, je sais analyser et interpréter des courbes de variation de la glycémie avant et après un repas
\item identifier les facteurs de variation de la glycémie et de la pression artérielle
\item expliquer, à partir de faits expérimentaux, les rôles du pancréas et du foie dans la régulation de la glycémie
\item réaliser un schéma fonctionnel de la régulation hormonale de la glycémie
\item interpréter les résultats d'analyses médicales (glycémie, HGPO, glycosurie) dans le cas du diabète
\item distinguer les types de diabète, leurs causes et les moyens de lutte
\item expliquer la technique de mesure de la pression artérielle au sphygmomanomètre
\item analyser des documents montrant les conditions de variation de la pression artérielle
\item exploiter des résultats médicaux pour diagnostiquer une hypertension ou une hypotension et proposer des mesures de prévention
\end{itemize}
\end{objectifs}

\begin{prerequis}
\begin{itemize}
\item Notion de milieu intérieur et d'homéostasie (constance du milieu intérieur)
\item Rôle des hormones et mode d'action hormonale via le sang (notion de glande endocrine)
\item Anatomie générale du foie, du pancréas et de l'appareil circulatoire (cœur, artères)
\item Lecture et exploitation de courbes expérimentales (acquise en Première)
\end{itemize}
\end{prerequis}

\newpage

\coursec{La glycémie : définition, mesure et facteurs de variation}

Au marché Mokolo de Yaoundé, une vendeuse de 52 ans s'évanouit un matin après avoir sauté le petit-déjeuner. Conduite au centre de santé, on lui mesure une glycémie de \SI{0,45}{g/L} : elle fait une \cle{hypoglycémie}. Le même jour, son voisin de stand, habitué des plats très salés et gras, apprend que sa tension artérielle est de 17/11 : il souffre d'\cle{hypertension artérielle}. Ces deux situations, de plus en plus fréquentes au Cameroun où les maladies métaboliques (diabète, hypertension) progressent, posent une question fondamentale : \textbf{comment l'organisme maintient-il constants le taux de glucose sanguin et la pression artérielle, et que se passe-t-il quand cette régulation déraille ?} Cette leçon répond à ces questions en s'appuyant sur des expériences et des données médicales réelles. Commençons par le paramètre le plus étudié : la glycémie.

\courssub{La glycémie : définition et valeur normale}

\textbf{Observation préliminaire.} Tout le monde a déjà ressenti, après un effort intense ou un jeûne prolongé, une sensation de faiblesse, des tremblements, voire un malaise. À l'inverse, après un repas riche en féculents (fufu, riz, igname), on se sent plein d'énergie. Ces sensations suggèrent qu'une substance énergétique circule dans le sang et que sa quantité varie au cours de la journée. Cette substance, c'est le \cle{glucose}.

\begin{definition}[La glycémie]
La \cle{glycémie} est le \textbf{taux (concentration) de glucose dans le sang}. Elle s'exprime en grammes par litre de sang (\si{g/L}) ou en millimoles par litre (\si{mmol/L}). Chez un sujet sain, la glycémie \textbf{à jeun} (au moins 8 h sans repas) est d'environ \SI{1}{g/L}, la fourchette normale s'étendant de \SI{0,7}{g/L} à \SI{1,1}{g/L} (soit 3,9 à \SI{5,5}{mmol/L}).
\end{definition}

\textbf{Pourquoi le glucose est-il si important ?} Le glucose est la \textbf{principale source d'énergie des cellules}. Par la respiration cellulaire, chaque cellule oxyde le glucose pour produire l'ATP indispensable à son fonctionnement, selon le bilan global :
\[ \ce{C6H12O6 + 6O2 -> 6CO2 + 6H2O} + \text{énergie (ATP)} \]
Certaines cellules sont particulièrement dépendantes du glucose : les \textbf{cellules nerveuses} (neurones) n'utilisent quasiment que lui comme carburant et ne possèdent aucune réserve. C'est pourquoi une chute brutale de la glycémie se traduit d'abord par des signes nerveux (tremblements, sueurs, confusion, malaise, coma).

\begin{savaistu}[Le cerveau, grand consommateur de glucose]
Le cerveau ne représente que 2 \% de la masse corporelle, mais il consomme à lui seul environ \textbf{120 g de glucose par jour}, soit près de la moitié des besoins totaux de l'organisme au repos. Cela explique pourquoi une hypoglycémie provoque rapidement des troubles de la conscience : privé de son carburant exclusif, le cerveau dysfonctionne en quelques minutes. C'est exactement ce qui est arrivé à la vendeuse du marché Mokolo.
\end{savaistu}

\begin{attention}[Ne pas confondre glycémie et glycosurie]
Erreur fréquente au Baccalauréat : la \cle{glycémie} est le taux de glucose \textbf{dans le sang}, tandis que la \cle{glycosurie} est la présence de glucose \textbf{dans les urines}. Chez un sujet sain, la glycosurie est nulle : le rein réabsorbe totalement le glucose filtré tant que la glycémie reste inférieure à un seuil (environ \SI{1,8}{g/L}, appelé seuil rénal du glucose). La glycosurie n'apparaît donc que lorsque la glycémie dépasse durablement ce seuil, comme chez le diabétique.
\end{attention}

\courssub{Mise en évidence expérimentale : la glycémie varie au cours de la journée}

\begin{experience}[Mesures de la glycémie avant et après un repas]
\textbf{Protocole.} On prélève un peu de sang (au bout du doigt, à l'aide d'un lecteur de glycémie, ou par prélèvement veineux au laboratoire) chez un sujet sain :
\begin{itemize}
\item à jeun, le matin (t = 0 h) ;
\item puis toutes les 30 min après un repas riche en glucides (riz, pain, banane mûre) pris immédiatement après la mesure t = 0 h.
\end{itemize}
\textbf{Résultats.} Les mesures sont consignées dans le tableau ci-dessous et traduites par une courbe.
\begin{center}
\begin{tabularx}{\linewidth}{|l|*{7}{>{\centering\arraybackslash}X|}}
\hline
\rowcolor{popblueL} Temps (h) & 0 & 0,5 & 1 & 1,5 & 2 & 2,5 & 3 \\
\hline
Glycémie (g/L) & 0,95 & 1,25 & 1,40 & 1,20 & 1,05 & 0,95 & 0,92 \\
\hline
\end{tabularx}
\end{center}
\textbf{Interprétation.} La glycémie à jeun vaut environ \SI{0,95}{g/L}. Après le repas, elle s'élève jusqu'à un pic d'environ \SI{1,4}{g/L} vers 1 h (phase d'\cle{hyperglycémie post-prandiale}), puis décroît et revient à la valeur initiale en 2 à 3 h.
\textbf{Conclusion.} La glycémie n'est pas fixe : elle varie en fonction de l'alimentation, mais ces variations restent \textbf{limitées et transitoires} chez le sujet sain.
\end{experience}

\begin{popfigure}
\begin{tikzpicture}
\begin{axis}[
width=\linewidth, height=6.5cm,
xlabel={Temps (h)}, ylabel={Glycémie (g/L)},
xmin=0, xmax=3.2, ymin=0.6, ymax=1.6,
xtick={0,0.5,1,1.5,2,2.5,3},
ytick={0.7,0.8,0.9,1.0,1.1,1.2,1.3,1.4,1.5},
grid=both, grid style={popline!40},
legend style={at={(0.97,0.97)}, anchor=north east, font=\small}
]
\addplot[popcurve, smooth, thick] coordinates {
(0,0.95) (0.5,1.25) (1,1.40) (1.5,1.20) (2,1.05) (2.5,0.95) (3,0.92)
};
\addlegendentry{Sujet sain}
\addplot[poporange, dashed, thick, domain=0:3.2] {1.0};
\addlegendentry{Valeur de référence : 1 g/L}
\addplot[popmuted, dotted, domain=0:3.2] {1.1};
\addplot[popmuted, dotted, domain=0:3.2] {0.7};
\end{axis}
\end{tikzpicture}
\legende{Évolution de la glycémie d'un sujet sain avant et après un repas riche en glucides (repas pris à t = 0 h). La glycémie dépasse transitoirement \SI{1,1}{g/L} (pic à \SI{1,4}{g/L} vers 1 h) puis revient vers \SI{1}{g/L} en moins de 3 h. Les pointillés délimitent la fourchette normale à jeun (0,7 à \SI{1,1}{g/L}).}
\end{popfigure}

\begin{exemplebox}[Lecture chiffrée de la courbe]
Chez le sujet sain étudié : la valeur à jeun est de \SI{0,95}{g/L} ; le pic post-prandial atteint environ \SI{1,4}{g/L} une heure après le repas ; le retour à une valeur inférieure à \SI{1,1}{g/L} s'effectue en moins de 2 h (à 2 h, glycémie = \SI{1,05}{g/L}). On conclut que ce sujet présente une \textbf{régulation glycémique efficace}.
\end{exemplebox}

\begin{methode}[Interpréter une courbe de glycémie]
Face à une courbe de glycémie au Baccalauréat, procède toujours dans le même ordre :
\begin{enumerate}
\item \textbf{Repérer la valeur à jeun} (ordonnée à l'origine, avant le repas) et la comparer à la norme (0,7 à \SI{1,1}{g/L}).
\item \textbf{Repérer le pic} : sa valeur et le temps auquel il survient.
\item \textbf{Mesurer le temps de retour à la normale} (retour sous \SI{1,1}{g/L}).
\item \textbf{Comparer à la courbe témoin} (sujet sain) et conclure : régulation normale, hyperglycémie durable (diabète) ou hypoglycémie.
\end{enumerate}
\end{methode}

\courssub{Les facteurs de variation de la glycémie}

L'expérience précédente montre que l'alimentation fait varier la glycémie. D'autres observations cliniques et expérimentales permettent de dresser la liste complète des facteurs.

\textbf{Facteurs d'augmentation de la glycémie (hyperglycémiantes).}
\begin{itemize}
\item \textbf{L'alimentation} : la digestion des glucides (amidon du fufu, sucres des fruits) libère du glucose dans l'intestin grêle, qui passe dans le sang par absorption intestinale. C'est le facteur majeur : la glycémie s'élève après chaque repas glucidique.
\item \textbf{Le stress et les émotions} : un choc émotionnel, la peur ou un examen provoquent la libération d'hormones (adrénaline) qui mobilisent les réserves de glucose du foie ; la glycémie s'élève transitoirement.
\end{itemize}

\textbf{Facteurs de diminution de la glycémie (hypoglycémiantes).}
\begin{itemize}
\item \textbf{Le jeûne} : en l'absence d'apport alimentaire, la glycémie tend à baisser ; c'est la situation de la vendeuse de Mokolo (\SI{0,45}{g/L} après une nuit et une matinée sans manger).
\item \textbf{L'exercice musculaire prolongé} : les muscles en activité consomment massivement le glucose sanguin pour produire de l'ATP ; un match de football ou un travail champêtre prolongé fait chuter la glycémie si aucun apport ne compense.
\end{itemize}

\begin{aretenir}[Le constat fondamental : la glycémie est un paramètre régulé]
Malgré les variations imposées par l'alimentation, le stress, le jeûne ou l'exercice, la glycémie d'un sujet sain \textbf{revient toujours vers une valeur moyenne constante} (environ \SI{1}{g/L}). On dit que la glycémie est un \textbf{paramètre régulé du milieu intérieur} : il existe dans l'organisme des mécanismes de \cle{régulation} capables de corriger les écarts. Ces mécanismes (rôles du pancréas et du foie, hormones insuline et glucagon) seront étudiés dans la section suivante.
\end{aretenir}

\courssub{Exploitation comparative : l'hyperglycémie provoquée par voie orale (HGPO)}

Pour tester l'efficacité de la régulation glycémique, les médecins réalisent une \cle{HGPO} (hyperglycémie provoquée par voie orale) : le sujet à jeun ingère \textbf{75 g de glucose} dissous dans de l'eau, et on mesure sa glycémie toutes les 30 min pendant 2 à 3 h. Comparons un sujet sain et un sujet diabétique.

\begin{center}
\begin{tabularx}{\linewidth}{|l|*{7}{>{\centering\arraybackslash}X|}}
\hline
\rowcolor{popblueL} Temps (h) & 0 & 0,5 & 1 & 1,5 & 2 & 2,5 & 3 \\
\hline
Sujet sain (g/L) & 0,90 & 1,30 & 1,40 & 1,15 & 0,95 & 0,90 & 0,90 \\
\hline
Sujet diabétique (g/L) & 1,60 & 2,20 & 2,60 & 2,50 & 2,30 & 2,10 & 1,95 \\
\hline
\end{tabularx}
\end{center}

\begin{popfigure}
\begin{tikzpicture}
\begin{axis}[
width=\linewidth, height=7cm,
xlabel={Temps (h)}, ylabel={Glycémie (g/L)},
xmin=0, xmax=3.2, ymin=0.6, ymax=2.9,
xtick={0,0.5,1,1.5,2,2.5,3},
ytick={0.7,1.0,1.4,1.8,2.2,2.6},
grid=both, grid style={popline!40},
legend style={at={(0.97,0.97)}, anchor=north east, font=\small}
]
\addplot[popcurve, smooth, thick] coordinates {
(0,0.90) (0.5,1.30) (1,1.40) (1.5,1.15) (2,0.95) (2.5,0.90) (3,0.90)
};
\addlegendentry{Sujet sain}
\addplot[color=poporange, very thick, smooth] coordinates {
(0,1.60) (0.5,2.20) (1,2.60) (1.5,2.50) (2,2.30) (2.5,2.10) (3,1.95)
};
\addlegendentry{Sujet diabétique}
\addplot[popmuted, dashed, domain=0:3.2] {1.1};
\addlegendentry{Limite supérieure à jeun : 1,1 g/L}
\addplot[popmuted, dotted, domain=0:3.2] {1.8};
\addlegendentry{Seuil rénal : 1,8 g/L}
\end{axis}
\end{tikzpicture}
\legende{Courbes d'hyperglycémie provoquée par voie orale (ingestion de 75 g de glucose à t = 0). Chez le sujet sain, le pic (\SI{1,4}{g/L}) est modéré et le retour à la normale s'effectue en moins de 2 h. Chez le sujet diabétique, la glycémie à jeun est déjà élevée (\SI{1,6}{g/L}), le pic dépasse \SI{2,6}{g/L} (au-delà du seuil rénal, d'où une glycosurie) et la glycémie ne revient pas à la normale au bout de 3 h : la régulation est défaillante.}
\end{popfigure}

\textbf{Interprétation.} Trois différences sautent aux yeux :
\begin{enumerate}
\item la glycémie \textbf{à jeun} du diabétique est déjà anormalement élevée (\SI{1,6}{g/L} $>$ \SI{1,1}{g/L}) ;
\item le \textbf{pic} est beaucoup plus haut (\SI{2,6}{g/L}) et dépasse le seuil rénal : du glucose apparaît alors dans les urines (glycosurie) ;
\item le \textbf{retour à la normale n'a pas lieu} en 3 h : la régulation glycémique est défaillante.
\end{enumerate}
On retiendra les critères diagnostiques usuels : une glycémie à jeun supérieure ou égale à \SI{1,26}{g/L} à deux reprises, ou supérieure à \SI{2}{g/L} deux heures après l'ingestion de 75 g de glucose, signe un diabète.

\begin{exoresolu}[Diagnostic d'une HGPO]
\textbf{Énoncé.} Lors d'une campagne de dépistage à l'hôpital de district de Bafoussam, un patient de 48 ans présente les résultats suivants après ingestion de 75 g de glucose : glycémie à jeun = \SI{1,35}{g/L} ; glycémie à 1 h = \SI{2,4}{g/L} ; glycémie à 2 h = \SI{2,2}{g/L}. Ce patient est-il diabétique ? Justifier.
\tcblower
\textbf{Solution détaillée.} Appliquons la méthode d'interprétation.
\begin{enumerate}
\item \emph{Valeur à jeun} : \SI{1,35}{g/L}, supérieure à la limite normale de \SI{1,1}{g/L} et au seuil diagnostique de \SI{1,26}{g/L}.
\item \emph{Pic} : \SI{2,4}{g/L} à 1 h, très supérieur au pic du sujet sain (\SI{1,4}{g/L}) et supérieur au seuil rénal (\SI{1,8}{g/L}) : le patient présente donc aussi une glycosurie.
\item \emph{Retour à la normale} : à 2 h, la glycémie vaut encore \SI{2,2}{g/L}, largement au-dessus de \SI{1,1}{g/L} ; le critère des 2 h (\SI{2}{g/L}) est dépassé.
\end{enumerate}
\textbf{Conclusion} : ce patient est diabétique, car sa glycémie à jeun dépasse \SI{1,26}{g/L} et sa glycémie à 2 h de l'HGPO dépasse \SI{2}{g/L}. Sa régulation glycémique est défaillante ; un suivi médical (régime alimentaire, activité physique, traitement éventuel) s'impose.
\end{exoresolu}

\begin{exercice}[À toi de jouer]
Un élève de Terminale C mesure sa glycémie à jeun : \SI{0,85}{g/L}. Après un petit-déjeuner composé de beignets et de bouillie sucrée, elle passe à \SI{1,35}{g/L} à 1 h, puis redescend à \SI{1,05}{g/L} à 2 h. (1) Sa glycémie à jeun est-elle normale ? (2) Sa régulation est-elle efficace ? Justifier chaque réponse par des valeurs chiffrées.
\end{exercice}

\begin{corrige}[Exercice : glycémie de l'élève]
(1) Oui : \SI{0,85}{g/L} appartient à la fourchette normale à jeun (0,7 à \SI{1,1}{g/L}). (2) Oui : le pic post-prandial (\SI{1,35}{g/L}) reste modéré et proche de celui du sujet sain (\SI{1,4}{g/L}), et la glycémie revient sous \SI{1,1}{g/L} en 2 h. La régulation glycémique de cet élève est donc efficace.
\end{corrige}

\textbf{Bilan de la section.} La glycémie, taux de glucose sanguin, vaut environ \SI{1}{g/L} à jeun. Elle varie sous l'effet de l'alimentation, du stress (augmentation), du jeûne et de l'exercice musculaire (diminution), mais elle revient toujours vers sa valeur de référence chez le sujet sain : c'est un paramètre \textbf{régulé} du milieu intérieur. La comparaison des courbes d'HGPO révèle que cette régulation est défaillante chez le diabétique. Reste à découvrir \emph{comment} l'organisme réalise cet exploit : c'est l'objet de la section suivante, consacrée aux rôles du pancréas et du foie.

\coursec{La régulation hormonale de la glycémie : rôles du pancréas et du foie}

Dans la section précédente, nous avons vu que la glycémie d'un individu sain reste remarquablement stable autour de \SI{1}{g/L} (environ \SI{5,5}{mmol/L}), malgré les apports discontinus de glucose liés aux repas et les consommations variables liées à l'effort. Après un repas riche en glucides, la glycémie monte transitoirement vers \SI{1,4}{g/L} puis revient à sa valeur de consigne en moins de deux heures ; à jeun prolongé, elle ne s'effondre pas. Une telle stabilité ne peut pas être un hasard : elle implique l'existence d'un \cle{système de régulation}. Quels organes en sont responsables ? Par quels messages agissent-ils ? C'est à ces questions que répondent des expériences historiques devenues classiques, que tu dois savoir analyser pour le Baccalauréat.

\courssub{Mise en évidence expérimentale du rôle du pancréas}

\textbf{Observation de départ.} Les médecins savent depuis longtemps que certains patients présentent une glycémie élevée en permanence et éliminent du glucose dans leurs urines (glycosurie) : c'est le diabète sucré. Mais quel organe est en cause ? La démarche expérimentale consiste à retirer chirurgicalement un organe suspect chez un animal et à observer les conséquences : c'est la méthode d'\cle{ablation} (ou résection), complétée ensuite par la méthode de \cle{greffe} ou d'\cle{injection d'extrait}.

\begin{experience}[La pancréatectomie : Mering et Minkowski (1889)]
\textbf{Protocole.} Mering et Minkowski réalisent l'ablation totale du pancréas (\cle{pancréatectomie}) chez un chien, puis suivent l'évolution de la glycémie et de la composition des urines de l'animal, comparées à celles d'un chien témoin non opéré.

\textbf{Résultats observés.} Le chien pancréatectomisé présente :
\begin{itemize}
\item une \cle{hyperglycémie} permanente : la glycémie dépasse durablement \SI{1,8}{g/L} ;
\item une \cle{glycosurie} : du glucose apparaît dans les urines (normalement absent) ;
\item une soif intense (polydipsie), des urines abondantes (polyurie), un amaigrissement rapide ;
\item la mort de l'animal en quelques semaines, dans un état proche du coma diabétique humain.
\end{itemize}

\textbf{Interprétation.} L'ensemble de ces symptômes reproduit exactement le diabète sucré humain. Le pancréas est donc nécessaire au maintien d'une glycémie normale : il exerce une action \cle{hypoglycémiante}.

\textbf{Conclusion.} Le pancréas intervient dans la régulation de la glycémie. Reste à savoir \emph{comment} : par son suc digestif (fonction exocrine) ou par une sécrétion interne (fonction endocrine) ?
\end{experience}

\begin{experience}[Ligature du canal pancréatique puis injection d'extrait : Banting et Best (1921)]
\textbf{Premier temps : la ligature du canal pancréatique.} On ligature le canal excréteur du pancréas chez un chien : le suc pancréatique ne peut plus être déversé dans le duodénum. La partie \cle{exocrine} de la glande (acini sécrétant les enzymes digestives) dégénère progressivement, mais les \cle{îlots de Langerhans}, amas de cellules disséminés dans la glande, restent intacts.

\textbf{Résultat.} Le chien ne présente \emph{aucune} hyperglycémie : sa glycémie reste normale, bien que sa digestion pancréatique soit abolie.

\textbf{Interprétation.} Le suc pancréatique (sécrétion exocrine) n'est pas en cause dans la régulation de la glycémie. L'action hypoglycémiante provient des îlots de Langerhans, qui déversent leur produit directement dans le sang : ce sont des structures \cle{endocrines}.

\textbf{Deuxième temps : l'injection d'extrait pancréatique.} Banting et Best préparent un extrait de pancréas (riche en sécrétion des îlots) et l'injectent à un chien préalablement pancréatectomisé, donc hyperglycémique.

\textbf{Résultat.} La glycémie de l'animal diminue rapidement et revient vers la normale ; le glucose disparaît des urines ; l'animal survit tant que les injections sont poursuivies.

\textbf{Conclusion.} Les îlots de Langerhans sécrètent dans le sang une substance hypoglycémiante : c'est l'\cle{insuline}, la première hormone hypoglycémiante identifiée.
\end{experience}

\begin{propriete}[Rôles du pancréas endocrine et du foie dans la régulation de la glycémie]
Le \cle{pancréas endocrine} (îlots de Langerhans) et le \cle{foie} sont les organes essentiels de la régulation de la glycémie. Cette propriété est établie expérimentalement par :
\begin{itemize}
\item la pancréatectomie (Mering et Minkowski, 1889) qui provoque une hyperglycémie mortelle ;
\item la ligature du canal pancréatique qui montre que seuls les îlots de Langerhans (tissu endocrine) sont impliqués ;
\item l'injection d'extrait pancréatique (Banting et Best, 1921) qui fait baisser la glycémie, prouvant l'existence d'une hormone hypoglycémiante, l'insuline.
\end{itemize}
L'analyse de ces expériences (protocole, résultats, conclusion) est exigible au Baccalauréat.
\end{propriete}

\begin{popfigure}
\begin{center}
\renewcommand{\arraystretch}{1,3}
\begin{tabularx}{\linewidth}{|l|X|X|X|}
\hline
\rowcolor{popblueL} \textbf{Expérience} & \textbf{Protocole} & \textbf{Résultat} & \textbf{Conclusion} \\
\hline
Pancréatectomie (Mering et Minkowski, 1889) & Ablation totale du pancréas chez le chien & Hyperglycémie permanente ($>$ \SI{1,8}{g/L}), glycosurie, mort de l'animal & Le pancréas est nécessaire au maintien d'une glycémie normale \\
\hline
Ligature du canal pancréatique & Le suc pancréatique n'est plus excrété ; les acini exocrines dégénèrent, les îlots de Langerhans restent intacts & Glycémie normale & La régulation ne dépend pas du suc exocrine mais des îlots de Langerhans (endocrine) \\
\hline
Injection d'extrait pancréatique (Banting et Best, 1921) & Extrait d'îlots injecté à un chien pancréatectomisé hyperglycémique & Chute rapide de la glycémie, disparition de la glycosurie, survie de l'animal & Les îlots sécrètent une hormone hypoglycémiante : l'insuline \\
\hline
\end{tabularx}
\end{center}
\legende{Tableau-synthèse des expériences historiques ayant établi le rôle endocrine du pancréas dans la régulation de la glycémie. Ce type de tableau est un excellent modèle de réponse à une question d'analyse de documents au Baccalauréat.}
\end{popfigure}

\courssub{Deux hormones pancréatiques antagonistes : insuline et glucagon}

L'étude histologique et biochimique des îlots de Langerhans révèle qu'ils contiennent deux types principaux de cellules endocrines, sécrétant deux hormones aux effets opposés : on dit qu'elles sont \cle{antagonistes}.

\begin{definition}[Insuline et glucagon]
L'\cle{insuline} est l'hormone \cle{hypoglycémiante} (elle fait baisser la glycémie) sécrétée par les \cle{cellules $\beta$} des îlots de Langerhans, en réponse à une élévation de la glycémie.

Le \cle{glucagon} est l'hormone \cle{hyperglycémiante} (elle fait monter la glycémie) sécrétée par les \cle{cellules $\alpha$} des îlots de Langerhans, en réponse à une baisse de la glycémie.
\end{definition}

L'insuline agit de trois façons complémentaires :
\begin{itemize}
\item elle favorise l'\cle{entrée du glucose dans les cellules} (notamment les cellules musculaires et adipeuses) en facilitant son passage à travers la membrane plasmique via les transporteurs GLUT4 ;
\item elle stimule l'\cle{utilisation du glucose} par les cellules (oxydation lors de la respiration cellulaire, source d'énergie) ;
\item elle active, dans le foie, le stockage du glucose sous forme de glycogène (voir plus loin).
\end{itemize}
Le glucagon, lui, agit essentiellement sur le foie en provoquant la libération de glucose dans le sang.

\begin{attention}[Ne pas confondre les deux hormones]
Piège classique du Baccalauréat : inverser les effets de l'insuline et du glucagon, ou attribuer leur sécrétion aux mauvaises cellules. Retiens le moyen mnémotechnique : \textbf{I}nsuline = \textbf{I}nférieure (elle fait descendre la glycémie) ; \textbf{G}lucagon = \textbf{G}lucose en haut (il fait monter la glycémie). Et : insuline = cellules $\beta$, glucagon = cellules $\alpha$. Autre erreur fréquente : écrire que l'insuline « détruit » le glucose. Elle ne détruit rien : elle favorise son stockage, son entrée dans les cellules et son utilisation.
\end{attention}

\courssub{Le rôle du foie : organe de stockage et de libération du glucose}

\textbf{Problème.} Le pancréas sécrète des hormones, mais sur quel organe agissent-elles principalement ? L'expérience consiste à comparer la glycémie du sang qui \emph{entre} dans le foie (par la veine porte hépatique, venant de l'intestin) et celle du sang qui en \emph{sort} (par les veines sus-hépatiques), dans deux situations : après un repas et à jeun.

\begin{experience}[Mesures de glycémie en amont et en aval du foie]
\textbf{Protocole.} Chez un animal (ou lors de prélèvements chirurgicaux), on mesure la glycémie dans la veine porte (sang entrant dans le foie) et dans la veine sus-hépatique (sang sortant du foie).

\textbf{Résultats.}

\begin{center}
\renewcommand{\arraystretch}{1,3}
\begin{tabularx}{\linewidth}{|l|X|X|X|}
\hline
\rowcolor{popblueL} \textbf{Situation} & \textbf{Veine porte (entrée)} & \textbf{Veine sus-hépatique (sortie)} & \textbf{Bilan} \\
\hline
Après un repas riche en glucides & \SI{1,8}{g/L} & \SI{1,1}{g/L} & Le foie \emph{retient} du glucose \\
\hline
À jeun (quelques heures après le repas) & \SI{0,9}{g/L} & \SI{1,0}{g/L} & Le foie \emph{libère} du glucose \\
\hline
\end{tabularx}
\end{center}

\textbf{Interprétation.} Après un repas, le sang arrivant au foie est très riche en glucose (absorption intestinale) ; le sang qui en sort est moins chargé : le foie a \cle{stocké} une partie du glucose. À jeun, le sang sortant du foie est plus riche en glucose que le sang entrant : le foie \cle{libère} du glucose dans le sang.

\textbf{Conclusion.} Le foie est un organe \cle{tampon} de la glycémie : il stocke le glucose en période d'abondance et le restitue en période de jeûne.
\end{experience}

\textbf{Les mécanismes hépatiques (niveau cellulaire, savoir admis).} Dans les hépatocytes (cellules du foie), le glucose est stocké sous la forme d'un polymère de réserve, le \cle{glycogène} :
\begin{itemize}
\item la \cle{glycogenèse} (ou glycogénogenèse) est la synthèse du glycogène à partir du glucose ; elle est stimulée par l'\cle{insuline} après un repas ;
\item la \cle{glycogénolyse} est la dégradation du glycogène qui libère du glucose dans le sang ; elle est stimulée par le \cle{glucagon} (et par l'adrénaline) à jeun ou lors d'un besoin urgent.
\end{itemize}

\begin{attention}[Le foie ne « produit » pas le glucose dans cette régulation rapide]
Erreur très fréquente dans les copies : écrire que le foie « fabrique » ou « produit » le glucose. Dans le cadre de cette régulation à court terme, le foie \textbf{stocke} le glucose (glycogenèse) et le \textbf{libère} (glycogénolyse) ; il ne le crée pas \emph{de novo}. Le glucose libéré à jeun provient du glycogène accumulé après les repas. Rédige donc : « le foie libère dans le sang le glucose provenant de l'hydrolyse de son glycogène ». La néoglucogenèse (synthèse de glucose à partir de précurseurs non glucidiques) intervient seulement lors de jeûne prolongé.
\end{attention}

\begin{popfigure}
\begin{center}
\begin{tikzpicture}[scale=1, every node/.style={font=\small}]
% Foie
\draw[fill=poporangeL, draw=popdark, thick] (0,0) ellipse (2.6 and 1.5);
\node[font=\bfseries] at (0,1.05) {Foie (hépatocytes)};
% Glycogène
\draw[fill=poppurpleL, draw=poppurple, rounded corners] (-1.1,-0.55) rectangle (1.1,0.15);
\node at (0,-0.2) {Glycogène};
% Veine porte (entrée)
\draw[->, very thick, popblue] (-5.2,0.6) -- (-2.7,0.6) node[midway, above, font=\small]{veine porte};
\node[font=\small, popblue] at (-4.3,0.15) {sang entrant};
% Veine sus-hépatique (sortie)
\draw[->, very thick, popblue] (2.7,0.6) -- (5.2,0.6) node[midway, above, font=\small]{veine sus-hépatique};
\node[font=\small, popblue] at (4.3,0.15) {sang sortant};
% Glycogenèse
\draw[->, very thick, popgreen] (-4.6,-1.6) .. controls (-3.4,-2.3) and (-1.6,-1.9) .. (-0.6,-0.8);
\node[font=\small, popgreen, align=center] at (-3.6,-2.5) {glycogenèse\\ (insuline, après un repas)};
% Glycogénolyse
\draw[->, very thick, poppink] (0.6,-0.8) .. controls (1.6,-1.9) and (3.4,-2.3) .. (4.6,-1.6);
\node[font=\small, poppink, align=center] at (3.6,-2.5) {glycogénolyse\\ (glucagon, à jeun)};
\end{tikzpicture}
\end{center}
\legende{Schéma du rôle tampon du foie. Après un repas, le sang de la veine porte est riche en glucose : sous l'effet de l'insuline, les hépatocytes réalisent la glycogenèse (stockage). À jeun, sous l'effet du glucagon, la glycogénolyse libère du glucose dans le sang sortant par la veine sus-hépatique.}
\end{popfigure}

\courssub{Le schéma fonctionnel de la régulation de la glycémie}

Nous disposons maintenant de tous les éléments : le paramètre réglé (la glycémie), le capteur (les cellules $\alpha$ et $\beta$ des îlots de Langerhans, sensibles au taux de glucose sanguin), les messages (deux hormones véhiculées par le sang : insuline et glucagon) et les effecteurs (le foie, mais aussi les muscles et le tissu adipeux). La régulation fonctionne en \cle{boucle fermée} : la conséquence de l'action (retour de la glycémie à la normale) supprime le stimulus qui l'a déclenchée. C'est un \cle{rétrocontrôle} (régulation par \emph{feed-back} négatif).

\begin{methode}[Construire un schéma fonctionnel de régulation]
Pour réaliser un schéma fonctionnel correct au Baccalauréat, procède en cinq étapes :
\begin{enumerate}
\item \textbf{Identifier le paramètre réglé} : ici, la glycémie (valeur de consigne $\approx$ \SI{1}{g/L}).
\item \textbf{Identifier le capteur} : la structure qui détecte les écarts ; ici, les cellules $\alpha$ et $\beta$ des îlots de Langerhans du pancréas.
\item \textbf{Identifier les messages} : leur nature et leur voie de transport ; ici, deux messages hormonaux (insuline, glucagon) véhiculés par le sang.
\item \textbf{Identifier les effecteurs} : les organes qui corrigent l'écart ; ici, le foie (glycogenèse/glycogénolyse), les muscles et le tissu adipeux (captation du glucose).
\item \textbf{Fermer la boucle} : flécher le retour du paramètre vers le capteur, avec un signe $-$ (rétrocontrôle négatif) : la correction de la glycémie arrête la sécrétion hormonale.
\end{enumerate}
Vérifie enfin que chaque flèche est annotée (nature du message) et que le sens des flèches est cohérent.
\end{methode}

\begin{popfigure}
\begin{center}
\begin{tikzpicture}[scale=1, every node/.style={font=\small}, >=stealth]
% Glycémie
\node[draw=popdark, fill=popcream, rounded corners, align=center, font=\bfseries] (gly) at (0,0) {Glycémie\\ (consigne $\approx$ \SI{1}{g/L})};
% Capteurs
\node[draw=poppurple, fill=poppurpleL, rounded corners, align=center] (beta) at (-5,2.6) {Cellules $\beta$\\ des îlots de Langerhans};
\node[draw=poppurple, fill=poppurpleL, rounded corners, align=center] (alpha) at (-5,-2.6) {Cellules $\alpha$\\ des îlots de Langerhans};
% Effecteurs
\node[draw=poporange, fill=poporangeL, rounded corners, align=center] (foie1) at (5,2.6) {Foie : glycogenèse\\ Muscles/Adipeux : captation\\ et utilisation du glucose};
\node[draw=poporange, fill=poporangeL, rounded corners, align=center] (foie2) at (5,-2.6) {Foie : glycogénolyse\\ (libération de glucose)};
% Flèches haut : hyperglycémie
\draw[->, very thick, popgreen] (gly.north west) -- (beta.south east) node[midway, above, sloped, font=\small]{hyperglycémie (après un repas)};
\draw[->, very thick, popgreen] (beta.east) -- (foie1.west) node[midway, above, font=\small]{insuline};
\draw[->, very thick, popgreen] (foie1.south west) -- (gly.north east) node[midway, above, sloped, font=\small]{baisse de la glycémie};
% Flèches bas : hypoglycémie
\draw[->, very thick, poppink] (gly.south west) -- (alpha.north east) node[midway, below, sloped, font=\small]{hypoglycémie (jeûne, effort)};
\draw[->, very thick, poppink] (alpha.east) -- (foie2.west) node[midway, below, font=\small]{glucagon};
\draw[->, very thick, poppink] (foie2.north west) -- (gly.south east) node[midway, below, sloped, font=\small]{élévation de la glycémie};
% Rétrocontrôle
\node[align=center, font=\small\itshape, popmuted] at (0,-4.1) {Le retour de la glycémie à la normale inhibe la sécrétion hormonale : rétrocontrôle négatif ($-$)};
\end{tikzpicture}
\end{center}
\legende{Schéma fonctionnel de la régulation de la glycémie. Le pancréas endocrine est à la fois le capteur (cellules $\alpha$ et $\beta$ sensibles au glucose sanguin) et la source des messages hormonaux (glucagon, insuline) ; le foie et les muscles sont les effecteurs. La boucle se referme par rétrocontrôle négatif : la normalisation de la glycémie stoppe la sécrétion.}
\end{popfigure}

\begin{exemplebox}[Deux situations concrètes]
\textbf{Situation 1.} À midi, un élève consomme un repas riche en foutou et en sauce sucrée. Sa glycémie passe de \SI{0,95}{g/L} à \SI{1,35}{g/L}. Les cellules $\beta$ de ses îlots de Langerhans libèrent de l'insuline : le foie stocke le glucose en glycogène (glycogenèse), les muscles le captent et l'oxydent. En moins de deux heures, la glycémie revient vers \SI{1}{g/L}, ce qui interrompt la sécrétion d'insuline.

\textbf{Situation 2.} Le même élève court 45 minutes en cours d'EPS à 16 h, loin du déjeuner. Sa glycémie tend à baisser : les cellules $\alpha$ sécrètent du glucagon, qui déclenche la glycogénolyse hépatique ; le glucose libéré compense la consommation musculaire et la glycémie se maintient.
\end{exemplebox}

\begin{savaistu}[Banting, Best et le prix Nobel de 1923]
Frederick Banting et John Macleod reçoivent le prix Nobel de physiologie ou médecine en 1923, à peine deux ans après la découverte de l'insuline (Banting, indigné que son jeune collaborateur Charles Best soit oublié, partage sa part du prix avec lui). La première insuline administrée à l'Homme, en janvier 1922 à un adolescent de 14 ans mourant de diabète, était extraite de pancréas de \textbf{bœuf} et de \textbf{porc} : l'insuline porcine ne diffère de l'insuline humaine que par un seul acide aminé ! Aujourd'hui, l'insuline est produite par des bactéries génétiquement modifiées (génie génétique), ce qui garantit un approvisionnement sûr, y compris pour les centres de santé camerounais.
\end{savaistu}

\begin{aretenir}[Complément Bac : l'adrénaline, hormone hyperglycémiante de stress]
L'\cle{adrénaline}, sécrétée par les glandes \cle{surrénales} (médullo-surrénales) lors d'un stress, d'une frayeur ou d'un effort violent, est également \cle{hyperglycémiante} : elle stimule, comme le glucagon, la glycogénolyse hépatique. C'est pourquoi une émotion intense (examen, accident) peut provoquer une \cle{hyperglycémie de stress}, transitoire et physiologique, à ne pas confondre avec un diabète. Ce point est un complément exigible : il explique que la glycémie est régulée par un système \emph{plurihormonal}, où l'insuline reste la \emph{seule} hormone hypoglycémiante.
\end{aretenir}

\begin{exoresolu}[Analyse d'expériences sur la régulation de la glycémie]
\textbf{Énoncé.} Chez un chien A, on pratique l'ablation totale du pancréas ; chez un chien B, on ligature uniquement le canal pancréatique ; un chien C sert de témoin. Trois semaines plus tard, on mesure les glycémies à jeun : A : \SI{2,4}{g/L} ; B : \SI{0,98}{g/L} ; C : \SI{1,02}{g/L}. On injecte ensuite au chien A un extrait d'îlots de Langerhans : sa glycémie tombe à \SI{1,1}{g/L} en deux heures.
\begin{enumerate}
\item Interpréter les résultats obtenus chez les chiens A et B.
\item Que prouve l'injection réalisée chez le chien A ?
\item Schématiser la boucle de régulation mise en défaut chez le chien A.
\end{enumerate}
\tcblower
\textbf{Solution détaillée.}
\begin{enumerate}
\item Le chien A, privé de tout son pancréas, présente une hyperglycémie majeure (\SI{2,4}{g/L} contre environ \SI{1}{g/L} chez le témoin C) : le pancréas est donc nécessaire au maintien d'une glycémie normale. Le chien B, dont le canal pancréatique est ligaturé (sécrétion exocrine abolie, îlots de Langerhans intacts), garde une glycémie normale (\SI{0,98}{g/L}) : la régulation glycémique ne dépend pas du suc pancréatique exocrine mais de la partie endocrine de la glande, les îlots de Langerhans.
\item L'injection d'extrait d'îlots au chien A fait chuter sa glycémie de \SI{2,4}{g/L} à \SI{1,1}{g/L} : cela prouve que les îlots de Langerhans sécrètent dans le sang une hormone hypoglycémiante, l'insuline, capable de corriger l'hyperglycémie due à l'absence de pancréas.
\item Chez le chien A, la boucle est rompue au niveau du capteur-message : l'hyperglycémie n'est plus détectée par les cellules $\beta$ (absentes), aucune insuline n'est libérée, le foie ne réalise plus la glycogenèse et les muscles ne captent plus correctement le glucose ; la glycémie reste donc durablement élevée. Le schéma attendu reprend les cinq éléments : glycémie $\uparrow$ $\rightarrow$ cellules $\beta$ $\rightarrow$ insuline $\rightarrow$ foie et muscles $\rightarrow$ glycémie $\downarrow$, avec la flèche de rétrocontrôle négatif fermant la boucle.
\end{enumerate}
\end{exoresolu}

\begin{exercice}[À toi de jouer]
Chez un sujet sain, on prélève du sang dans la veine porte et dans la veine sus-hépatique, une heure après un repas glucidique puis après 12 heures de jeûne.
\begin{enumerate}
\item Prévoir, en le justifiant, le signe de la différence (glycémie porte $-$ glycémie sus-hépatique) dans chaque situation.
\item Nommer l'hormone dominante dans chaque situation et le processus hépatique qu'elle contrôle.
\item Expliquer pourquoi, lors d'une forte émotion, on peut observer une hyperglycémie transitoire chez un sujet non diabétique.
\end{enumerate}
\end{exercice}

\begin{corrige}[Exercice : rôle tampon du foie]
\begin{enumerate}
\item Une heure après le repas : glycémie porte $>$ glycémie sus-hépatique (différence positive), car le foie retient une partie du glucose absorbé par l'intestin pour le stocker. Après 12 h de jeûne : glycémie porte $<$ glycémie sus-hépatique (différence négative), car le foie libère du glucose dans le sang.
\item Après le repas, l'hormone dominante est l'insuline, qui stimule la glycogenèse. Au jeûne, c'est le glucagon, qui stimule la glycogénolyse.
\item L'émotion provoque la sécrétion d'adrénaline par les médullo-surrénales ; l'adrénaline stimule la glycogénolyse hépatique, d'où une libération massive de glucose et une hyperglycémie transitoire, qui disparaît dès que le stress cesse (rétrocontrôle par l'insuline).
\end{enumerate}
\end{corrige}

\coursec{Le diabète et l'hypoglycémie : conséquences des dérèglements de la glycémie}

Dans la section précédente, nous avons vu que la glycémie est maintenue autour de \SI{1}{g/L} grâce à une régulation hormonale fine : le pancréas sécrète l'\cle{insuline} (hormone hypoglycémiante) lorsque la glycémie s'élève, et le \cle{glucagon} (hormone hyperglycémiante) lorsqu'elle s'abaisse, le foie étant l'organe effecteur principal. Mais que se passe-t-il lorsque cette régulation est défaillante ? Lorsqu'un patient se présente au Centre Hospitalier d'Essos à Yaoundé en se plaignant d'uriner sans cesse, d'avoir toujours soif et de maigrir malgré un bon appétit, le médecin prescrit immédiatement un dosage de la glycémie. Cette situation clinique, de plus en plus fréquente dans nos villes camerounaises, illustre les conséquences des dérèglements de la glycémie : le \cle{diabète} lorsque la glycémie reste durablement trop élevée, et l'\cle{hypoglycémie} lorsqu'elle s'effondre dangereusement.

\courssub{Le diabète : une hyperglycémie chronique}

\textbf{Observation et définition.} Chez un sujet sain, même après un repas riche en glucides, la glycémie ne dépasse guère \SI{1,40}{g/L} et revient à sa valeur de base en moins de deux heures. Chez certains patients, en revanche, la glycémie reste élevée en permanence, même à jeun. C'est cette constatation qui fonde la définition médicale du diabète.

\begin{definition}[Le diabète]
Le \cle{diabète} (diabète sucré) est une maladie caractérisée par une \textbf{hyperglycémie chronique}, c'est-à-dire une élévation permanente du taux de glucose dans le sang. On parle de diabète lorsque la \textbf{glycémie à jeun est supérieure ou égale à \SI{1,26}{g/L} (\SI{7}{mmol/L}) à deux reprises}, c'est-à-dire lors de deux dosages effectués à des jours différents. Une glycémie à jeun comprise entre \SI{1,10}{g/L} et \SI{1,25}{g/L} traduit une « intolérance au glucose » (état pré-diabétique).
\end{definition}

\textbf{Les symptômes cardinaux.} Le diabète se manifeste par un ensemble de signes que l'on peut tous relier logiquement à l'hyperglycémie :

\begin{itemize}
\item la \cle{polyurie} : le patient urine abondamment et fréquemment (parfois plus de \SI{3}{L} par jour). Lorsque la glycémie dépasse le \textbf{seuil rénal de réabsorption du glucose} (environ \SI{1,80}{g/L}), les reins ne peuvent plus réabsorber tout le glucose filtré : le glucose passe dans les urines et entraîne de l'eau avec lui (diurèse osmotique) ;
\item la \cle{glycosurie} : présence de glucose dans les urines (normalement nulle), conséquence directe du dépassement du seuil rénal. C'est d'ailleurs cette urine « sucrée » qui a donné son nom à la maladie ;
\item la \cle{polydipsie} : soif intense et permanente, conséquence de la déshydratation provoquée par la polyurie ;
\item la \cle{polyphagie} : faim excessive. Bien que le sang soit riche en glucose, les cellules ne l'utilisent pas correctement (faute d'insuline ou par résistance à son action) : elles « jeûnent au milieu de l'abondance » et envoient des signaux de faim ;
\item l'\cle{amaigrissement} : malgré la polyphagie, le patient maigrit, car l'organisme, incapable d'utiliser le glucose, puise dans ses réserves de lipides et de protéines.
\end{itemize}

\begin{attention}[Ne pas confondre glycémie et glycosurie]
La \textbf{glycémie} est le taux de glucose dans le \emph{sang} ; la \textbf{glycosurie} est le taux de glucose dans les \emph{urines}. Une glycosurie n'apparaît que si la glycémie dépasse le seuil rénal (environ \SI{1,80}{g/L}). Un diabétique peut donc avoir une glycémie de \SI{1,50}{g/L} (anormale !) sans glycosurie. Au Baccalauréat, n'écris jamais que « le diabète se définit par la présence de sucre dans les urines » : c'est la glycémie à jeun qui fait le diagnostic.
\end{attention}

\courssub{Les deux grands types de diabète}

\textbf{Problème scientifique :} tous les diabètes sont-ils dus à la même cause ? L'observation clinique montre que non : certains diabétiques sont des enfants ou adolescents maigres dont la maladie débute brutalement, d'autres sont des adultes en surpoids chez qui la maladie s'installe insidieusement. Les dosages hormonaux ont permis de comprendre cette différence.

\begin{experience}[Mise en évidence expérimentale des deux types de diabète]
\textbf{Protocole :} on dose l'insulinémie (taux d'insuline sanguine) à jeun chez deux groupes de diabétiques : des sujets jeunes (groupe A) et des sujets adultes en surpoids (groupe B).

\textbf{Résultats :}
\begin{center}
\begin{tabularx}{\linewidth}{|l|X|X|}
\hline
\rowcolor{popblueL} \textbf{Paramètre} & \textbf{Groupe A (jeunes)} & \textbf{Groupe B (adultes)} \\
\hline
Insulinémie à jeun & quasi nulle & normale ou élevée \\
\hline
Effet d'une injection d'insuline & glycémie normalisée & effet faible ou nul \\
\hline
Effet d'un régime + exercice & insuffisant seul & glycémie souvent améliorée \\
\hline
\end{tabularx}
\end{center}

\textbf{Interprétation :} chez le groupe A, le pancréas ne produit plus d'insuline ; chez le groupe B, l'insuline est présente mais inefficace : les cellules cibles y sont résistantes.

\textbf{Conclusion :} il existe (au moins) deux mécanismes distincts conduisant à l'hyperglycémie chronique, d'où la classification en diabète de type 1 et diabète de type 2.
\end{experience}

\textbf{Le diabète de type 1 (DID).} Le \cle{diabète de type 1}, ou \textbf{diabète insulinodépendant (DID)}, résulte de la \textbf{destruction des cellules $\beta$ des îlots de Langerhans} du pancréas, le plus souvent par un mécanisme auto-immun (l'organisme fabrique des anticorps contre ses propres cellules $\beta$). Il en résulte une \textbf{absence totale de sécrétion d'insuline}. Il survient classiquement chez le \textbf{sujet jeune} (enfant, adolescent, adulte jeune) et son début est brutal. Le seul traitement possible est l'apport quotidien d'insuline par \textbf{injections sous-cutanées} (l'insuline est une protéine : prise par voie orale, elle serait digérée et détruite), associé à une surveillance régulière de la glycémie et à un régime alimentaire adapté.

\textbf{Le diabète de type 2 (DNID).} Le \cle{diabète de type 2}, ou \textbf{diabète non insulinodépendant (DNID)}, est dû à une \textbf{résistance des cellules cibles (muscles, foie, tissu adipeux) à l'action de l'insuline} : l'insuline est sécrétée, parfois même en excès au début, mais les récepteurs cellulaires y répondent mal, si bien que le glucose pénètre mal dans les cellules et s'accumule dans le sang. Il survient chez le \textbf{sujet adulte} (généralement après 40 ans) et est fortement associé à l'\textbf{obésité}, à la \textbf{sédentarité} et à une alimentation trop riche ; il existe aussi une composante héréditaire. Le traitement repose d'abord sur les \textbf{mesures hygiéno-diététiques} (régime équilibré, réduction des sucres rapides et des graisses, exercice physique régulier), puis sur des \textbf{antidiabétiques oraux} si nécessaire.

\begin{attention}[Erreur fréquente au Baccalauréat]
Ne dis \textbf{jamais} que « le diabète de type 2 est dû à un manque d'insuline ». C'est faux : dans le DNID, \textbf{l'insuline est présente} (parfois en quantité normale ou supérieure à la normale), mais \textbf{les cellules y résistent}. Le défaut se situe au niveau des récepteurs et de la réponse cellulaire, non au niveau de la sécrétion. C'est le diabète de \textbf{type 1} qui est dû à l'absence d'insuline.
\end{attention}

\begin{popfigure}
\begin{center}
\begin{tikzpicture}[font=\small]
% En-têtes
\fill[popblue!85] (0,0) rectangle (4.4,0.8);
\fill[poporange!85] (4.6,0) rectangle (9.4,0.8);
\fill[popgreen!70] (9.6,0) rectangle (14.8,0.8);
\node[white,font=\bfseries] at (2.2,0.4) {Critère};
\node[white,font=\bfseries] at (7.0,0.4) {Diabète de type 1 (DID)};
\node[white,font=\bfseries] at (12.2,0.4) {Diabète de type 2 (DNID)};
% Lignes
\foreach \y/\a/\b/\c in {
  -0.9/{Cause}/{Destruction auto-immune des cellules $\beta$}/{Résistance des cellules à l'insuline},
  -2.0/{Insuline sanguine}/{Absente (sécrétion nulle)}/{Présente (normale ou élevée)},
  -3.1/{Âge d'apparition}/{Sujet jeune (enfant, adolescent)}/{Adulte (souvent $>$ 40 ans)},
  -4.2/{Facteurs associés}/{Prédisposition génétique, auto-immunité}/{Obésité, sédentarité, hérédité},
  -5.3/{Traitement}/{Injections quotidiennes d'insuline}/{Régime, exercice, antidiabétiques oraux}}
{
  \node[anchor=west,font=\bfseries\small,text=popdark] at (0.1,\y) {\a};
  \node[anchor=west,text width=4.5cm,align=flush left] at (4.7,\y) {\b};
  \node[anchor=west,text width=4.9cm,align=flush left] at (9.7,\y) {\c};
  \draw[popline] (0,\y+0.55) -- (14.8,\y+0.55);
}
\draw[popline] (0,-5.85) -- (14.8,-5.85);
\draw[popline] (4.5,0) -- (4.5,-5.85);
\draw[popline] (9.5,0) -- (9.5,-5.85);
\end{tikzpicture}
\end{center}
\legende{Tableau comparatif des deux grands types de diabète sucré. Le DID résulte d'un défaut de \emph{sécrétion} d'insuline, le DNID d'un défaut d'\emph{action} de l'insuline.}
\end{popfigure}

\begin{exemplebox}[Le diabète au Cameroun : un problème de santé publique croissant]
Au Cameroun, la prévalence du diabète en milieu urbain est estimée entre 5 et \SI{7}{\%} de la population adulte, et elle augmente régulièrement avec la sédentarisation des modes de vie : urbanisation, alimentation riche en sucres rapides et en graisses, diminution de l'activité physique. Le diabète de type 2, autrefois rare, devient l'un des motifs de consultation les plus fréquents dans les hôpitaux de Yaoundé et de Douala, et touche des sujets de plus en plus jeunes.
\end{exemplebox}

\courssub{Interprétation des analyses médicales : glycémie à jeun et HGPO}

Le diagnostic du diabète repose sur des examens biologiques précis qu'il faut savoir interpréter :

\begin{itemize}
\item la \textbf{glycémie à jeun} : normale entre \SI{0,70}{g/L} et \SI{1,10}{g/L} ; diabète si $\geq \SI{1,26}{g/L}$ à deux reprises ;
\item l'\cle{HGPO} (hyperglycémie provoquée par le glucose), ou test de tolérance au glucose : le patient à jeun ingère \SI{75}{g} de glucose dissous dans l'eau, puis on dose la glycémie à $T_0$ (à jeun), $T_{60}$ (1 h après) et $T_{120}$ (2 h après). Chez un sujet sain, la glycémie à 2 h est revenue sous \SI{1,40}{g/L} ; \textbf{une glycémie $\geq \SI{2}{g/L}$ à 2 h signe le diabète} ;
\item la \textbf{recherche de glucose et de corps cétoniques dans les urines} : la glycosurie traduit le dépassement du seuil rénal ; la présence de corps cétoniques (cétonurie) signe une dégradation massive des lipides, typique du diabète de type 1 non traité (acidocétose).
\end{itemize}

\begin{methode}[Interpréter une HGPO]
Pour interpréter les résultats d'une HGPO, procède en étapes :
\begin{enumerate}
\item Repère la \textbf{glycémie à jeun} ($T_0$) et compare-la aux seuils : $< \SI{1,10}{g/L}$ normale ; entre \SI{1,10}{g/L} et \SI{1,25}{g/L} intolérance au glucose ; $\geq \SI{1,26}{g/L}$ diabète.
\item Repère la \textbf{glycémie à 2 h} ($T_{120}$) : $< \SI{1,40}{g/L}$ normale ; entre \SI{1,40}{g/L} et \SI{1,99}{g/L} intolérance au glucose ; $\geq \SI{2}{g/L}$ diabète.
\item Observe la \textbf{cinétique} : chez le sujet sain, le pic est modéré ($< \SI{1,80}{g/L}$) et le retour à la normale est rapide ; chez le diabétique, le pic est élevé et la glycémie reste haute.
\item Conclus en citant les valeurs chiffrées du document, jamais de façon vague.
\end{enumerate}
\end{methode}

\begin{popfigure}
\begin{center}
\begin{tikzpicture}
\begin{axis}[
  width=0.85\linewidth, height=6.5cm,
  xlabel={Temps (min)}, ylabel={Glycémie (g/L)},
  xmin=-10, xmax=150, ymin=0.6, ymax=2.6,
  xtick={0,30,60,90,120}, ytick={0.8,1.0,1.26,1.4,1.8,2.0,2.4},
  yticklabels={0,80, 1,00, 1,26, 1,40, 1,80, 2,00, 2,40},
  grid=both, grid style={popline!40},
  legend style={at={(0.98,0.05)},anchor=south east,font=\small,draw=popline}
]
% Seuil diabète à jeun
\addplot[popmuted,dashed,domain=-10:150]{1.26};
% Seuil 2 g/L
\addplot[popmuted,dashed,domain=-10:150]{2.0};
% Sujet sain
\addplot[popgreen,very thick,smooth,mark=*] coordinates {(0,0.90) (30,1.35) (60,1.30) (90,1.05) (120,0.95)};
% Patient diabétique
\addplot[poporange,very thick,smooth,mark=square*] coordinates {(0,1.45) (30,2.20) (60,2.40) (90,2.30) (120,2.15)};
\legend{,,Sujet sain,Patient M.\ N.}
\node[font=\scriptsize,text=popmuted,anchor=west] at (axis cs:128,1.26) {1,26 g/L};
\node[font=\scriptsize,text=popmuted,anchor=west] at (axis cs:128,2.0) {2,00 g/L};
\end{axis}
\end{tikzpicture}
\end{center}
\legende{Résultats d'une HGPO (ingestion de \SI{75}{g} de glucose à $T_0$) chez un sujet sain (courbe verte) et chez le patient M.\ N.\ (courbe orange). Chez le patient : glycémie à jeun de \SI{1,45}{g/L} ($\geq \SI{1,26}{g/L}$) et glycémie de \SI{2,15}{g/L} à 2 h ($\geq \SI{2}{g/L}$) : le diagnostic de diabète est posé.}
\end{popfigure}

\begin{exoresolu}[Interprétation d'un bilan biologique]
\textbf{Énoncé.} M.\ N., 52 ans, commerçant à Douala, en surpoids, consulte pour fatigue et soif permanente. Résultats : glycémie à jeun : \SI{1,45}{g/L} ; HGPO : $T_0 = \SI{1,45}{g/L}$, $T_{60} = \SI{2,40}{g/L}$, $T_{120} = \SI{2,15}{g/L}$ ; urines : glucose présent, corps cétoniques absents ; insulinémie : élevée. 1) Le patient est-il diabétique ? Justifie. 2) De quel type de diabète s'agit-il probablement ? Argumente. 3) Propose une conduite thérapeutique.
\tcblower
\textbf{Solution.}
1) Oui, M.\ N.\ est diabétique : sa glycémie à jeun (\SI{1,45}{g/L}) dépasse le seuil de \SI{1,26}{g/L}, et sa glycémie à 2 h de l'HGPO (\SI{2,15}{g/L}) dépasse le seuil de \SI{2}{g/L}. La glycosurie confirme le dépassement du seuil rénal (\SI{1,80}{g/L}).

2) Il s'agit probablement d'un \textbf{diabète de type 2} : le patient est un adulte de 52 ans, en surpoids ; son insulinémie est \emph{élevée}, ce qui exclut un défaut de sécrétion d'insuline (type 1) et traduit une résistance des cellules à l'insuline ; l'absence de corps cétoniques dans les urines plaide également contre un diabète de type 1 décompensé.

3) Conduite thérapeutique : mesures hygiéno-diététiques d'abord (régime pauvre en sucres rapides et en graisses, perte de poids, exercice physique régulier), puis antidiabétiques oraux si la glycémie reste élevée, avec surveillance régulière de la glycémie et éducation thérapeutique.
\end{exoresolu}

\courssub{Complications du diabète non traité et moyens de lutte}

Un diabète mal contrôlé expose, à long terme, l'excès de glucose qui altère les parois des petits vaisseaux sanguins et les nerfs. Les principales complications sont :

\begin{itemize}
\item les \textbf{atteintes rénales} (néphropathie diabétique) : altération progressive des reins pouvant aller jusqu'à l'insuffisance rénale chronique, nécessitant des dialyses ;
\item les \textbf{atteintes oculaires} (rétinopathie diabétique) : lésions des vaisseaux de la rétine, première cause de cécité acquise chez l'adulte ;
\item les \textbf{atteintes nerveuses} (neuropathie) : perte de sensibilité, fourmillements, douleurs, surtout aux pieds ;
\item les \textbf{atteintes cardiovasculaires} : l'hyperglycémie chronique accélère l'athérosclérose, d'où un risque accru d'infarctus du myocarde et d'accidents vasculaires cérébraux ;
\item le \textbf{pied diabétique} : par perte de sensibilité (neuropathie) et mauvaise circulation, une simple plaie du pied passe inaperçue, s'infecte et peut évoluer vers la gangrène et l'amputation. C'est une complication fréquente et redoutée dans nos hôpitaux.
\end{itemize}

\textbf{Les moyens de lutte} contre le diabète reposent sur :
\begin{enumerate}
\item le \textbf{dépistage} : dosage de la glycémie à jeun chez les personnes à risque (obésité, antécédents familiaux, âge $>$ 40 ans) ; un diabète dépisté tôt se contrôle bien ;
\item l'\textbf{hygiène alimentaire} : limiter les sucres rapides, les boissons sucrées et les graisses ; privilégier les légumes, les céréales complètes et les repas réguliers ;
\item l'\textbf{activité physique régulière} : elle améliore la sensibilité des cellules à l'insuline ;
\item l'\textbf{observance du traitement} : injections d'insuline ou antidiabétiques oraux pris rigoureusement, sans interruption ;
\item l'\textbf{éducation thérapeutique} : apprendre au patient à surveiller sa glycémie, à soigner ses pieds, à reconnaître les signes d'alerte.
\end{enumerate}

\courssub{L'hypoglycémie : conséquences d'une baisse anormale de la glycémie}

\begin{definition}[L'hypoglycémie]
L'\cle{hypoglycémie} est une baisse anormale de la glycémie, définie par une glycémie \textbf{inférieure à \SI{0,60}{g/L}}. Le cerveau, qui utilise presque exclusivement le glucose comme source d'énergie, est le premier organe menacé.
\end{definition}

\textbf{Manifestations.} L'hypoglycémie se traduit par un \textbf{malaise} avec \textbf{sueurs froides}, \textbf{tremblements}, palpitations, faim intense, pâleur, puis, si elle s'aggrave, des \textbf{troubles de la conscience} (confusion, comportement bizarre) pouvant aller jusqu'au \textbf{coma hypoglycémique}, qui engage le pronostic vital.

\textbf{Causes principales :}
\begin{itemize}
\item un \textbf{excès d'insuline} chez un diabétique traité (dose trop forte, repas sauté après l'injection) : c'est la cause la plus fréquente ;
\item un \textbf{jeûne prolongé} ou une alimentation insuffisante ;
\item un \textbf{effort physique intense} non compensé par un apport glucidique.
\end{itemize}

\textbf{Conduite à tenir :} si le sujet est conscient, lui faire absorber rapidement du \textbf{sucre} (morceaux de sucre, boisson sucrée, jus de fruit) ; si le malaise persiste ou si le sujet perd conscience, le mettre en position latérale de sécurité et appeler les secours en urgence.

\begin{savaistu}[Urgence : sucre, jamais d'insuline !]
Le coma hypoglycémique chez un diabétique sous insuline est une \textbf{urgence vitale}. Face à un diabétique inconscient, la conduite à tenir est de lui donner du \textbf{sucre} (par voie orale s'il est conscient, ou glucose par voie veineuse par le personnel soignant) — et \textbf{jamais d'insuline}, qui aggraverait l'hypoglycémie et pourrait être mortelle. En cas de doute entre coma hypoglycémique et hyperglycémique, c'est toujours le sucre qu'il faut donner : il sauve dans l'hypoglycémie et n'aggrave pas durablement une hyperglycémie.
\end{savaistu}

\begin{aretenir}[L'essentiel de la section]
\begin{itemize}
\item Le diabète est une \textbf{hyperglycémie chronique} : glycémie à jeun $\geq \SI{1,26}{g/L}$ à deux reprises ; symptômes : polyurie, polydipsie, polyphagie, amaigrissement, glycosurie.
\item \textbf{Type 1 (DID)} : destruction des cellules $\beta$, insuline absente, sujet jeune, traitement par injections d'insuline. \textbf{Type 2 (DNID)} : résistance des cellules à l'insuline (insuline présente !), sujet adulte souvent obèse et sédentaire, traitement par régime, exercice et antidiabétiques oraux.
\item Diagnostic : glycémie à jeun, HGPO (diabète si $\geq \SI{2}{g/L}$ à 2 h), recherche de glucose et de corps cétoniques urinaires.
\item Complications : reins, yeux, nerfs, c\oe ur et vaisseaux, pied diabétique ; lutte par dépistage, hygiène de vie, observance et éducation thérapeutique.
\item L'\textbf{hypoglycémie} ($< \SI{0,60}{g/L}$) : malaise, sueurs, tremblements, coma ; conduite à tenir : apport rapide de sucre, jamais d'insuline.
\end{itemize}
\end{aretenir}

\coursec{La pression artérielle : définition et mesure}

Après avoir exploré les mécanismes fins de la régulation de la glycémie et les pathologies métaboliques majeures que sont le diabète et l'hypoglycémie, nous abordons maintenant un autre paramètre vital dont la stabilité conditionne la survie de l'organisme : la \cle{pression artérielle}. Tout comme la glycémie, la pression artérielle doit être maintenue dans des limites étroites ; tout écart durable, qu'il soit vers le haut (hypertension) ou vers le bas (hypotension), expose à des complications graves, particulièrement fréquentes dans nos contextes urbains camerounais où le stress, l'alimentation riche en sel et la sédentarité sont en hausse. Commençons par définir ce qu'elle est, comment on la mesure et comment interpréter les valeurs obtenues.

\courssub{Définition et valeurs physiologiques de référence}

\begin{definition}[Pression artérielle (ou tension artérielle)]
La \textbf{pression artérielle} est la pression (force par unité de surface) exercée par le sang sur la paroi des artères. Elle résulte directement de l'activité de pompe du cœur (contraction ventriculaire) et de la résistance vasculaire périphérique. En langage courant, on parle souvent de « tension artérielle » (notée TA).
\end{definition}

Le cœur fonctionne comme une pompe cyclique. À chaque cycle cardiaque, la pression dans l'aorte et les grosses artères oscille entre deux valeurs extrêmes :

\begin{itemize}
    \item La \textbf{pression systolique (ou maxima, notée $P_{sys}$ ou MAX)} : c'est la valeur maximale atteinte lors de la \emph{systole ventriculaire} (contraction des ventricules, éjection du sang). Elle correspond à l'onde de pression qui se propage le long des artères (le pouls).
    \item La \textbf{pression diastolique (ou minima, notée $P_{dias}$ ou MIN)} : c'est la valeur minimale atteinte lors de la \emph{diastole ventriculaire} (relâchement des ventricules, remplissage). Elle ne tombe pas à zéro grâce à l'élasticité des grosses artères (effet \emph{windkessel} ou réservoir élastique) qui se déforment pendant la systole et se restaurent pendant la diastole, maintenant ainsi un flux continu.
\end{itemize}

\begin{exemplebox}[Valeurs de référence au repos chez l'adulte]
Chez un sujet adulte sain, au repos, en position assise ou allongée, les valeurs usuelles sont :
\begin{itemize}
    \item Pression systolique (MAX) : \textbf{12 à 13 cmHg} (soit 120 à 130 mmHg).
    \item Pression diastolique (MIN) : \textbf{7 à 8 cmHg} (soit 70 à 80 mmHg).
\end{itemize}
On note conventionnellement : \textbf{13/8 cmHg} (ou 130/80 mmHg). L'unité officielle au Cameroun dans les services de santé reste souvent le \textbf{cmHg} (centimètre de mercure), bien que le mmHg soit la norme internationale (1 cmHg = 10 mmHg).

\textbf{Pression différentielle (ou pression de pouls)} : $PP = MAX - MIN$. Elle reflète le volume d'éjection systolique et la compliance artérielle. Exemple : $13,0 - 8,0 = 5,0$ cmHg (50 mmHg).
\end{exemplebox}

\begin{attention}[Pièges de notation et d'unités]
\begin{itemize}
    \item \textbf{Ne jamais additionner} les deux valeurs (ex. : 13 + 8 = 21 n'a aucun sens physiologique).
    \item \textbf{Ne pas confondre} avec la glycémie : la tension s'exprime en cmHg (ou mmHg), \emph{jamais} en g/L ou mmol/L.
    \item L'ordre est immuable : \textbf{MAX / MIN} (systolique / diastolique). Écrire « 8/13 » est une erreur grave.
    \item Le terme « tension » est un abus de langage accepté (tension = force par unité de longueur ; pression = force par unité de surface), mais « pression artérielle » est le terme scientifiquement exact.
\end{itemize}
\end{attention}

\courssub{La mesure au sphygmomanomètre : méthode auscultatoire (technique de Korotkoff)}

La méthode de référence, non invasive, utilise un \textbf{sphygmomanomètre à mercure} (ou anéroïde) et un \textbf{stéthoscope}. C'est une manipulation que vous devez savoir décrire, justifier et interpréter (savoir-faire exigible au Baccalauréat).

\begin{experience}[Mesure de la pression artérielle par la méthode auscultatoire]
\textbf{Matériel :} Sphygmomanomètre (brassard gonflable, poire de gonflage, manomètre à colonne de mercure ou cadran), stéthoscope.
\textbf{Sujet :} Au repos depuis au moins 5 à 10 min, assis, bras nu soutenu à la hauteur du cœur (niveau du 4\textsuperscript{e} espace intercostal), paume vers le haut.

\textbf{Protocole :}
\begin{enumerate}
    \item \textbf{Mise en place :} Enrouler le brassard autour du bras, son bord inférieur à 2--3 cm au-dessus du pli du coude. Placer le pavillon du stéthoscope sur l'\textbf{artère humérale} (dans le pli du coude, face interne du bras), sans appuyer excessivement. Fermer la vis de la poire.
    \item \textbf{Gonflage :} Gonfler rapidement le brassard en pressant la poire jusqu'à ce que la pression indiquée par le manomètre dépasse d'environ \textbf{20 à 30 mmHg (soit 2 à 3 cmHg)} la pression systolique estimée (généralement jusqu'à \textbf{18--20 cmHg} au total). \emph{Observation :} Les bruits cardiaques disparaissent complètement ; l'artère est totalement écrasée, le flux sanguin est nul (stase).
    \item \textbf{Dégonflage progressif :} Ouvrir légèrement la vis de la poire pour laisser l'air s'échapper \emph{lentement} (vitesse idéale : 2 à 3 mmHg/s, soit 0,2 à 0,3 cmHg/s).
    \item \textbf{Auscultation :} Écouter attentivement les bruits au niveau de l'artère humérale pendant la descente de la pression.
    \item \textbf{Lecture des valeurs :}
    \begin{itemize}
        \item \textbf{Apparition des premiers bruits (bruits de Korotkoff, phase I) :} La pression lue sur le manomètre à cet instant est la \textbf{pression systolique (MAX)}.
        \item \textbf{Disparition définitive des bruits (phase V) :} La pression lue à cet instant est la \textbf{pression diastolique (MIN)}.
    \end{itemize}
    \item \textbf{Dégonflage complet :} Ouvrir totalement la vis pour dégonfler le brassard rapidement et le retirer.
\end{enumerate}

\textbf{Interprétation physique des bruits de Korotkoff :}
\begin{itemize}
    \item \textbf{Pression du brassard $>$ MAX :} Artère totalement occlusée $\rightarrow$ flux nul $\rightarrow$ silence.
    \item \textbf{MAX $>$ Pression du brassard $>$ MIN :} Artère partiellement comprimée $\rightarrow$ flux turbulent, intermittent (systole seulement) $\rightarrow$ \textbf{bruits de Korotkoff} (claquement, souffle rythmique).
    \item \textbf{Pression du brassard $<$ MIN :} Artère perméable en permanence $\rightarrow$ flux laminaire, continu $\rightarrow$ silence (ou bruits très faibles non pris en compte pour la MIN en pratique courante).
\end{itemize}
\end{experience}

\begin{popfigure}
\begin{tikzpicture}[scale=0.9, every node/.style={font=\small}]
    % Couleurs autorisées : poppink pour le sang artériel (remplace popred)
    \colorlet{arteryred}{poppink!80!popdark}
    \colorlet{cuffblue}{popblue!70!popdark}
    \colorlet{skin}{poporange!20!popcream}
    \colorlet{bone}{popcream!80!popmuted}
    
    % Bras (coupe transversale simplifiée)
    \draw[fill=skin, draw=popmuted, thick] (-3,-2) rectangle (3,2);
    % Os (humérus)
    \draw[fill=bone, draw=popmuted] (0,0) ellipse (0.8 and 0.5);
    % Artère humérale
    \draw[fill=arteryred, draw=arteryred!50!popdark, thick] (0,0.8) circle (0.35);
    \node[arteryred, font=\bfseries] at (0,0.8) {Art. humérale};
    
    % Brassard (autour du bras)
    \draw[fill=cuffblue!30, draw=cuffblue, very thick, rounded corners=3pt] (-3.5,-2.5) rectangle (3.5,2.5);
    \node[cuffblue, font=\bfseries, rotate=90] at (-3.8,0) {Brassard};
    
    % Tube vers manomètre
    \draw[cuffblue, very thick] (3.5,0) -- (5,0) -- (5,2.5);
    % Manomètre (schématique)
    \draw[fill=popcream, draw=popdark, thick] (4.5,2.5) rectangle (7,4.5);
    \draw[popdark, thick] (5.75,3.5) -- (5.75,4.2); % Aiguille
    \node[popdark, font=\bfseries] at (5.75,3.0) {Manomètre};
    \node[popdark, font=\tiny] at (5.75,4.3) {cmHg};
    
    % Stéthoscope
    \draw[poppurple, very thick] (-1.5,-2) -- (-1.5,-3.5) -- (-3.5,-3.5) -- (-3.5,-4.5);
    \draw[fill=poppurple!30, draw=poppurple, thick] (-3.5,-4.5) circle (0.4);
    \node[poppurple, font=\bfseries] at (-1.5,-2.3) {Pavillon};
    \node[poppurple, font=\bfseries] at (-3.5,-5.2) {Pavillon};
    
    % Flèches indicatives
    \draw[->, popdark, thick] (0,2.5) -- (0,3.5) node[above, font=\bfseries] {Pression brassard};
    \draw[->, arteryred, thick] (1.5,1.5) -- (2.5,2.5) node[right, font=\bfseries] {Flux sanguin};
    
    % Légendes positions
    \node[align=center, font=\tiny, popmuted] at (0,-3) {Coupe transversale du bras\\ (niveau pli du coude)};
\end{tikzpicture}
\legende{Schéma du dispositif de mesure de la pression artérielle : brassard positionné sur le bras (coupe transversale), stéthoscope posé sur l'artère humérale au pli du coude, raccordé au manomètre.}
\end{popfigure}

\begin{popfigure}
\begin{tikzpicture}[scale=1.0, every node/.style={font=\small}]
    % Trois phases : A (écrasée), B (intermittente), C (libre)
    \def\artx{0}
    \def\arty{0}
    \def\artr{0.6}
    \def\cuffr{1.2}
    \colorlet{arteryred}{poppink!80!popdark} % Couleur sang
    
    % --- Phase A : Brassard > MAX ---
    \begin{scope}[xshift=-7cm]
        \node[font=\bfseries, popdark] at (0,2.2) {Phase A : $P_{brassard} > MAX$};
        % Brassard
        \draw[fill=popblue!20, draw=popblue, very thick] (-\cuffr,-\cuffr) rectangle (\cuffr,\cuffr);
        % Artère écrasée (ligne)
        \draw[arteryred, ultra thick] (-0.8,0) -- (0.8,0);
        \node[arteryred, font=\bfseries\small] at (0,-1.1) {Artère occlusée};
        % Flux
        \draw[arteryred, ->, thick, dashed] (-1.5,0) -- (-0.9,0);
        \draw[arteryred, ->, thick, dashed] (0.9,0) -- (1.5,0);
        \node[popmuted, font=\tiny] at (0,1.5) {Flux = 0};
        % Manomètre
        \node[align=center, font=\tiny, popblue] at (0,-1.7) {Manomètre : ex. 18 cmHg};
        \node[align=center, font=\tiny, popdark] at (0,-2.1) {Silence au stéthoscope};
    \end{scope}
    
    % --- Phase B : MAX > Brassard > MIN ---
    \begin{scope}[xshift=0cm]
        \node[font=\bfseries, popdark] at (0,2.2) {Phase B : $MAX > P_{brassard} > MIN$};
        % Brassard
        \draw[fill=poporange!20, draw=poporange, very thick] (-\cuffr,-\cuffr) rectangle (\cuffr,\cuffr);
        % Artère partiellement ouverte (forme ovoïde irrégulière)
        \draw[fill=arteryred!30, draw=arteryred, thick] (0,0) ellipse (0.5 and 0.25);
        \node[arteryred, font=\bfseries\small] at (0,0) {Turbulence};
        % Flux intermittent (flèches saccadées)
        \foreach \y in {-0.15, 0.15} {
            \draw[arteryred, ->, thick] (-1.5,\y) -- (-0.7,\y);
            \draw[arteryred, ->, thick] (0.7,\y) -- (1.5,\y);
        }
        \node[popmuted, font=\tiny] at (0,1.5) {Flux intermittent (systole only)};
        % Manomètre
        \node[align=center, font=\tiny, poporange] at (0,-1.7) {Manomètre : descend 13 $\to$ 8};
        \node[align=center, font=\tiny, popdark] at (0,-2.1) {Bruits de Korotkoff entendus};
    \end{scope}
    
    % --- Phase C : Brassard < MIN ---
    \begin{scope}[xshift=7cm]
        \node[font=\bfseries, popdark] at (0,2.2) {Phase C : $P_{brassard} < MIN$};
        % Brassard
        \draw[fill=popgreen!20, draw=popgreen, very thick] (-\cuffr,-\cuffr) rectangle (\cuffr,\cuffr);
        % Artère ouverte (cercle)
        \draw[fill=arteryred!30, draw=arteryred, thick] (0,0) circle (0.5);
        \node[arteryred, font=\bfseries\small] at (0,0) {Libre};
        % Flux continu (flèches parallèles)
        \foreach \y in {-0.35, -0.1, 0.15, 0.4} {
            \draw[arteryred, ->, thick] (-1.5,\y) -- (1.5,\y);
        }
        \node[popmuted, font=\tiny] at (0,1.5) {Flux continu, laminaire};
        % Manomètre
        \node[align=center, font=\tiny, popgreen] at (0,-1.7) {Manomètre : < 8 cmHg};
        \node[align=center, font=\tiny, popdark] at (0,-2.1) {Silence (fin des bruits)};
    \end{scope}
    
    % Flèches de progression
    \draw[popdark, ->, very thick] (-3.5,-2.8) -- (-0.5,-2.8) node[midway, below, font=\small] {Dégonflage lent};
    \draw[popdark, ->, very thick] (3.5,-2.8) -- (6.5,-2.8) node[midway, below, font=\small] {Dégonflage lent};
\end{tikzpicture}
\legende{Les trois phases physiologiques lors du dégonflage du brassard. A : Artère totalement écrasée (silence). B : Artère partiellement comprimée, flux turbulent pendant la systole uniquement $\rightarrow$ bruits de Korotkoff (apparition = MAX, disparition = MIN). C : Artère perméable en permanence, flux laminaire (silence).}
\end{popfigure}

\begin{methode}[Réaliser et interpréter une mesure de tension artérielle]
\begin{enumerate}
    \item \textbf{Préparer le sujet :} Repos 10 min, position assise, bras nu au niveau du cœur, pas de café/tabac/exercice récent, vessie vide.
    \item \textbf{Choisir le brassard :} Largeur $\approx$ 40\% de la circonférence du bras (brassard adulte standard $\approx$ 12--13 cm de large). Un brassard trop petit surestime la TA ; trop grand la sous-estime.
    \item \textbf{Gonfler au-dessus de la MAX :} Jusqu'à disparition du pouls radial + 20--30 mmHg (soit 2--3 cmHg).
    \item \textbf{Dégonfler lentement :} 2--3 mmHg/s (0,2--0,3 cmHg/s). Trop rapide $\rightarrow$ erreur de lecture.
    \item \textbf{Noter :} Pression au 1\textsuperscript{er} bruit clair (Phase I Korotkoff) = \textbf{MAX}. Pression à la disparition totale des bruits (Phase V) = \textbf{MIN}.
    \item \textbf{Répéter :} Faire 2 à 3 mesures à 1--2 min d'intervalle ; garder la moyenne des deux dernières si la première est plus élevée (phénomène d'alarme).
    \item \textbf{Noter le résultat :} Ex. : \textbf{TA = 13,2 / 8,1 cmHg} (ou 132/81 mmHg), bras droit, position assise, heure.
\end{enumerate}
\end{methode}

\begin{savaistu}[Les bruits de Korotkoff : une découverte russe]
C'est le chirurgien militaire russe \textbf{Nikolaï Korotkoff} qui, en 1905, décrit pour la première fois les bruits entendus au stéthoscope lors du dégonflage du brassard. Il identifie 5 phases (I à V). En pratique clinique courante (et au programme), on retient :
\begin{itemize}
    \item \textbf{Phase I :} Apparition de bruits clairs, rythmiques (toniques) $\rightarrow$ \textbf{Pression Systolique (MAX)}.
    \item \textbf{Phase V :} Disparition brutale et définitive des bruits $\rightarrow$ \textbf{Pression Diastolique (MIN)}.
\end{itemize}
Les phases II (souffle), III (bruits plus forts) et IV (étouffement) sont des nuances acoustiques liées à l'évolution de la turbulence, non utilisées pour la définition standard de la MIN aujourd'hui (sauf cas particuliers comme l'insuffisance aortique où la phase IV est préférée).
\end{savaistu}

\begin{exemplebox}[Conversion d'unités et lecture courante]
Un tensiomètre électronique affiche souvent en \textbf{mmHg}.
\begin{itemize}
    \item Lecture : \textbf{130 / 80 mmHg}.
    \item Conversion en cmHg (division par 10) : \textbf{13,0 / 8,0 cmHg}.
    \item Notation au dossier médical (contexte camerounais) : \textbf{TA = 13/8 cmHg}.
\end{itemize}
\end{exemplebox}

\begin{exoresolu}[Analyse d'un protocole de mesure]
\textbf{Énoncé :} Un élève de Terminale mesure la tension artérielle d'un camarade au repos. Il gonfle le brassard à 15 cmHg, puis dégonfle rapidement (environ 10 mmHg/s). Il note l'apparition des bruits à 12,5 cmHg et leur disparition à 7,5 cmHg. Il note : « TA = 20 cmHg ».
\begin{enumerate}
    \item Identifier trois erreurs méthodologiques dans ce protocole.
    \item Donner la valeur correcte de la tension artérielle avec l'unité et la notation adaptées.
    \item Expliquer pourquoi la vitesse de dégonflage influence la précision de la mesure.
\end{enumerate}
\tcblower
\textbf{Solution détaillée :}
\begin{enumerate}
    \item \textbf{Erreurs identifiées :}
    \begin{itemize}
        \item \textbf{Pression de gonflage insuffisante :} 15 cmHg est trop bas si la MAX réelle est 12,5 cmHg (marge < 20--30 mmHg = 2--3 cmHg). Risque de ne pas occlure totalement l'artère au départ, faussant le point de départ.
        \item \textbf{Vitesse de dégonflage trop rapide :} 10 mmHg/s = 1 cmHg/s $\gg$ 0,2--0,3 cmHg/s recommandés. Cela entraîne une erreur de lecture sur le manomètre (inertie de l'œil/mercure) et une sous-estimation de la MAX / surestimation de la MIN.
        \item \textbf{Notation erronée :} « TA = 20 cmHg » est absurde (addition des deux valeurs ? valeur unique ?). La tension se note \textbf{MAX/MIN}.
    \end{itemize}
    \item \textbf{Valeur correcte :} \textbf{TA = 12,5 / 7,5 cmHg} (ou 125/75 mmHg). Notation : « 12,5/7,5 cmHg » ou « 125/75 mmHg ». Préciser le bras (droit/gauche) et la position.
    \item \textbf{Justification vitesse :} Une vitesse lente (2--3 mmHg/s) permet au mercure (ou à l'aiguille) de suivre fidèlement la pression intra-brassard et à l'observateur de synchroniser l'événement auditif (bruit) avec la lecture visuelle. Une vitesse rapide crée un décalage temporel : on entend le bruit \emph{après} avoir lu la valeur, ou inversement, introduisant une erreur systématique de plusieurs mmHg.
\end{enumerate}
\end{exoresolu}

\begin{aretenir}[Essentiel sur la mesure de la PA]
\begin{itemize}
    \item La pression artérielle (TA) = pression du sang sur la paroi artérielle. Deux valeurs : \textbf{MAX (systolique)} et \textbf{MIN (diastolique)}.
    \item Unité : \textbf{cmHg} (ou mmHg). Notation : \textbf{MAX/MIN} (ex. 13/8 cmHg).
    \item Mesure de référence : \textbf{Méthode auscultatoire (Korotkoff)} : brassard + stéthoscope sur artère humérale.
    \item \textbf{Protocole rigoureux :} Repos, brassard au niveau du cœur, gonflage $>$ MAX + 2--3 cmHg, dégonflage \textbf{lent} (2--3 mmHg/s).
    \item \textbf{Interprétation :} 1\textsuperscript{er} bruit = MAX ; Disparition bruits = MIN.
    \item Conditions de fiabilité : moyenne de 2--3 mesures, brassard adapté, sujet calme.
\end{itemize}
\end{aretenir}

\courssub{Automatisation et automesure : tensions artérielles « blanches » et « masquées »}

Les tensiomètres électroniques (oscillométriques) sont désormais omniprésents (pharmacies, domiciles, consultations). Ils ne détectent pas les bruits de Korotkoff mais les \textbf{oscillations de pression} dans le brassard transmises par la paroi artérielle. L'algorithme interne calcule la MAX et la MIN (souvent en surestimant légèrement la MIN par rapport à l'auscultation).

\begin{attention}[Automesure et pièges de l'électronique]
\begin{itemize}
    \item \textbf{Effet « blouse blanche » :} TA élevée au cabinet (stress) mais normale à domicile. L'automesure (règle des 3 : 3 mesures le matin, 3 le soir, 3 jours de suite) est le standard pour confirmer une hypertension.
    \item \textbf{Hypertension masquée :} TA normale au cabinet mais élevée à domicile (risque cardiovasculaire majoré). Seule l'automesure ou la MAPA (Mesure Ambulatoire de la Pression Artérielle sur 24h) la dépiste.
    \item \textbf{Valider son appareil :} Seuls les appareils validés par des protocoles internationaux (ESH, AAMI, BHS) sont fiables. Vérifier la mention « cliniquement validé » sur la notice. Les appareils poignet sont moins fiables que les appareils bras (positionnement difficile au niveau du cœur).
\end{itemize}
\end{attention}

\begin{exercice}[Entraînement au Baccalauréat - Analyse de document]
\textbf{Document :} Un médecin mesure la TA d'un patient de 55 ans au cabinet. Il utilise un sphygmomanomètre à mercure. Le brassard est gonflé à 180 mmHg. Le dégonflage se fait à 2 mmHg/s. Les bruits apparaissent à 148 mmHg et disparaissent à 96 mmHg. Le patient est assis, bras soutenu. Deux mesures supplémentaires à 2 min d'intervalle donnent : 146/94 et 144/92 mmHg.
\begin{enumerate}
    \item Exprimer la tension artérielle moyenne de ce patient en cmHg, avec la notation conventionnelle.
    \item Classer cette tension selon la classification OMS (normale, normale élevée, hypertension grade 1, 2, 3).
    \item Justifier l'intérêt de réaliser trois mesures successives.
    \item Le médecin suspecte une « hypertension artérielle blanche ». Quelle conduite tenir pour lever le doute ?
\end{enumerate}
\end{exercice}

\begin{corrige}[Exercice Entraînement Bac]
\begin{enumerate}
    \item \textbf{Moyennes :} MAX = $(148 + 146 + 144)/3 = 146$ mmHg. MIN = $(96 + 94 + 92)/3 = 94$ mmHg.
    Conversion : \textbf{TA = 14,6 / 9,4 cmHg} (notation : 14,6/9,4 cmHg).
    \item \textbf{Classification OMS 1999 (toujours au programme) :}
    \begin{itemize}
        \item Optimale : $<$ 120/80 mmHg ($<$ 12/8 cmHg).
        \item Normale : 120--129 / 80--84 mmHg (12--12,9 / 8--8,4 cmHg).
        \item Normale élevée : 130--139 / 85--89 mmHg (13--13,9 / 8,5--8,9 cmHg).
        \item \textbf{HTA Grade 1 (légère) : 140--159 / 90--99 mmHg (14--15,9 / 9--9,9 cmHg).}
        \item HTA Grade 2 (modérée) : 160--179 / 100--109 mmHg.
        \item HTA Grade 3 (sévère) : $\ge$ 180 / $\ge$ 110 mmHg.
    \end{itemize}
    Le patient (146/94 mmHg) est en \textbf{Hypertension Artérielle Grade 1}.
    \item \textbf{Intérêt des 3 mesures :} La première mesure est souvent plus élevée (stress, effet d'alerte, non-repos suffisant). La moyenne des deux dernières (ou des trois) est plus représentative de la TA basale. Cela évite de surestimer la TA et de traiter à tort.
    \item \textbf{Conduite face à HTA blanche suspectée :} Prescrire une \textbf{automesure tensionnelle (AMT)} à domicile (règle des 3) ou une \textbf{MAPA (Mesure Ambulatoire de la Pression Artérielle sur 24h)}. Si les valeurs hors cabinet sont normales ($<$ 135/85 mmHg en automesure, $<$ 130/80 mmHg en MAPA jour), c'est une HTA blanche (pas de traitement médicamenteux immédiat, mais surveillance). Si elles restent élevées : HTA soutenue $\rightarrow$ prise en charge.
\end{enumerate}
\end{corrige}

Cette section a posé les bases physiques et techniques de la pression artérielle. La suite logique (Section 5) consistera à comprendre \emph{physiologiquement} ce qui fait varier cette pression (débit cardiaque, résistance périphérique, volume sanguin, viscosité) et à analyser les pathologies majeures que sont l'hypertension artérielle (HTA) et l'hypotension, fléaux de santé publique au Cameroun.

\coursec{Variations de la pression artérielle, hypertension et hypotension}

Après avoir défini la pression artérielle (PA) et décrit sa mesure au brassard (section précédente), nous devons maintenant comprendre \textbf{quoi} fait varier cette pression d'un instant à l'autre et d'un individu à l'autre, et \textbf{quelles en sont les conséquences pathologiques} lorsque la régulation échoue. C'est l'objet de cette dernière section.

\courssub{Paramètres physiologiques déterminant la pression artérielle}

\begin{experience}[Modélisation hydraulique de la circulation sanguine]
\textbf{Matériel :} Une pompe à eau (cœur), un tuyau rigide (aorte), un réseau de tuyaux fins à diamètre variable (artérioles), un manomètre (pression), un réservoir (volume sanguin), du colorant (traceur).
\textbf{Protocole :}
\begin{enumerate}
    \item Mettre la pompe en marche à débit constant. Mesurer la pression dans le tuyau principal.
    \item Augmenter la fréquence de la pompe (fréquence cardiaque) puis la force de chaque impulsion (volume d'éjection). Noter la pression.
    \item Revenir au régime initial. Rétrécir les tuyaux fins (vasoconstriction des artérioles). Noter la pression.
    \item Revenir au régime initial. Ajouter de l'eau dans le réservoir (augmentation du volume sanguin). Noter la pression.
\end{enumerate}
\textbf{Observations :} La pression augmente dans les trois cas : (1) quand le débit de la pompe augmente, (2) quand le calibre des tuyaux fins diminue, (3) quand le volume total de liquide augmente.
\textbf{Interprétation :} La pression artérielle est le résultat de l'interaction entre le \textbf{débit cardiaque} (pompe) et les \textbf{résistances périphériques} (calibre des artérioles), modulée par le \textbf{volume sanguin} (remplissage).
\end{experience}

\begin{propriete}[Relation fondamentale de l'hémodynamique (admise)]
La pression artérielle moyenne (PAM) est proportionnelle au produit du débit cardiaque ($Q_c$) par les résistances vasculaires périphériques totales ($R_{vp}$) :
\[
\text{PAM} \approx Q_c \times R_{vp}
\]
\emph{Rappel :} Le débit cardiaque $Q_c = \text{Fréquence cardiaque (FC)} \times \text{Volume d'éjection systolique (VES)}$.
Les résistances périphériques dépendent principalement du \textbf{diamètre des artérioles} (loi de Poiseuille : $R \propto 1/r^4$) et de l'\textbf{élasticité des grosses artères} (compliance artérielle).
\end{propriete}

\begin{popfigure}
\begin{tikzpicture}[
    node distance=1.5cm and 2cm,
    box/.style={draw=popdark, thick, rounded corners, fill=popcream, minimum width=3cm, minimum height=1.2cm, align=center},
    arrow/.style={-Stealth, thick, popblue},
    factor/.style={draw=poporange, thick, rounded corners, fill=poporangeL, minimum width=2.5cm, minimum height=1cm, align=center, font=\small}
]
% Nœuds principaux
\node[box, align=center] (coeur){CŒUR\\Débit cardiaque ($Q_c$)\\\small FC $\times$ VES};
\node[box, right=of coeur, align=center] (arteres){ARTÈRES / ARTÉRIOLES\\Résistances périphériques ($R_{vp}$)\\\small Calibre $\downarrow$ $\Rightarrow$ $R_{vp} \uparrow$};
\node[box, below=of arteres, align=center] (reins){REINS\\Volume sanguin ($V_s$)\\\small Sel $\uparrow$ $\Rightarrow$ Eau $\uparrow$ $\Rightarrow$ $V_s \uparrow$};
\node[box, below=of coeur, align=center] (pa){PRESSION ARTÉRIELLE (PA)\\\small PAM $\approx Q_c \times R_{vp}$};

% Flèches
\draw[arrow] (coeur) -- node[above, font=\small] {Éjecte le sang} (arteres);
\draw[arrow] (arteres) -- node[right, font=\small] {S'oppose au flux} (pa);
\draw[arrow] (coeur) -- node[left, font=\small] {Génère la pression} (pa);
\draw[arrow] (reins) -- node[right, font=\small] {Régule le remplissage} (arteres);
\draw[arrow] (reins) -| node[near start, below, font=\small] {Rétention Na$^+$/Eau} (coeur);

% Facteurs modulateurs
\node[factor, above=of coeur, xshift=-1.5cm] (fc) {Fréquence cardiaque};
\node[factor, above=of coeur, xshift=1.5cm] (ves) {Force de contraction};
\node[factor, above=of arteres, align=center] (calibre){Diamètre artérioles\\ (Système nerveux, Hormones)};
\node[factor, below=of reins] (sel) {Apport sodé (Sel)};

\draw[arrow, poporange] (fc) -- (coeur);
\draw[arrow, poporange] (ves) -- (coeur);
\draw[arrow, poporange] (calibre) -- (arteres);
\draw[arrow, poporange] (sel) -- (reins);
\end{tikzpicture}
\legende{Schéma fonctionnel des déterminants de la pression artérielle. Le cœur (débit) et les artérioles (résistances) sont les deux moteurs immédiats ; les reins régulent le volume sanguin sur le long terme via l'équilibre sodé.}
\end{popfigure}

\courssub{Variations physiologiques de la pression artérielle}

La PA n'est pas une constante ; elle s'adapte en permanence aux besoins de l'organisme. L'analyse de documents expérimentaux (courbes tensionnelles) met en évidence quatre situations types.

\begin{experience}[Variation de la PA au cours d'un effort physique (document type Bac)]
\textbf{Document :} Enregistrement continu de la PA systolique (PAS) et diastolique (PAD) chez un sujet sain au repos, pendant un exercice sur cyclo-ergomètre (charge croissante), et pendant la récupération.
\textbf{Résultats (valeurs moyennes) :}
\begin{center}
\begin{tabularx}{\linewidth}{|l|X|X|X|}\hline
\textbf{Phase} & \textbf{PAS (cmHg)} & \textbf{PAD (cmHg)} & \textbf{Interprétation physiologique} \\ \hline
Repos & 12 & 8 & Régulation basale (SNA orthosympathique modéré) \\ \hline
Effort modéré & 14 & 8 & $\uparrow$ Débit cardiaque (FC $\uparrow$, VES $\uparrow$) ; $\downarrow$ Résistances périphériques (vasodilatation musculaire) $\Rightarrow$ $\uparrow$ PAS, PAD stable \\ \hline
Effort intense & 18-20 & 8-9 & $\uparrow \uparrow$ Débit cardiaque ; $\uparrow$ Résistances (vasoconstriction zones non actives) $\Rightarrow$ $\uparrow \uparrow$ PAS, $\uparrow$ PAD \\ \hline
Récupération (5 min) & 11 & 7 & $\downarrow$ FC, $\downarrow$ Catécholamines, vasodilatation résiduelle $\Rightarrow$ PA $<$ repos (hypotension post-effort transitoire) \\ \hline
\end{tabularx}
\end{center}
\textbf{Conclusion :} L'effort physique élève la PA systolique proportionnellement à l'intensité (augmentation du débit cardiaque), tandis que la diastolique varie peu ou augmente modérément. La chute post-effort témoigne de l'inertie de la vasodilatation métabolique.
\emph{Note :} 1 cmHg = 10 mmHg. Les seuils cliniques usuels sont exprimés en cmHg au Cameroun (ex : 14/9 cmHg = 140/90 mmHg).
\end{experience}

\begin{popfigure}
\begin{tikzpicture}
\begin{axis}[
    width=\linewidth, height=8cm,
    axis lines=middle,
    xlabel={Temps (min)}, ylabel={Pression artérielle (cmHg)},
    xmin=0, xmax=30, ymin=5, ymax=22,
    xtick={0,10,15,20,22,30}, 
    xticklabels={Repos, D\'ebut effort, Effort intense, Fin effort, D\'ebut r\'ecup., Fin r\'ecup.},
    ytick={8,12,14,16,18,20},
    legend style={at={(0.02,0.98)}, anchor=north west, fill=popcream, draw=popdark},
    popgrid
]
% PAS
\addplot[popcurve, popblue, very thick, mark=*] coordinates {
    (0,12) (5,12) (10,14) (15,19) (20,19) (22,11) (25,11.5) (30,12)
};
\addlegendentry{PAS (Systolique)}
% PAD
\addplot[popcurve, poporange, very thick, mark=square*] coordinates {
    (0,8) (5,8) (10,8) (15,8.5) (20,9) (22,7) (25,7.5) (30,8)
};
\addlegendentry{PAD (Diastolique)}
% Zones
\draw[fill=popblueL, opacity=0.3, draw=none] (axis cs:10,5) rectangle (axis cs:20,22);
\node[popblue, font=\small, rotate=90] at (axis cs:15,21) {Effort physique};
\end{axis}
\end{tikzpicture}
\legende{Variation type de la pression artérielle systolique (bleu) et diastolique (orange) au repos, pendant un effort physique croissant et lors de la récupération. L'écart entre PAS et PAD (pression différentielle) s'élargit à l'effort.}
\end{popfigure}

\begin{aretenir}[Autres variations physiologiques]
\begin{itemize}
    \item \textbf{Sommeil :} Diminution de la PA (PAS $\approx$ 10-11 cmHg, PAD $\approx$ 6-7 cmHg) liée à la prédominance parasympathique, baisse du métabolisme et de la température.
    \item \textbf{Émotions / Stress aigu :} Montée brutale de la PA (PAS peut dépasser 18-20 cmHg) par décharge d'adrénaline/ noradrénaline ($\uparrow$ FC, $\uparrow$ VES, vasoconstriction cutanée/splanchnique). C'est l'\textbf{effet « blouse blanche »} en consultation.
    \item \textbf{Âge :} Augmentation progressive et continue de la PA systolique (raidissement artériel, perte d'élasticité, athérosclérose) ; la diastolique augmente jusqu'à $\approx$ 50-60 ans puis stagne ou baisse (augmentation de la pression différentielle).
\end{itemize}
\end{aretenir}

\begin{popfigure}
\begin{tikzpicture}
\begin{axis}[
    width=\linewidth, height=7cm,
    axis lines=left,
    xlabel={Âge (années)}, ylabel={Pression artérielle moyenne (cmHg)},
    xmin=0, xmax=80, ymin=8, ymax=18,
    xtick={0,20,40,60,80}, ytick={8,10,12,14,16,18},
    legend style={at={(0.02,0.98)}, anchor=north west, fill=popcream, draw=popdark},
    popgrid
]
\addplot[popcurve, popblue, very thick] coordinates {
    (0,9) (10,10) (20,11) (30,11.5) (40,12.5) (50,13.5) (60,14.5) (70,15.5) (80,16)
};
\addlegendentry{PAS (Systolique) : hausse continue}
\addplot[popcurve, poporange, very thick] coordinates {
    (0,6) (10,7) (20,7.5) (30,8) (40,8.5) (50,9) (60,9) (70,8.5) (80,8)
};
\addlegendentry{PAD (Diastolique) : plateau puis baisse}
\addplot[popcurve, popgreen, very thick, dashed] coordinates {
    (0,7) (10,8) (20,8.7) (30,9.2) (40,9.8) (50,10.5) (60,11) (70,11) (80,10.7)
};
\addlegendentry{PAM (Moyenne) : hausse modérée}
\end{axis}
\end{tikzpicture}
\legende{Évolution moyenne de la pression artérielle avec l'âge. La perte d'élasticité des grosses artères (rigidité aortique) explique la divergence croissante entre systolique et diastolique chez le sujet âgé.}
\end{popfigure}

\courssub{Le rôle des reins et du volume sanguin : le facteur sel}

\begin{experience}[Corrélation épidémiologique sel / PA (Étude INTERSALT, données camerounaises)]
\textbf{Observation :} Comparaison de populations à apport sodé très différent.
\begin{itemize}
    \item Population traditionnelle rurale (Cameroun profond, Yanomami d'Amazonie) : apport NaCl $<$ 1 g/j $\rightarrow$ PA moyenne $\approx$ 10/6 cmHg, pas d'augmentation avec l'âge, HTA quasi absente.
    \item Population urbaine (Yaoundé, Douala) : apport NaCl 8-12 g/j (sel de cuisine, cubes d'assaisonnement, conserves, pain) $\rightarrow$ PA moyenne $\approx$ 13/8 cmHg, prévalence HTA $\approx$ 30\%.
    \item Essais cliniques (DASH-sodium) : réduction de 9 à 3 g/j de sel $\rightarrow$ baisse de la PAS de 5 à 7 mmHg (0,5-0,7 cmHg) chez normotendus, 10-12 mmHg chez hypertendus.
\end{itemize}
\textbf{Interprétation mécanistique :} Excès de Na$^+$ $\rightarrow$ $\uparrow$ Osmolarité plasmatique $\rightarrow$ Soif + $\uparrow$ ADH (vasopressine) + $\uparrow$ Aldostérone $\rightarrow$ Rétention d'eau $\rightarrow$ $\uparrow$ Volume sanguin circulant ($V_s$) $\rightarrow$ $\uparrow$ Volume d'éjection $\rightarrow$ $\uparrow$ Débit cardiaque $\rightarrow$ $\uparrow$ PA. À long terme, l'hypervolémie entretient une hypertrophie artériolaire (remodelage vasculaire) qui fixe l'hypertension.
\end{experience}

\begin{attention}[Piège fréquent : Confondre corrélation et causalité unique]
Le sel n'est pas la \emph{seule} cause de l'HTA, mais c'est un \textbf{facteur majeur, modifiable et prouvé}. Dire « je ne mange pas de sel, donc je n'aurai pas d'HTA » est faux (facteurs génétiques, obésité, âge, alcool...). Dire « le sel n'a pas d'effet car mon oncle en mange beaucoup et va bien » est un raisonnement anecdote non scientifique (variabilité interindividuelle de la « sensibilité au sel »).
\end{attention}

\courssub{L'hypertension artérielle (HTA) : le « tueur silencieux »}

\begin{definition}[Hypertension Artérielle (HTA)]
L'HTA est une \textbf{élévation permanente} de la pression artérielle au repos, confirmée par \textbf{plusieurs mesures} sur plusieurs consultations :
\begin{itemize}
    \item \textbf{HTA confirmée (cabinet) :} PAS $\ge$ \textbf{14 cmHg} (140 mmHg) \textbf{ET/OU} PAD $\ge$ \textbf{9 cmHg} (90 mmHg).
    \item \textbf{HTA grade 1 (modérée) :} 14-15,9 / 9-9,9 cmHg.
    \item \textbf{HTA grade 2 (sévère) :} 16-17,9 / 10-10,9 cmHg.
    \item \textbf{HTA grade 3 (très sévère) :} $\ge$ 18 / $\ge$ 11 cmHg.
\end{itemize}
\emph{Note :} En automesure (domicile), les seuils sont abaissés de 5 mmHg (0,5 cmHg) : $\ge$ 13,5 / 8,5 cmHg.
\end{definition}

\begin{savaistu}[L'HTA au Cameroun : un défi de santé publique]
Selon l'OMS et les enquêtes STEPS (2017-2018), \textbf{environ 30\% des adultes camerounais (1 sur 3) sont hypertendus}. Plus de \textbf{la moitié l'ignorent} (asymptomatique). Seule une minorité des traités est contrôlée. L'HTA est la première cause de consultation en cardiologie aux urgences de l'Hôpital Central de Yaoundé et de l'Hôpital Laquintinie de Douala. Elle est la première cause d'AVC et d'insuffisance cardiaque dans nos services.
\end{savaistu}

\begin{propriete}[Signes cliniques : une maladie le plus souvent silencieuse]
L'HTA essentielle (90-95\% des cas, sans cause unique identifiée) est \textbf{asymptomatique} pendant des années.
Les signes (céphalées matinales occipitales, vertiges, acouphènes, saignements de nez \emph{épistaxis}, fatigue visuelle) sont \textbf{inconstants, non spécifiques et tardifs}. Ils apparaissent souvent seulement lors de complications (crise hypertensive, AVC).
$\Rightarrow$ \textbf{Seule la mesure tensionnelle répétée permet le diagnostic.}
\end{propriete}

\begin{aretenir}[Facteurs de risque et causes de l'HTA]
\begin{center}
\begin{tabularx}{\linewidth}{|l|X|}\hline
\textbf{Type} & \textbf{Facteurs} \\ \hline
\textbf{Modifiables (Hygiène de vie)} & Excès de sel (NaCl) $\uparrow$ ; Obésité / Surpoids (IMC $\ge$ 25) ; Sédentarité ; Alcool (chronique) ; Tabac (rigidité artérielle) ; Stress chronique ; Apport potassique insuffisant (peu fruits/légumes). \\ \hline
\textbf{Non modifiables} & Âge ($\uparrow$ rigidité artérielle) ; Hérédité (antécédents parents 1$^{\text{er}}$ degré) ; Sexe (homme $>$ femme avant 60 ans, puis inversion) ; Ethnie (prédisposition génétique chez l'homme noir). \\ \hline
\textbf{Secondaires (5-10\%)} & Maladie rénale (sténose artère rénale, glomérulonéphrite) ; Endocrinienne (phéochromocytome, hyperaldostéronisme, syndrome de Cushing) ; Médicamenteuse (corticoïdes, AINS, pilule œstro-progestative, décongestionnants). \\ \hline
\end{tabularx}
\end{center}
\end{aretenir}

\begin{propriete}[Complications majeures de l'HTA non contrôlée (lésions d'organes cibles)]
L'hyperpression mécanique chronique détruit les parois vasculaires :
\begin{itemize}
    \item \textbf{Cerveau :} \textbf{Accident Vasculaire Cérébral (AVC)} — hémorragie (rupture micro-anévrysme charcot-bouchard) ou ischémie (athérothrombose). 1$^{\text{re}}$ cause de décès par HTA au Cameroun.
    \item \textbf{Cœur :} \textbf{Insuffisance cardiaque} (hypertrophie ventriculaire gauche $\rightarrow$ dysfonction diastolique puis systolique) ; \textbf{Infarctus du myocarde} (athérosclérose coronarienne accélérée).
    \item \textbf{Reins :} \textbf{Insuffisance rénale chronique} (néphro-angiosclérose hyalinique $\rightarrow$ protéinurie $\rightarrow$ dialyse).
    \item \textbf{Yeux :} \textbf{Rétinopathie hypertensive} (rétrécissement artérioles, hémorragies, exsudats, papillite) $\rightarrow$ cécité.
    \item \textbf{Artères :} Anévrysme de l'aorte (dissection), artériopathie des membres inférieurs.
\end{itemize}
\end{propriete}

\begin{methode}[Exploiter un résultat tensionnel (compétence Bac)]
\begin{enumerate}
    \item \textbf{Vérifier les conditions de mesure :} Repos 5-10 min, position assise, bras au niveau du cœur, brassard adapté, pas de café/tabac/effort 30 min avant, 3 mesures à 1-2 min d'intervalle (moyenne des 2 dernières).
    \item \textbf{Comparer aux seuils :} PAS $\ge$ 14 \textbf{ET/OU} PAD $\ge$ 9 cmHg (cabinet).
    \item \textbf{Exiger la répétition :} Au moins \textbf{2 consultations} espacées de 1 à 4 semaines (sauf urgence hypertensive $>$ 18/11 avec signes).
    \item \textbf{Rechercher le contexte :} Âge, facteurs de risque, traitement en cours, automesure (règle des 3 : 3 mesures matin/soir pendant 3 jours).
    \item \textbf{Conclure :} « HTA confirmée grade X » ou « PA normale / limite » ou « Effet blouse blanche (PA cabinet $\ge$ 14/9, PA domicile $<$ 13,5/8,5) ».
\end{enumerate}
\end{methode}

\begin{exoresolu}[Interprétation d'un bilan tensionnel]
\textbf{Énoncé :} M. K., 52 ans, commerçant à Douala, consulte pour « contrôle ». Il ne prend aucun traitement. Mesures au cabinet (repos 10 min, brassard adapté) :
M1 : 15,2 / 9,5 cmHg \quad M2 : 14,8 / 9,2 cmHg \quad M3 : 14,5 / 9,0 cmHg.
Il rapporte une automesure domicile (3 jours, règle des 3) : Moyenne matin 13,8 / 8,6 cmHg ; Moyenne soir 14,2 / 8,8 cmHg.
IMC = 29 kg/m$^2$. Alcool : 2 verres vin de palme/jour. Sel : « je sale peu mais je mange beaucoup de cubes Maggi et poisson fumé ». Père décédé d'AVC à 60 ans.
\textbf{Analyse et conclusion.}

\textbf{Solution détaillée :}
\begin{enumerate}
    \item \textbf{Conditions de mesure :} Respectées (repos, 3 mesures, moyenne des 2 dernières). Moyenne cabinet = (14,8+14,5)/2 = \textbf{14,65 / 9,1 cmHg} pour la systolique/diastolique (on moyenne séparément).
    \item \textbf{Seuils cabinet :} PAS 14,65 $\ge$ 14 \textbf{ET} PAD 9,1 $\ge$ 9 $\rightarrow$ \textbf{HTA Grade 1 (limite Grade 2 pour la diastolique)}.
    \item \textbf{Automesure :} Moyenne globale $\approx$ 14,0 / 8,7 cmHg. Seuil domicile $\ge$ 13,5 / 8,5. \textbf{Confirmée} (pas d'effet blouse blanche majeur).
    \item \textbf{Facteurs de risque :} Obésité (IMC 29), Alcool (excès), Sel caché (cubes, poisson fumé = NaCl), Hérédité (père AVC), Âge, Sexe masculin.
    \item \textbf{Conclusion :} \textbf{HTA essentielle Grade 1 confirmée, à risque cardiovasculaire élevé (obésité + antécédent familial + âge).} Prise en charge : règles hygiéno-diététiques strictes (régime DASH, arrêt alcool, marche 30 min/j), réévaluation à 3 mois, traitement médicamenteux si PA non contrôlée.
\end{enumerate}
\end{exoresolu}

\courssub{Moyens de lutte contre l'hypertension artérielle}

La prise en charge associe \textbf{mesures hygiéno-diététiques (premier niveau, indispensables)} et \textbf{traitement médicamenteux (si objectifs non atteints ou risque élevé)}.

\begin{aretenir}[Stratégie thérapeutique de l'HTA (recommandations OMS / Société Camerounaise de Cardiologie)]
\begin{enumerate}
    \item \textbf{Réduction de l'apport sodé :} $<$ 5 g NaCl/j (1 c. à café rase). Supprimer sel de table, limiter cubes, conserves, charcuteries, poisson fumé/salé, pain. Privilégier épices, ail, oignon, citron. \textbf{Cible :} $\downarrow$ 5-10 mmHg PAS.
    \item \textbf{Lutte contre le surpoids/obésité :} IMC cible $<$ 25 kg/m$^2$. Perte de 5-10\% du poids $\rightarrow$ $\downarrow$ 5-20 mmHg PAS.
    \item \textbf{Activité physique régulière :} Endurance modérée (marche rapide, vélo, natation) 30 min/j, 5-7 j/sem. $\downarrow$ 4-9 mmHg PAS.
    \item \textbf{Alcool :} Limitation stricte $\le$ 2 verres/j (homme), $\le$ 1 verre/j (femme). Arrêt si possible. $\downarrow$ 2-4 mmHg.
    \item \textbf{Arrêt du tabac :} Indispensable pour risque vasculaire global (ne baisse pas directement la PA mais réduit complications).
    \item \textbf{Régime type DASH / Méditerranéen :} Riches en fruits/légumes (K$^+$, Mg$^{2+}$, fibres), produits laitiers allégés, pauvres en graisses saturées.
    \item \textbf{Traitement antihypertenseur :} Si PA $\ge$ 14/9 malgré 3 mois d'hygiène de vie (ou d'emblée si PA $\ge$ 16/10 ou risque très élevé). Classes principales : IEC / ARA2 (1$^{\text{re}}$ intention souvent), Diurétiques thiazidiques, Calcium-antagonistes, Bêta-bloquants. \textbf{Observance à vie indispensable.}
    \item \textbf{Contrôles :} Consultation tous les 3-6 mois si contrôlée ; automesure mensuelle. Bilan rénal (créatininémie, clairance), ionogramme, ECG, fond d'œil annuels.
\end{enumerate}
\end{aretenir}

\courssub{L'hypotension artérielle : une baisse permanente}

\begin{definition}[Hypotension Artérielle]
C'est une \textbf{baisse permanente} de la pression artérielle au repos :
\begin{itemize}
    \item \textbf{Hypotension « constitutionnelle » :} PAS $<$ \textbf{9 cmHg} (90 mmHg) \textbf{ET/OU} PAD $<$ \textbf{6 cmHg} (60 mmHg) chez un sujet asymptomatique, souvent jeune femme, sans pathologie. \emph{Bénigne, ne nécessite pas de traitement.}
    \item \textbf{Hypotension symptomatique / pathologique :} Même seuils, mais \textbf{avec signes cliniques} (fatigue majeure, vertiges, lipothymies, malaises à l'orthostatisme, vision trouble, sueurs froides).
\end{itemize}
\end{definition}

\begin{propriete}[Causes principales de l'hypotension symptomatique]
\begin{itemize}
    \item \textbf{Hypovolémie :} Déshydratation (diarrhée, vomissements, fièvre, insuffisance d'apport — fréquent chez le sujet âgé au Cameroun en saison sèche), Hémorragie (accouchement, traumatisme, ulcère), Brûlures.
    \item \textbf{Médicamenteuse :} Surdosage antihypertenseur (cause n°1 chez le traité), Diurétiques, Vasodilatateurs, Antidépresseurs, Neuroleptiques, Antiparkinsoniens.
    \item \textbf{Orthostatique (chute $\ge$ 2 cmHg PAS ou $\ge$ 1 cmHg PAD en station debout 3 min) :} Dysautonomie (diabète, Parkinson, vieillissement), Déshydratation, Médicaments (alpha-bloquants).
    \item \textbf{Autres :} Insuffisance surrénalienne (Addison), Insuffisance cardiaque sévère, Sepsis (choc septique), Embolie pulmonaire massive, Anaphylaxie.
\end{itemize}
\end{propriete}

\begin{methode}[Prise en charge de l'hypotension symptomatique]
\begin{enumerate}
    \item \textbf{Identifiez la cause :} Interrogatoire (médicaments, pertes liquides), Examen clinique (marbrures, veines jugulaires, pouls), Bilan (ionogramme, hémoglobine, cortisol si évoqué).
    \item \textbf{Mesures immédiates :} Allongement, surélévation des membres inférieurs (autotransfusion veineuse), arrêt du médicament suspect si iatrogénie.
    \item \textbf{Réhydratation :} Per os (eau salée, bouillons, solutés de réhydratation orale - SRO) ou IV (NaCl 0,9\% isotonique) si vomissements/obnubilation.
    \item \textbf{Adaptation du traitement antihypertenseur :} Réduction des doses, changement de classe, prise le soir (si hypotension matinale).
    \item \textbf{Conseils hygiéno-diététiques :} Fractionner les repas (éviter hypotension post-prandiale), lever progressif (assis 1 min au bord du lit avant de se lever), bas de contention (classe 2) si orthostatique, apport sodé \textbf{non restreint} (sauf HTA associée complexe).
\end{enumerate}
\end{methode}

\begin{exoresolu}[Cas clinique : Malaise chez un hypertendu traité]
\textbf{Énoncé :} Mme B., 68 ans, connue pour HTA traitée par Amlodipine 10 mg (calcium-antagoniste) + Hydrochlorothiazide 25 mg (diurétique) + Perindopril 5 mg (IEC). Elle consulte pour « malaises répétés en se levant le matin, fatigue ». PA cabinet : 10,5 / 6,0 cmHg (allongée) ; 8,5 / 5,5 cmHg (debout 3 min). Poids stable. Pas de fièvre, pas de diarrhée. Ionogramme : Na$^+$ 138, K$^+$ 3,2 mmol/L (hypokaliémie).
\textbf{Analyse et proposition.}

\textbf{Solution détaillée :}
\begin{enumerate}
    \item \textbf{Diagnostic :} \textbf{Hypotension orthostatique symptomatique iatrogène} (chute PAS 2 cmHg, PAD 0,5 cmHg + signes) \textbf{associée à une hypokaliémie} (diurétique thiazidique).
    \item \textbf{Mécanisme :} Diurétique $\rightarrow$ pertes Na$^+$/Eau $\rightarrow$ hypovolémie + pertes K$^+$ $\rightarrow$ hypokaliémie (aggravant la fatigue). IEC + Calcium-antagoniste $\rightarrow$ vasodilatation. Somme = PA trop basse pour la statique veineuse en orthostatisme.
    \item \textbf{Conduite à tenir :}
    \begin{itemize}
        \item Arrêt de l'hydrochlorothiazide (cause majeure hypovolémie + hypokaliémie).
        \item Supplémentation potassique (KCl per os) et alimentaire (bananes, avocat, épinards).
        \item Réévaluation PA à J+7 : si PA normale (12-13/7-8), maintenir bi-thérapie (IEC + Calcium-antagoniste).
        \item Si PA reste basse $\rightarrow$ réduire dose Amlodipine (5 mg) ou Perindopril (2,5 mg).
        \item Éducation : lever progressif, bas de contention, fractionner repas, boire 1,5-2 L/j.
    \end{itemize}
\end{enumerate}
\end{exoresolu}

\begin{attention}[Erreurs fréquentes à ne pas commettre au Bac]
\begin{itemize}
    \item \textbf{Confondre « PA basse » et « Hypotension pathologique » :} Une jeune fille de 20 ans à 10/6 cmHg sans symptôme est \textbf{normale}, pas « hypotendue malade ».
    \item \textbf{Oublier l'orthostatisme :} Ne pas mesurer la PA \textbf{allongée PUIS debout 1 min et 3 min} devant tout malaise, vertige, chute chez un sujet âgé ou traité.
    \item \textbf{Traiter un chiffre isolé :} Une PA à 15/9,5 cmHg \textbf{une seule fois} chez un patient stressé $\neq$ HTA. Il faut \textbf{répéter les mesures} (automesure ou MAPA).
    \item \textbf{Négliger le sel caché :} Au Cameroun, « je ne sale pas » ne veut pas dire « je mange peu de sel ». Cubes (Maggi, Jumbo), poisson fumé/salé, saucisses, pain, plats de rue sont les vraies sources.
    \item \textbf{Arrêter le traitement « car la PA est normale » :} L'HTA est chronique. L'arrêt $\rightarrow$ rebond tensionnel (effet rebond bêta-bloquants, clonidine) + réapparition complications. \textbf{Traitement à vie sauf avis médical contraire.}
\end{itemize}
\end{attention}

\begin{exercice}[Entraînement Bac - Variations de la PA]
\textbf{Document :} Évolution de la prévalence de l'HTA selon l'âge et le sexe au Cameroun (Enquête STEPS 2017).
\begin{center}
\begin{tabular}{|c|c|c|}\hline
\textbf{Tranche d'âge} & \textbf{Hommes (\% HTA)} & \textbf{Femmes (\% HTA)} \\ \hline
18-29 ans & 12,5 & 8,2 \\ \hline
30-44 ans & 28,4 & 22,1 \\ \hline
45-59 ans & 45,6 & 48,3 \\ \hline
60-69 ans & 58,2 & 62,7 \\ \hline
\end{tabular}
\end{center}
\textbf{Questions :}
\begin{enumerate}
    \item Décrire l'évolution de la prévalence de l'HTA avec l'âge pour les deux sexes.
    \item Comparer les deux sexes avant et après 60 ans. Proposer une explication physiologique.
    \item Un homme de 55 ans, IMC 31, consommateur de 3 verres d'alcool/jour, fumeur, se plaint de « maux de tête le matin ». Sa PA au cabinet est de 16,2 / 10,4 cmHg. Quelle est la classification de son HTA ? Citez 3 facteurs de risque modifiables présents. Quelle est la première mesure non médicamenteuse à lui conseiller ?
\end{enumerate}
\end{exercice}

\begin{corrige}[Exercice n°1]
\begin{enumerate}
    \item \textbf{Évolution :} Pour les deux sexes, la prévalence augmente \textbf{régulièrement et fortement} avec l'âge. Elle est multipliée par $\approx$ 5 entre 18-29 ans et 60-69 ans.
    \item \textbf{Comparaison :} Avant 60 ans (18-44 ans), les \textbf{hommes} sont plus touchés que les femmes. Dès la tranche 45-59 ans, la prévalence devient \textbf{supérieure chez les femmes} (48,3\% vs 45,6\%) et l'écart s'accentue après 60 ans (62,7\% vs 58,2\%).
    \textbf{Explication :} Avant la ménopause, les œstrogènes confèrent une protection vasculaire (vasodilatation, effet anti-inflammatoire, meilleur profil lipidique). La ménopause (vers 50 ans) entraîne une chute brutale des œstrogènes, une rigidification artérielle accélérée et une perte de cette protection, expliquant l'inversion de la prévalence en défaveur des femmes après 50 ans.
    \item \textbf{Classification :} PAS 16,2 $\ge$ 16 ET PAD 10,4 $\ge$ 10 $\rightarrow$ \textbf{HTA Grade 2 (sévère)}.
    \textbf{Facteurs modifiables :} Obésité (IMC 31), Alcool (3 verres/j > 2), Tabac. (Sédentarité probable, sel probable).
    \textbf{Première mesure non médicamenteuse :} \textbf{Perte de poids} (objectif IMC $<$ 25) associée à \textbf{réduction majeure de l'alcool} (arrêt idéal) et \textbf{arrêt du tabac}. Régime pauvre en sel (DASH) et reprise activité physique.
\end{enumerate}
\end{corrige}

\begin{aretenir}[Synthèse : La pression artérielle, un équilibre dynamique à surveiller]
\begin{itemize}
    \item \textbf{Physiologie :} PA = f(Débit cardiaque, Résistances périphériques, Volume sanguin). Elle s'adapte (effort, stress, sommeil, âge).
    \item \textbf{Pathologie majeure :} \textbf{HTA} ($\ge$ 14/9 cmHg permanent). Maladie silencieuse, fréquente (1/3 adultes au Cameroun), facteur de risque n°1 d'AVC, d'insuffisance cardiaque et rénale.
    \item \textbf{Facteurs clés :} Sel (volume), Obésité, Sédentarité, Alcool, Tabac, Âge, Hérédité.
    \item \textbf{Prise en charge :} Hygiène de vie (sel, poids, sport, alcool, tabac) \textbf{TOUJOURS} + Traitement médicamenteux (si nécessaire) \textbf{À VIE} + Observance + Contrôles réguliers.
    \item \textbf{Hypotension :} $<$ 9/6 cmHg. Souvent iatrogène (médicaments) ou hypovolémique. Symptomatique = prise en charge étiologique + réhydratation + adaptation traitement.
\end{itemize}
\end{aretenir}

\newpage

\coursec{Activité d'intégration}

\courssub{Objectifs de l'activité}

À l'issue de cette activité d'intégration, tu seras capable de :

\begin{itemize}
\item analyser et interpréter une courbe d'hyperglycémie provoquée par le glucose (HGPO) en la comparant aux valeurs de référence ;
\item exploiter des résultats d'analyses médicales (glycémie, bandelette urinaire) pour argumenter un diagnostic de diabète ;
\item interpréter une série de mesures de pression artérielle et identifier une hypertension artérielle ;
\item expliquer, à l'aide d'un schéma fonctionnel, le mécanisme physiologique en cause (régulation hormonale de la glycémie, paramètres de la pression artérielle) ;
\item rédiger une fiche de conseils de lutte et de prévention adaptée au contexte camerounais ;
\item travailler en équipe, argumenter oralement et confronter tes conclusions à celles des autres groupes.
\end{itemize}

\courssub{Situation}

Tu es membre d'une équipe médicale du centre de santé intégré d'un quartier de Yaoundé. Ce matin, deux patients se présentent à la consultation. Le médecin chef te confie, à toi et à ton groupe, l'exploitation de leurs dossiers.

\textbf{Dossier 1 — Mme A., 48 ans, vendeuse au marché.} Elle se plaint depuis plusieurs mois d'une \cle{soif permanente} (polydipsie), d'urines abondantes et fréquentes (polyurie) et d'un \cle{amaigrissement} malgré un appétit conservé. Le laboratoire du centre a réalisé une \cle{HGPO} : après une nuit de jeûne, la patiente ingère \SI{75}{g} de glucose dissous dans l'eau, et sa glycémie est mesurée toutes les 30 minutes pendant 2 heures. Une bandelette urinaire révèle la présence de \cle{glucose dans les urines} (glycosurie), normalement absent.

\begin{center}
\begin{tabularx}{\linewidth}{|l|X|X|X|X|X|}
\hline
\rowcolor{popblueL} Temps (min) & 0 (à jeun) & 30 & 60 & 90 & 120 \\
\hline Glycémie de Mme A. (g/L) & 1,50 & 2,10 & 2,45 & 2,40 & 2,30 \\
\hline Témoin sain (g/L) & 0,90 & 1,45 & 1,30 & 1,10 & 0,95 \\
\hline
\end{tabularx}
\end{center}

\begin{popfigure}
\begin{center}
\begin{tikzpicture}
\begin{axis}[width=0.85\linewidth,height=6.2cm,
xlabel={Temps (min)},ylabel={Glycémie (g/L)},
xmin=0,xmax=130,ymin=0,ymax=2.8,
xtick={0,30,60,90,120},ytick={0,0.5,1,1.26,1.5,2,2.5},
grid=both,grid style={popline!40},
legend style={at={(0.03,0.97)},anchor=north west,font=\small}]
\addplot[popcurve,color=popink,mark=*,mark size=1.5pt] coordinates {(0,1.5) (30,2.1) (60,2.45) (90,2.4) (120,2.3)};
\addplot[popcurve,color=popgreen,mark=square*,mark size=1.5pt] coordinates {(0,0.9) (30,1.45) (60,1.3) (90,1.1) (120,0.95)};
\addplot[domain=0:130,dashed,color=popmuted] {1.26};
\addplot[domain=0:130,dashed,color=poporange] {2.0};
\legend{Mme A., Témoin sain, Seuil à jeun (\SI{1,26}{g/L}), Seuil à 2 h (\SI{2,00}{g/L})}
\end{axis}
\end{tikzpicture}
\end{center}
\legende{Courbes d'HGPO de Mme A. et d'un témoin sain. Rappel des critères : glycémie à jeun $\geq \SI{1,26}{g/L}$ et/ou glycémie $\geq \SI{2,00}{g/L}$ deux heures après la charge en glucose.}
\end{popfigure}

\textbf{Dossier 2 — M. B., 55 ans, chauffeur de taxi.} Il consulte pour des \cle{maux de tête} fréquents, des vertiges et des bourdonnements d'oreilles. L'infirmier a mesuré sa pression artérielle au brassard, au repos et en position assise, à trois jours différents :

\begin{center}
\begin{tabularx}{\linewidth}{|l|X|X|X|}
\hline
\rowcolor{popblueL} Mesure & Jour 1 & Jour 2 & Jour 3 \\
\hline Pression systolique (cm de Hg) & 16 & 15 & 17 \\
\hline Pression diastolique (cm de Hg) & 10 & 11 & 10 \\
\hline
\end{tabularx}
\end{center}

Le questionnaire révèle une alimentation \cle{très salée} (cube d'assaisonnement en excès, poisson fumé quotidien), une consommation régulière d'alcool, l'absence totale d'activité physique (sédentarité liée au métier) et un père décédé d'une « crise » à 60 ans. On rappelle qu'un adulte au repos est considéré comme \cle{hypertendu} lorsque la pression artérielle est durablement supérieure ou égale à $14/9$ cm de Hg.

\textbf{Problème posé :} pour chaque patient, ton équipe doit exploiter les documents, poser un diagnostic argumenté, expliquer le mécanisme physiologique en cause et proposer une fiche de conseils de lutte et de prévention adaptée à la vie quotidienne du quartier.

\courssub{Consignes}

Par groupes de quatre à six élèves, traitez les deux dossiers en répondant aux questions suivantes. Chaque groupe préparera une restitution orale de cinq minutes par dossier.

\textbf{Dossier 1 — Mme A.}

\begin{enumerate}
\item Décris l'allure de la courbe d'HGPO de Mme A. et compare-la à celle du témoin sain (valeurs à jeun, pic, retour à la normale).
\item À l'aide des valeurs de référence, pose un diagnostic argumenté. Quel signe urinaire confirme ce diagnostic et pourquoi apparaît-il ?
\item Explique les trois symptômes de Mme A. (soif, urines abondantes, amaigrissement) en t'appuyant sur le mécanisme physiologique en cause.
\item Réalise le schéma fonctionnel de la régulation hormonale de la glycémie (pancréas, insuline, glucagon, foie) et indique sur ce schéma l'étape défaillante chez Mme A.
\item Rédige une fiche de conseils (hygiène alimentaire, activité physique, suivi médical) adaptée à une vendeuse de marché de Yaoundé.
\end{enumerate}

\textbf{Dossier 2 — M. B.}

\begin{enumerate}
\item Rappelle brièvement la technique de mesure de la pression artérielle au brassard et la signification des deux nombres d'une mesure « 16/10 ».
\item Compare les trois mesures de M. B. aux valeurs de référence. Pourquoi le diagnostic exige-t-il plusieurs mesures à des jours différents ? Conclus.
\item Identifie dans le questionnaire les facteurs favorisants et relie chacun d'eux à un paramètre physiologique de la pression artérielle (débit cardiaque, volume sanguin, résistance des vaisseaux).
\item Cite deux complications possibles d'une hypertension non traitée et propose une fiche de conseils de lutte adaptée au contexte local.
\end{enumerate}

\courssub{Ressources à mobiliser}

\begin{aretenir}[Repères indispensables]
\begin{itemize}
\item \cle{Glycémie} normale à jeun : entre \SI{0,70}{g/L} et \SI{1,10}{g/L} ; elle ne doit pas dépasser environ \SI{1,40}{g/L} après un repas et revient à la normale en 2 h grâce à l'\cle{insuline}.
\item Critères du \cle{diabète} : glycémie à jeun $\geq \SI{1,26}{g/L}$ (deux fois) ou glycémie $\geq \SI{2,00}{g/L}$ à 2 h de l'HGPO. Le glucose apparaît dans les urines quand la glycémie dépasse le \cle{seuil rénal} (environ \SI{1,80}{g/L}).
\item Régulation : hyperglycémie $\Rightarrow$ sécrétion d'\cle{insuline} (îlots de Langerhans, cellules $\beta$) $\Rightarrow$ stockage du glucose en glycogène dans le foie et les muscles (glycogénèse) ; hypoglycémie $\Rightarrow$ sécrétion de \cle{glucagon} (cellules $\alpha$) $\Rightarrow$ libération du glucose hépatique (glycogénolyse et néoglucogenèse).
\item \cle{Pression artérielle} normale au repos : environ $12/7$ à $13/8$ cm de Hg ; hypertension si $\geq 14/9$ cm de Hg de façon permanente ; hypotension si $\leq 10/6$ cm de Hg avec symptômes.
\item La pression artérielle (PA) dépend de trois paramètres physiologiques : le \cle{débit cardiaque} (DC), le \cle{volume sanguin} (VS) et l'\cle{état de résistance des artérioles} (élasticité, diamètre) ou résistances vasculaires périphériques (RVP) ; relation fondamentale : $\text{PA} \propto \text{DC} \times \text{RVP}$.
\end{itemize}
\end{aretenir}

\courssub{Démarche et corrigé commenté}

\begin{exoresolu}[Dossier 1 — Le cas de Mme A.]
Exploiter la courbe et les analyses, poser le diagnostic, expliquer les symptômes et proposer des conseils.
\tcblower
\textbf{Étape 1 — Exploitation de la courbe d'HGPO.} À jeun (temps 0), la glycémie de Mme A. vaut \SI{1,50}{g/L}, très au-dessus de la normale ($\leq \SI{1,10}{g/L}$). Après ingestion de la charge de glucose (\SI{75}{g}), elle culmine à \SI{2,45}{g/L} vers 60 min et reste anormalement élevée à \SI{2,30}{g/L} à 120 min : il n'y a \textbf{pas de retour à la normale} en 2 heures. Chez le témoin sain, la glycémie à jeun est normale (\SI{0,90}{g/L}), le pic est modéré (\SI{1,45}{g/L} à 30 min) et le retour à la normale (\SI{0,95}{g/L}) est complet à 120 min, témoignant d'une régulation hormonale efficace.

\textbf{Étape 2 — Diagnostic argumenté.} La glycémie à jeun (\SI{1,50}{g/L}) est $\geq \SI{1,26}{g/L}$ et la glycémie à 2 h (\SI{2,30}{g/L}) est $\geq \SI{2,00}{g/L}$ : les deux critères biologiques du \textbf{diabète sucré} sont réunis. La glycosurie (glucose dans les urines) confirme le diagnostic : la glycémie dépassant le \textbf{seuil rénal} de réabsorption du glucose (environ \SI{1,80}{g/L}), le tubule proximal ne peut plus réabsorber tout le glucose filtré, celui-ci passe dans les urines en entraînant de l'eau par osmose. Vu l'âge (48 ans), l'installation progressive et le contexte de surcharge pondérale fréquent chez la vendeuse sédentaire, il s'agit d'un \textbf{diabète de type 2} (insulino-résistance), forme la plus fréquente au Cameroun.

\textbf{Étape 3 — Explication physiologique des symptômes.}
\begin{itemize}
\item \textbf{Polyurie / Polydipsie :} Le glucose urinaire (glycosurie) exerce un effet osmotique qui retient l'eau dans la lumière tubulaire $\rightarrow$ diurèse osmotique importante (polyurie) $\rightarrow$ déshydratation cellulaire et stimulation des osmorécepteurs hypothalamiques $\rightarrow$ soif intense (polydipsie).
\item \textbf{Amaigrissement :} Malgré l'hyperglycémie, le glucose ne pénètre pas efficacement dans les cellules cibles (foie, muscle, tissu adipeux) du fait de la résistance à l'insuline. L'organisme se croit en jeûne : il active la lipolyse (dégradation des triglycérides) et la protéolyse (dégradation des protéines musculaires) pour fournir de l'énergie $\rightarrow$ perte de poids.
\end{itemize}

\textbf{Étape 4 — Schéma fonctionnel de la régulation et étape défaillante.}
\begin{popfigure}
\begin{center}
\begin{tikzpicture}[node distance=1.8cm, >=Stealth, auto,
  proc/.style={rectangle, draw=popdark, thick, fill=popblueL, rounded corners, minimum height=1.1cm, text width=2.6cm, align=center, font=\small},
  arrow/.style={->, thick, popdark},
  fail/.style={->, thick, popink, dashed},
  annot/.style={draw=popink, thick, fill=poppinkL, rounded corners, font=\small, align=center, text width=3cm}]
\node[proc, align=center] (hg){Hyperglycémie\\après repas};
\node[proc, right=of hg, align=center] (panc){Pancréas endocrine\\Îlots de Langerhans\\(Cellules $\beta$)};
\node[proc, right=of panc, align=center] (ins){Sécrétion\\d'Insuline};
\node[proc, right=of ins, align=center] (cible){Cellules cibles\\(Foie, Muscle, Adipose)\\\textit{Glycogénèse / Lipogenèse}};
\node[proc, right=of cible, align=center] (ng){Normoglycémie\\$[\approx 1,0 \text{ g/L}]$};
\draw[arrow] (hg) -- node[above] {Stimule} (panc);
\draw[arrow] (panc) -- node[above] {Sécrète} (ins);
\draw[arrow] (ins) -- node[above] {Se lie récepteur} (cible);
\draw[arrow] (cible) -- node[above] {Stockage glucose} (ng);
\draw[arrow] (ng) |- node[near start, left] {Rétroaction -} (panc);
% Annotation défaillance
\node[annot, below=0.8cm of cible, align=center] (defaut){Chez Mme A. (Diabète T2) :\\Résistance à l'insuline\\$\Rightarrow$ Translocation GLUT4 bloquée\\$\Rightarrow$ Glycogénèse inefficace};
\draw[fail] (defaut.north) -- (cible.south);
\end{tikzpicture}
\end{center}
\legende{Schéma fonctionnel de la régulation de la glycémie par l'insuline. La flèche pointillée rouge indique l'étape défaillante chez Mme A. : la résistance cellulaire à l'insuline empêche l'entrée du glucose et son stockage.}
\end{popfigure}

\textbf{Étape 5 — Fiche de conseils pour Mme A. (contexte Yaoundé).}
\begin{itemize}
\item \textbf{Alimentation :} Fractionner les repas (3 repas + 1 collation). Réduire les sucres rapides (boissons gazeuses, jus industriels, beignets, pain blanc). Privilégier les féculents à index glycémique bas (igname, macabo, plantain peu mûr, niébé, maïs) en portions contrôlées (poing fermé). Augmenter les légumes verts (feuilles d'éru, folong, gombo) et les fibres. Cuisiner avec peu d'huile (préférer l'huile rouge non chauffée ou l'arachide), limiter le cube d'assaisonnement.
\item \textbf{Activité physique :} Marcher activement 30 min minimum par jour (aller-retour marché, monter les escaliers). Profiter des temps d'attente au stand pour faire des flexions ou marcher sur place.
\item \textbf{Suivi médical :} Contrôle glycémique mensuel au centre de santé, bilan rénal et ophtalmologique annuel. Prise rigoureuse du traitement oral prescrit (metformine le plus souvent). Dépistage du diabète chez les enfants et frères/sœurs (facteur héréditaire).
\item \textbf{Hygiène de vie :} Arrêt du tabac si fumeuse, hydratation à l'eau pure (1,5--2 L/jour), sommeil suffisant.
\end{itemize}
\end{exoresolu}

\begin{exoresolu}[Dossier 2 — Le cas de M. B.]
Interpréter les mesures tensionnelles, identifier les facteurs favorisants et proposer des conseils.
\tcblower
\textbf{Étape 1 — Technique de mesure (Méthode auscultatoire de Korotkoff).} On place un brassard gonflable autour du bras nu, à hauteur du cœur, le patient assis, au calme depuis 5 min. On gonfle le brassard au-dessus de la pression systolique attendue (environ 20 cm Hg au-dessus de la disparition du pouls radial). On dégonfle lentement (\SI{2}{--}\SI{3}{mmHg/s} soit \SI{0,2}{--}\SI{0,3}{cm Hg/s}). Le \textbf{premier bruit} perçu au stéthoscope (bruits de Korotkoff) correspond à la \textbf{pression artérielle systolique (PAS)} : pression maximale lors de la systole ventriculaire (éjection du sang). La \textbf{disparition des bruits} (5e phase) correspond à la \textbf{pression artérielle diastolique (PAD)} : pression minimale lors de la diastole (remplissage ventriculaire). La mesure « 16/10 » signifie donc une PAS de \textbf{16 cm de Hg} et une PAD de \textbf{10 cm de Hg} (équivalent à 160/100 mm Hg).

\textbf{Étape 2 — Diagnostic d'hypertension artérielle (HTA).} Les trois mesures successives (\textbf{J1 : 16/10}, \textbf{J2 : 15/11}, \textbf{J3 : 17/10}) sont \textbf{toutes supérieures} au seuil diagnostique de $14/9$ cm de Hg. Une seule mesure élevée ne suffit pas (stress de la consultation, « effet blouse blanche », variations physiologiques). C'est la \textbf{permanence} de l'élévation sur trois jours différents, au repos, qui permet de conclure à une \textbf{hypertension artérielle essentielle (primaire) stade 2}. Les symptômes (céphalées matinales, vertiges, acouphènes) sont évocateurs d'une HTA mal contrôlée.

\textbf{Étape 3 — Facteurs favorisants et paramètres physiologiques.}
\begin{center}
\begin{tabularx}{\linewidth}{|l|X|X|}
\hline
\rowcolor{popblueL} Facteur de risque (Questionnaire) & Mécanisme physiologique & Paramètre de la PA augmenté \\
\hline Excès de sel (cube, poisson fumé) & Rétention sodée $\rightarrow$ rétention d'eau & \textbf{Volume sanguin (VS)} $\uparrow$ \\
\hline Alcool régulier & Prise de poids, activation système nerveux sympathique, rigidité artérielle & \textbf{Débit cardiaque (DC)} $\uparrow$ et \textbf{RVP} $\uparrow$ \\
\hline Sédentarité (métier chauffeur) & Prise de poids, déconditionnement cardiovasculaire, rigidité vasculaire & \textbf{Débit cardiaque (DC)} $\uparrow$ et \textbf{RVP} $\uparrow$ \\
\hline Hérédité (père décédé « crise ») & Prédisposition génétique : structure vasculaire, sensibilité au sel, système rénine-angiotensine & \textbf{Résistances vasculaires périphériques (RVP)} $\uparrow$ \\
\hline Âge (55 ans) & Perte d'élasticité des grosses artères (rigidité aortique) & \textbf{RVP} $\uparrow$ (surtout PAS) \\
\hline
\end{tabularx}
\end{center}
Rappel : $\text{PA} \propto \text{DC} \times \text{RVP}$ et $\text{DC} = \text{Fréquence cardiaque} \times \text{Volume d'éjection systolique}$ (dépendant du VS). L'association de ces facteurs explique l'élévation durable des deux chiffres tensionnels.

\textbf{Étape 4 — Complications et fiche de conseils.}
\textbf{Complications majeures si non traitée :} \textbf{Accident Vasculaire Cérébral (AVC)} ischémique ou hémorragique (première cause de mortalité liée à l'HTA au Cameroun), \textbf{Insuffisance cardiaque} (hypertrophie ventriculaire gauche puis dilatation), \textbf{Insuffisance rénale chronique} (néphroangiosclérose), \textbf{Rétinopathie hypertensive} (cécité).

\textbf{Fiche de conseils pour M. B. (contexte chauffeur Yaoundé) :}
\begin{itemize}
\item \textbf{Alimentation :} \textbf{Réduire drastiquement le sel} : cuisiner \textbf{sans cube} (utiliser ail, oignon, gingembre, persil, céleri, poivre), rincer le poisson fumé/salé avant cuisson, éviter les saucisses, charcuteries, conserves. 5 g de sel/jour max (1 cuillère à café rase). Augmenter potassium (banane plantain, avocat, légumes verts) pour favoriser l'excrétion sodée.
\item \textbf{Alcool :} Arrêt total ou réduction maximale (\textless 2 verres/jour, pas tous les jours). L'alcool vide les réserves de magnésium et active le système sympathique.
\item \textbf{Activité physique :} Marcher 30 min continu \textbf{tous les jours} : descendre du taxi 2 arrêts avant le terminus, faire le tour du parcage au lieu de rester assis, monter les escaliers. Objectif : perte de 5--10\% du poids si surcharge.
\item \textbf{Suivi médical :} Contrôle tensionnel \textbf{hebdomadaire} au centre de santé (automesure validée si possible). \textbf{Observance stricte du traitement antihypertenseur} (ne jamais arrêter sans avis médical même si « ça va mieux »). Bilan rénal (créatininémie, ionogramme), ECG, fond d'œil à la découverte puis annuel.
\item \textbf{Gestion du stress :} Respiration abdominale aux feux rouges, pauses régulières hors du véhicule.
\end{itemize}
\end{exoresolu}

\begin{attention}[Pièges fréquents à éviter au Bac]
\begin{itemize}
\item Ne conclus jamais au diabète sur une \textbf{seule} valeur de glycémie : le diagnostic nécessite \textbf{deux} dosages à jeun $\geq \SI{1,26}{g/L}$ \textbf{OU} une HGPO avec glycémie à 2 h $\geq \SI{2,00}{g/L}$ (confirmée par un second examen si pas de symptômes).
\item Ne confonds pas \textbf{diabète de type 1} (destruction auto-immune des cellules $\beta$, sujet jeune, maigre, insulinodépendant immédiat, acido-cétose fréquente) et \textbf{type 2} (insulino-résistance + défaut de sécrétion secondaire, adulte, surpoids fréquent, non insulinodépendant au début, contexte mode de vie).
\item Une tension élevée \textbf{un seul jour} ne suffit pas pour diagnostiquer une HTA : il faut \textbf{au moins 3 mesures sur 3 consultations distinctes} (ou automesure validée sur 3--7 jours). L'« effet blouse blanche » est fréquent.
\item Dans la notation « 16/10 », le premier nombre est la \textbf{systolique} (PAS), le second la \textbf{diastolique} (PAD). Ne les inverse jamais.
\item La glycosurie n'apparaît que si la glycémie dépasse le \textbf{seuil rénal} ($\approx \SI{1,80}{g/L}$) : une glycémie à \SI{1,50}{g/L} ne donne \textbf{pas} de glycosurie.
\item L'hypertension artérielle est \textbf{asymptomatique} dans la grande majorité des cas : l'absence de maux de tête ne signifie pas une PA normale. Seule la mesure le prouve.
\end{itemize}
\end{attention}

\courssub{Bilan de l'activité}

Cette activité t'a permis de mobiliser simultanément tous les savoir-faire de la leçon dans une situation authentique de consultation au centre de santé. Tu as vérifié qu'un \textbf{diagnostic médical} ne se devine pas : il se \textbf{démontre} rigoureusement en comparant des données chiffrées (courbe d'HGPO, bandelette urinaire, mesures tensionnelles répétées) à des valeurs de référence validées. Tu as relié chaque symptôme clinique à un mécanisme physiologique précis et étayé : dérèglement de la régulation hormonale de la glycémie (résistance à l'insuline $\rightarrow$ hyperglycémie $\rightarrow$ glycosurie osmotique $\rightarrow$ polyurie/polydipsie ; défaut d'utilisation du glucose $\rightarrow$ catabolisme protéino-lipidique $\rightarrow$ amaigrissement) pour Mme A. ; augmentation durable du volume sanguin (sel) et des résistances vasculaires périphériques (âge, alcool, sédentarité, hérédité) pour M. B. Enfin, tu as constaté que le diabète de type 2 et l'hypertension artérielle, deux maladies métaboliques et cardiovasculaires en recrudescence alarmante au Cameroun (transition nutritionnelle, urbanisation, vieillissement), se préviennent et se contrôlent largement par des mesures simples, accessibles et peu coûteuses : alimentation équilibrée et peu salée, activité physique quotidienne, dépistage systématique et suivi régulier au centre de santé de proximité. La confrontation orale des conclusions des groupes a enfin montré qu'en médecine comme en sciences, c'est la qualité de l'argumentation structurée à partir des documents fournis qui fait la valeur scientifique d'une conclusion.

\newpage

\coursec{Méthodes et exercices résolus}

\courssub{Méthode 1 : Analyser et interpréter les variations de la glycémie avant et après un repas}

\begin{methode}[Lire une courbe de glycémie postprandiale]
\begin{enumerate}
\item \textbf{Identifier les grandeurs} : en abscisse, le temps (en h ou en min, l'origine étant souvent le début du repas) ; en ordonnée, la glycémie (en g/L ou en mmol/L, avec \SI{1}{g/L} = \SI{5,5}{mmol/L}).
\item \textbf{Repérer la valeur de référence} : la glycémie à jeun normale est comprise entre \num{0,70} et \SI{1,10}{g/L}. Tracer mentalement cette bande sur le graphique.
\item \textbf{Décrire la courbe en trois temps} : valeur initiale (à jeun), montée après le repas (pic, valeur et heure du pic), retour à la valeur de base (durée du retour).
\item \textbf{Comparer les courbes} si plusieurs sujets : un sujet diabétique présente une glycémie à jeun élevée ($\geq \SI{1,26}{g/L}$ à deux reprises), un pic plus haut et un retour lent ou absent.
\item \textbf{Interpréter} en mobilisant les savoirs : la montée correspond à l'absorption intestinale du glucose ; le retour à la normale traduit l'action de l'\cle{insuline} (stockage du glucose en glycogène hépatique et musculaire, entrée du glucose dans les cellules).
\item \textbf{Conclure} par une phrase qui relie l'observation au mécanisme de régulation.
\end{enumerate}
\end{methode}

\begin{exoresolu}[Glycémies de deux sujets après ingestion de glucose]
Lors d'une épreuve d'\cle{hyperglycémie provoquée} (HGPO) réalisée à l'hôpital, deux sujets A et B, à jeun depuis 12 h, ingèrent à $t = 0$ une solution contenant \SI{75}{g} de glucose. Leur glycémie est mesurée toutes les 30 min. Les résultats sont les suivants :

\begin{center}
\begin{tabularx}{\linewidth}{|l|X|X|X|X|X|X|}\hline
\rowcolor{popblueL} Temps (min) & 0 & 30 & 60 & 90 & 120 & 180 \\ \hline
Glycémie de A (g/L) & 0,85 & 1,40 & 1,20 & 1,00 & 0,90 & 0,85 \\ \hline
Glycémie de B (g/L) & 1,45 & 2,20 & 2,40 & 2,25 & 2,00 & 1,80 \\ \hline
\end{tabularx}
\end{center}

\begin{enumerate}
\item Tracer les deux courbes sur un même graphique.
\item Décrire et comparer les variations de la glycémie des deux sujets.
\item Indiquer, en le justifiant, lequel des deux sujets est diabétique.
\end{enumerate}

\tcblower

\textbf{1. Tracé des courbes.} On place le temps en abscisse et la glycémie en ordonnée, avec une échelle adaptée.

\begin{popfigure}
\begin{center}
\begin{tikzpicture}
\begin{axis}[
width=0.9\linewidth, height=6cm,
xlabel={Temps (min)}, ylabel={Glycémie (g/L)},
xmin=0, xmax=190, ymin=0.5, ymax=2.6,
xtick={0,30,60,90,120,180},
grid=both, grid style={popline!40},
legend style={at={(0.97,0.97)}, anchor=north east, font=\small},
]
\addplot[popcurve, color=popblue, mark=*, mark size=1.5pt] coordinates {(0,0.85) (30,1.40) (60,1.20) (90,1.00) (120,0.90) (180,0.85)};
\addplot[popcurve, color=poporange, mark=square*, mark size=1.5pt] coordinates {(0,1.45) (30,2.20) (60,2.40) (90,2.25) (120,2.00) (180,1.80)};
\addplot[domain=0:190, samples=2, dashed, color=popgreen] {1.26};
\legend{Sujet A, Sujet B, Seuil diabète (\SI{1,26}{g/L})}
\end{axis}
\end{tikzpicture}
\end{center}
\legende{Épreuve d'hyperglycémie provoquée chez deux sujets. Le sujet A (bleu) revient à sa glycémie de base en 2 h ; le sujet B (orange) reste durablement au-dessus du seuil de diabète.}
\end{popfigure}

\textbf{2. Description et comparaison.}

\emph{Sujet A} : la glycémie à jeun est normale (\SI{0,85}{g/L}). Elle s'élève après l'ingestion pour atteindre un pic de \SI{1,40}{g/L} à $t = 30$ min, puis décroît et revient à la valeur initiale (\SI{0,85}{g/L}) en 180 min. La régulation est efficace.

\emph{Sujet B} : la glycémie à jeun est déjà anormalement élevée (\SI{1,45}{g/L} $>$ \SI{1,26}{g/L}). Le pic est beaucoup plus haut (\SI{2,40}{g/L} à 60 min) et la décroissance est lente : à 180 min, la glycémie vaut encore \SI{1,80}{g/L}, loin de la normale.

\emph{Comparaison chiffrée} : l'augmentation maximale est de $1,40 - 0,85 = \SI{0,55}{g/L}$ chez A contre $2,40 - 1,45 = \SI{0,95}{g/L}$ chez B ; le retour à la normale s'effectue en moins de 2 h chez A et n'est pas achevé à 3 h chez B.

\textbf{3. Identification du sujet diabétique.} Le sujet B est diabétique, car sa glycémie à jeun (\SI{1,45}{g/L}) dépasse le seuil diagnostique de \SI{1,26}{g/L} et sa glycémie reste supérieure ou égale à \SI{2,00}{g/L} deux heures après la charge en glucose (critère de l'épreuve d'HGPO). Cela traduit une insuffisance d'insuline (diabète de type 1) ou une résistance des cellules à l'insuline (diabète de type 2) : le glucose absorbé n'est pas capté ni stocké normalement, d'où une hyperglycémie prolongée.

\textbf{Conclusion} : chez le sujet sain, l'insuline ramène la glycémie à sa valeur de référence en moins de 2 h ; l'absence de ce retour à la normale est la signature biologique du diabète.
\end{exoresolu}

\courssub{Méthode 2 : Réaliser un schéma fonctionnel de la régulation hormonale de la glycémie}

\begin{methode}[Construire une boucle de régulation hormonale]
\begin{enumerate}
\item \textbf{Identifier les éléments de la boucle} : la variable réglée (la glycémie), le capteur (cellules des îlots de Langerhans du pancréas), les effecteurs (foie, muscles, tissu adipeux) et les messagers (hormones).
\item \textbf{Distinguer les deux hormones antagonistes} : l'\cle{insuline} (cellules $\beta$ des îlots), hypoglycémiante, sécrétée quand la glycémie monte ; le \cle{glucagon} (cellules $\alpha$), hyperglycémiant, sécrété quand la glycémie baisse.
\item \textbf{Représenter chaque branche} par des flèches étiquetées : hyperglycémie $\rightarrow$ sécrétion d'insuline $\rightarrow$ glycogénèse hépatique + entrée du glucose dans les cellules $\rightarrow$ baisse de la glycémie.
\item \textbf{Indiquer le signe de la rétroaction} : le retour de la glycémie vers la normale inhibe la sécrétion hormonale (rétroaction négative, signe $-$).
\item \textbf{Vérifier la cohérence} : chaque flèche doit porter un verbe d'action (stimule, inhibe, provoque) et le schéma doit « boucler ».
\end{enumerate}
\end{methode}

\begin{exoresolu}[Schéma fonctionnel de la régulation de la glycémie]
Un élève de Terminale C doit présenter, sous forme de schéma fonctionnel, la régulation de la glycémie après un repas riche en glucides puis après un effort prolongé. Construire ce schéma et commenter chaque branche.

\tcblower

\textbf{Schéma fonctionnel :}

\begin{popfigure}
\begin{center}
\begin{tikzpicture}[
node distance=1.1cm and 2.2cm,
boite/.style={draw=popblue, thick, rounded corners, fill=popblueL, text width=3.4cm, align=center, minimum height=1cm, font=\small},
horm/.style={draw=poporange, thick, rounded corners, fill=poporangeL, text width=2.6cm, align=center, minimum height=0.9cm, font=\small},
fleche/.style={-{Stealth[length=3mm]}, thick, color=popdark},
]
\node[boite] (hyper) {Hyperglycémie (après un repas) $>$ \SI{1,10}{g/L}};
\node[horm, right=of hyper] (insu) {Insuline (cellules $\beta$ des îlots de Langerhans)};
\node[boite, below=of hyper] (hypo) {Hypoglycémie (jeûne, effort) $<$ \SI{0,70}{g/L}};
\node[horm, right=of hypo] (gluca) {Glucagon (cellules $\alpha$ des îlots de Langerhans)};
\node[boite, right=of insu, yshift=-1.65cm] (foie) {Foie, muscles, tissu adipeux};
\node[boite, draw=popgreen, fill=popgreenL, left=of foie, xshift=-0.2cm] (norm) {Glycémie normale \num{0,70}--\SI{1,10}{g/L}};

\draw[fleche] (hyper) -- node[above, font=\scriptsize]{stimule} (insu);
\draw[fleche] (hypo) -- node[above, font=\scriptsize]{stimule} (gluca);
\draw[fleche] (insu) -- node[above, sloped, font=\scriptsize, align=center]{glycogénèse, entrée\\ du glucose cellulaire} (foie);
\draw[fleche] (gluca) -- node[below, sloped, font=\scriptsize, align=center]{glycogénolyse,\\ néoglucogenèse} (foie);
\draw[fleche, color=popgreen] (foie) -- node[above, font=\scriptsize]{retour à la normale} (norm);
\draw[fleche, dashed, color=poppurple] (norm.north) |- node[pos=0.25, left, font=\scriptsize]{rétroaction $-$} (hyper.west);
\draw[fleche, dashed, color=poppurple] (norm.south) |- node[pos=0.25, left, font=\scriptsize]{rétroaction $-$} (hypo.west);
\end{tikzpicture}
\end{center}
\legende{Boucle de régulation hormonale de la glycémie. 1 : détection par le pancréas ; 2 : sécrétion d'insuline ou de glucagon ; 3 : action sur le foie et les tissus ; 4 : rétroaction négative qui arrête la sécrétion.}
\end{popfigure}

\textbf{Commentaire de la branche hyperglycémie} : après un repas glucidique, l'absorption intestinale fait passer la glycémie au-dessus de \SI{1,10}{g/L}. Les cellules $\beta$ des îlots de Langerhans libèrent de l'insuline, qui provoque la transformation du glucose en glycogène dans le foie (glycogénèse) et l'entrée du glucose dans les cellules musculaires et adipeuses : la glycémie redescend.

\textbf{Commentaire de la branche hypoglycémie} : au cours d'un jeûne ou d'un effort, la glycémie descend sous \SI{0,70}{g/L}. Les cellules $\alpha$ libèrent du glucagon, qui déclenche la dégradation du glycogène hépatique (glycogénolyse) et la fabrication de glucose à partir d'autres métabolites (néoglucogenèse) : la glycémie remonte.

\textbf{Conclusion} : la glycémie est maintenue constante par un double contrôle hormonal antagoniste à rétroaction négative, le pancréas jouant le rôle de capteur et de glande, le foie celui d'effecteur principal.
\end{exoresolu}

\courssub{Méthode 3 : Interpréter les résultats d'analyses médicales dans le cas du diabète}

\begin{methode}[Exploiter un bilan biologique de diabète]
\begin{enumerate}
\item \textbf{Comparer chaque paramètre aux valeurs de référence} : glycémie à jeun (normale : \num{0,70}--\SI{1,10}{g/L} ; diabète si $\geq \SI{1,26}{g/L}$ à deux reprises), hémoglobine glyquée HbA$_{1c}$ (normale $< 6\,\%$ ; diabète si $\geq 6{,}5\,\%$), présence de glucose (glycosurie) ou de corps cétoniques (cétonurie) dans les urines.
\item \textbf{Expliquer la glycosurie} : quand la glycémie dépasse le \cle{seuil rénal} de réabsorption du glucose ($\approx \SI{1,80}{g/L}$ en moyenne, mais variable d'un individu à l'autre et abaissé chez certains patients), les reins ne réabsorbent plus tout le glucose filtré, qui passe alors dans les urines.
\item \textbf{Typer le diabète} : diabète de \cle{type 1} (jeune, maigre, destruction auto-immune des cellules $\beta$, absence d'insuline, traitement par injections d'insuline) ; diabète de \cle{type 2} (adulte, souvent obèse, insulinorésistance, traitement par régime, exercice physique et antidiabétiques oraux).
\item \textbf{Relier les symptômes aux analyses} : polyurie, polydipsie, amaigrissement s'expliquent par l'hyperglycémie et la glycosurie.
\item \textbf{Conclure} par le diagnostic et proposer les moyens de lutte adaptés.
\end{enumerate}
\end{methode}

\begin{exoresolu}[Deux bilans biologiques à l'hôpital de district]
Deux patients consultent dans un hôpital de district du Cameroun pour fatigue, soif intense et urines abondantes. Les résultats de leurs analyses sont consignés ci-dessous.

\begin{center}
\begin{tabularx}{\linewidth}{|l|X|X|X|}\hline
\rowcolor{popblueL} Paramètre & Patient 1 (14 ans, 40 kg) & Patient 2 (52 ans, 95 kg) & Valeurs de référence \\ \hline
Glycémie à jeun (g/L) & 2,85 & 1,65 & 0,70 -- 1,10 \\ \hline
HbA$_{1c}$ (\%) & 11,2 & 8,1 & $< 6$ \\ \hline
Glucose dans les urines & présent & présent & absent \\ \hline
Corps cétoniques urinaires & présents & absents & absents \\ \hline
Insulinémie & très faible & normale à élevée & normale \\ \hline
\end{tabularx}
\end{center}

\begin{enumerate}
\item Montrer que les deux patients sont diabétiques.
\item Déterminer, en le justifiant, le type de diabète de chaque patient.
\item Proposer pour chacun un moyen de lutte adapté.
\end{enumerate}

\tcblower

\textbf{1. Les deux patients sont diabétiques.} Le patient 1 présente une glycémie à jeun de \SI{2,85}{g/L} et le patient 2 de \SI{1,65}{g/L} : ces deux valeurs dépassent le seuil diagnostique de \SI{1,26}{g/L}. De plus, leurs HbA$_{1c}$ (11,2\,\% et 8,1\,\%) sont supérieures au seuil de 6,5\,\%, ce qui confirme une hyperglycémie chronique (l'HbA$_{1c}$ reflète la glycémie moyenne des 2 à 3 derniers mois). Enfin, la glycosurie observée chez les deux s'explique : chez le patient 1, la glycémie (\SI{2,85}{g/L}) dépasse largement et en permanence le seuil rénal de réabsorption ($\approx \SI{1,80}{g/L}$) ; chez le patient 2, la glycémie à jeun de \SI{1,65}{g/L} s'accompagne d'hyperglycémies postprandiales plus élevées qui franchissent ce seuil après les repas, d'autant que le seuil rénal peut être abaissé chez certains individus. Le glucose non réabsorbé est éliminé dans les urines, entraînant par entraînement osmotique une polyurie, donc une déshydratation responsable de la polydipsie.

\textbf{2. Typage des diabètes.}

\emph{Patient 1 : diabète de type 1 (insulino-dépendant).} Justification : sujet jeune (14 ans) et maigre, insulinémie très faible (les cellules $\beta$ des îlots de Langerhans sont détruites par un processus auto-immun), présence de corps cétoniques dans les urines (en l'absence d'insuline, les cellules ne peuvent utiliser le glucose et l'organisme dégrade massivement les lipides, d'où une cétose).

\emph{Patient 2 : diabète de type 2 (non insulino-dépendant).} Justification : sujet adulte (52 ans) et obèse (95 kg), insulinémie normale voire élevée : le pancréas produit de l'insuline mais les cellules cibles y sont résistantes (insulinorésistance liée à l'obésité et à la sédentarité) ; absence de cétose car l'insuline résiduelle empêche la dégradation massive des graisses.

\textbf{3. Moyens de lutte.}

\emph{Patient 1} : injections quotidiennes d'\cle{insuline} (l'hormone manquante ne peut être remplacée que par voie injectable, car c'est une protéine digérée si elle est avalée), surveillance régulière de la glycémie, éducation diététique.

\emph{Patient 2} : régime alimentaire pauvre en sucres rapides et en graisses, pratique régulière d'exercice physique, réduction du poids, antidiabétiques oraux (qui améliorent la sensibilité des cellules à l'insuline) ; l'insuline n'est nécessaire qu'en cas d'échec.

\textbf{Conclusion} : la confrontation des analyses (glycémie, HbA$_{1c}$, insulinémie, cétonurie) au profil du patient permet de poser le diagnostic de diabète et d'en préciser le type, condition d'un traitement adapté.
\end{exoresolu}

\courssub{Méthode 4 : Expliquer la technique de mesure de la pression artérielle}

\begin{methode}[Décrire la mesure de la pression artérielle au sphygmomanomètre]
\begin{enumerate}
\item \textbf{Nommer le matériel} : sphygmomanomètre (tensiomètre) à brassard gonflable, associé à un stéthoscope (méthode auscultatoire).
\item \textbf{Décrire le protocole en étapes} : brassard autour du bras à hauteur du c\oe ur, sujet au repos ; gonflage du brassard au-dessus de la pression systolique (l'artère humérale est occluse, aucun bruit) ; dégonflage progressif.
\item \textbf{Identifier les deux valeurs} : la \cle{pression systolique} (maxima) correspond à l'apparition des premiers bruits (le sang repasse par saccades lors de la contraction ventriculaire) ; la \cle{pression diastolique} (minima) correspond à la disparition des bruits (le sang circule à nouveau librement pendant le relâchement cardiaque).
\item \textbf{Exprimer le résultat} en cmHg ou mmHg : valeur normale au repos $\approx \SI{12}{cmHg}$ (systolique) / \SI{8}{cmHg} (diastolique), notée « 12/8 » ; on parle d'hypertension au-delà de 14/9.
\item \textbf{Préciser les conditions de validité} : sujet assis, calme depuis 5 min, plusieurs mesures, car stress, effort et émotions font varier la pression.
\end{enumerate}
\end{methode}

\begin{exoresolu}[Mesure de la pression artérielle en consultation]
Au centre de santé, l'infirmier place un brassard autour du bras d'un patient et le gonfle jusqu'à \SI{18}{cmHg} : aucun pouls n'est perceptible au poignet et aucun bruit n'est entendu au stéthoscope placé au pli du coude. Il dégonfle lentement le brassard. À \SI{13,5}{cmHg}, des bruits frappés apparaissent ; à \SI{9}{cmHg}, ils disparaissent.

\begin{enumerate}
\item Expliquer pourquoi aucun bruit n'est entendu lorsque le brassard est gonflé à \SI{18}{cmHg}.
\item Donner la pression artérielle du patient et justifier chaque valeur.
\item Dire si ce patient est hypertendu.
\end{enumerate}

\tcblower

\textbf{1. Absence de bruit à \SI{18}{cmHg}.} La pression exercée par le brassard (\SI{18}{cmHg}) est supérieure à la pression maximale du sang dans l'artère humérale. L'artère est donc complètement écrasée (occluse) : le sang ne passe plus, même pendant la systole. Sans écoulement sanguin, il n'y a ni pouls au poignet ni bruit audible au stéthoscope.

\textbf{2. Pression artérielle du patient.}

\emph{Pression systolique (maxima) = \SI{13,5}{cmHg}} : lors du dégonflage, dès que la pression du brassard devient inférieure à la pression maximale du sang, celui-ci franchit l'obstacle par saccades au moment de la contraction des ventricules ; ce flux intermittent et turbulent produit les premiers bruits entendus.

\emph{Pression diastolique (minima) = \SI{9}{cmHg}} : lorsque la pression du brassard devient inférieure à la pression minimale du sang, l'artère reste ouverte en permanence, même pendant le relâchement cardiaque ; l'écoulement redevient régulier et silencieux, les bruits disparaissent.

La pression artérielle du patient est donc de \textbf{13,5/9 cmHg} (soit 135/90 mmHg).

\textbf{3. Interprétation.} La valeur normale au repos est d'environ 12/8 et le seuil d'hypertension est fixé à 14/9. La systolique (13,5) est inférieure à 14, mais la diastolique atteint 9 : le patient présente une pression à la limite haute (dite « normale haute »). Une seule mesure ne suffit pas à conclure : il faut confirmer par plusieurs mesures, au repos, à des jours différents, avant de parler d'hypertension artérielle.

\textbf{Conclusion} : la méthode auscultatoire repose sur l'apparition (systolique) puis la disparition (diastolique) des bruits liés au passage du sang dans l'artère comprimée ; le patient mesuré à 13,5/9 doit faire l'objet d'une surveillance.
\end{exoresolu}

\courssub{Méthode 5 : Analyser les conditions de variation de la pression artérielle}

\begin{methode}[Interpréter un document sur les variations de la pression artérielle]
\begin{enumerate}
\item \textbf{Identifier le facteur étudié} (effort physique, émotion, posture, âge, température, ingestion de sel) et la grandeur mesurée (systolique, diastolique, fréquence cardiaque).
\item \textbf{Décrire quantitativement} la variation : citer les valeurs avant/après, calculer l'écart ($\Delta P = P_{\text{après}} - P_{\text{avant}}$).
\item \textbf{Expliquer le mécanisme} en mobilisant les paramètres physiologiques : la pression artérielle dépend du \cle{débit cardiaque} (volume d'éjection systolique $\times$ fréquence cardiaque), du \cle{volume sanguin} et de l'\cle{élasticité des artères} (résistances périphériques).
\item \textbf{Distinguer variation physiologique et pathologique} : une élévation transitoire à l'effort ou à l'émotion est normale et réversible ; une élévation permanente au repos ($\geq 14/9$) définit l'hypertension artérielle.
\item \textbf{Conclure} en reliant le facteur au mécanisme et, si demandé, en proposant une mesure de prévention.
\end{enumerate}
\end{methode}

\begin{exoresolu}[Pression artérielle et effort physique]
On mesure la pression artérielle systolique d'un élève de Terminale dans différentes conditions :

\begin{center}
\begin{tabularx}{\linewidth}{|l|X|X|X|}\hline
\rowcolor{popblueL} Condition & Repos assis & Après 5 min de course & 10 min après l'arrêt \\ \hline
Pression systolique (cmHg) & 12,0 & 16,5 & 12,3 \\ \hline
Fréquence cardiaque (battements/min) & 70 & 145 & 78 \\ \hline
\end{tabularx}
\end{center}

\begin{enumerate}
\item Calculer la variation de la pression systolique entre le repos et la fin de l'effort.
\item Expliquer cette variation à l'aide des paramètres physiologiques de la pression artérielle.
\item Montrer que cette élévation n'est pas une hypertension artérielle.
\end{enumerate}

\tcblower

\textbf{1. Calcul de la variation.}
\[
\Delta P = P_{\text{effort}} - P_{\text{repos}} = 16{,}5 - 12{,}0 = \SI{4,5}{cmHg}.
\]
La pression systolique augmente de \SI{4,5}{cmHg} pendant l'effort, soit une élévation relative de $\dfrac{4{,}5}{12{,}0} \times 100 = 37{,}5\,\%$.

\textbf{2. Explication du mécanisme.} Pendant la course, les muscles ont des besoins accrus en dioxygène et en glucose. Le c\oe ur répond en augmentant le \cle{débit cardiaque} : la fréquence cardiaque passe de 70 à 145 battements/min et le volume d'éjection systolique augmente également. Le sang étant chassé plus vite et en plus grande quantité dans les artères, la pression exercée sur la paroi artérielle s'élève. La vasodilatation des artérioles musculaires limite toutefois cette élévation (la diastolique varie peu), ce qui dirige le sang vers les muscles en activité.

\textbf{3. Ce n'est pas une hypertension.} L'hypertension artérielle est définie comme une élévation \emph{permanente} de la pression au \emph{repos} ($\geq 14/9$ cmHg, confirmée par plusieurs mesures). Or ici : la valeur de repos est normale (\SI{12,0}{cmHg}), l'élévation est \emph{transitoire} et liée à l'effort, et la pression revient à la normale 10 min après l'arrêt (\SI{12,3}{cmHg}). Il s'agit donc d'une adaptation physiologique normale et réversible, non d'une maladie.

\textbf{Conclusion} : la pression artérielle varie avec le débit cardiaque ; l'élévation observée à l'effort est une réponse physiologique adaptée, à ne pas confondre avec l'hypertension artérielle chronique.
\end{exoresolu}

\courssub{Méthode 6 : Exploiter les résultats d'analyses médicales dans le cas de l'hyper/hypotension}

\begin{methode}[Conduire le raisonnement diagnostique sur un trouble de la pression artérielle]
\begin{enumerate}
\item \textbf{Comparer les valeurs mesurées aux seuils} : hypertension si pression au repos $\geq 14/9$ cmHg à plusieurs reprises ; hypotension si systolique $< \SI{9}{cmHg}$ environ, avec symptômes (vertiges, malaise).
\item \textbf{Rechercher les causes} dans le dossier : hypertension (excès de sel, obésité, sédentarité, stress chronique, alcool, hérédité, vieillissement des artères qui perdent leur élasticité) ; hypotension (déshydratation, hémorragie, jeûne prolongé, certains médicaments, passage brutal à la position debout).
\item \textbf{Relier les symptômes au mécanisme} : l'hypertension est souvent silencieuse (« tueur silencieux ») mais fatigue le c\oe ur et fragilise les vaisseaux (risques : accident vasculaire cérébral, infarctus, insuffisance rénale) ; l'hypotension provoque une mauvaise irrigation du cerveau (vertiges, évanouissement).
\item \textbf{Proposer les moyens de lutte} : hygiène de vie (réduction du sel, exercice, lutte contre l'obésité et le stress), traitement médicamenteux (antihypertenseurs), suivi régulier ; pour l'hypotension : hydratation, apports adaptés, lever progressif.
\item \textbf{Rédiger une conclusion} complète : diagnostic + cause probable + conduite à tenir.
\end{enumerate}
\end{methode}

\begin{exoresolu}[Deux cas cliniques au centre de santé]
\emph{Cas 1} : M. N., 55 ans, commerçant sédentaire, consomme des plats très salés et fume. Trois mesures effectuées à une semaine d'intervalle donnent : 16/10 ; 15,5/9,5 ; 16/10 cmHg. Il se dit pourtant « en pleine forme ».

\emph{Cas 2} : Mlle A., 20 ans, jeûne depuis la veille pour un examen et se lève brusquement du banc : elle ressent des vertiges et s'évanouit. Sa pression mesurée allongée est de 10/6,5 cmHg ; debout, elle chute à 8/5 cmHg.

\begin{enumerate}
\item Établir le diagnostic pour chaque cas en justifiant par les valeurs.
\item Expliquer l'absence de symptômes chez M. N. et les risques qu'il encourt.
\item Expliquer le malaise de Mlle A. et proposer une conduite à tenir pour chaque cas.
\end{enumerate}

\tcblower

\textbf{1. Diagnostics.}

\emph{Cas 1 : hypertension artérielle.} Les trois mesures au repos sont toutes supérieures au seuil de 14/9 cmHg (16/10 ; 15,5/9,5 ; 16/10). La répétition des valeurs élevées à des jours différents élimine une élévation transitoire due au stress de la consultation : le diagnostic d'hypertension artérielle permanente est confirmé.

\emph{Cas 2 : hypotension (orthostatique).} La pression allongée (10/6,5) est déjà basse, et elle chute à 8/5 cmHg au passage en position debout, avec une systolique inférieure à \SI{9}{cmHg} accompagnée de symptômes (vertiges, évanouissement) : c'est une hypotension orthostatique, favorisée par le jeûne prolongé.

\textbf{2. Hypertension silencieuse et risques.} L'hypertension artérielle ne provoque souvent aucun symptôme apparent : on l'appelle le « tueur silencieux ». Pourtant, la pression élevée exerce en permanence une contrainte excessive sur la paroi des artères, qui s'épaississent, perdent leur élasticité et peuvent se rompre ou se boucher. Le c\oe ur, qui doit vaincre une résistance accrue, se fatigue et s'hypertrophie. M. N. court donc des risques graves : accident vasculaire cérébral (AVC), infarctus du myocarde, insuffisance rénale, troubles de la vision. Ses facteurs de risque (sel, tabac, sédentarité) aggravent le pronostic.

\textbf{3. Malaise de Mlle A. et conduites à tenir.}

\emph{Mécanisme du malaise} : au passage brutal en position debout, le sang s'accumule par gravité dans les membres inférieurs ; le retour veineux vers le c\oe ur diminue, donc le débit cardiaque et la pression artérielle chutent. Le cerveau, momentanément mal irrigué et mal oxygéné, provoque vertiges puis perte de connaissance. Le jeûne (hypoglycémie associée et volume sanguin réduit) aggrave le phénomène.

\emph{Conduite à tenir pour M. N.} : réduire fortement la consommation de sel, arrêter le tabac, pratiquer une activité physique régulière (marche quotidienne), perdre du poids si nécessaire, et suivre un traitement antihypertenseur prescrit par le médecin avec contrôles réguliers de la pression.

\emph{Conduite à tenir pour Mlle A.} : l'allonger jambes surélevées lors du malaise, la réhydrater et lui faire reprendre une alimentation normale ; à l'avenir, éviter les jeûnes prolongés et se lever progressivement.

\textbf{Conclusion} : l'interprétation des mesures de pression artérielle exige la comparaison aux seuils (14/9 pour l'hypertension), la répétition des mesures et la prise en compte du contexte ; la prévention de l'hypertension repose d'abord sur l'hygiène de vie.
\end{exoresolu}

\begin{attention}[Les 5 erreurs qui coûtent des points au Bac]
\begin{enumerate}
\item \textbf{Confondre insuline et glucagon.} L'insuline est \emph{hypoglycémiante} (sécrétée quand la glycémie monte, elle fait stocker le glucose) ; le glucagon est \emph{hyperglycémiant} (sécrété quand la glycémie baisse, il libère le glucose du foie). Inverser les deux dans un schéma de régulation annule la logique de toute la copie.
\item \textbf{Oublier les unités et les seuils chiffrés.} Une glycémie s'exprime en g/L (ou mmol/L) : écrire « glycémie = 1,26 » sans unité est une faute. Les seuils à citer : glycémie normale à jeun \num{0,70}--\SI{1,10}{g/L}, diabète si $\geq \SI{1,26}{g/L}$, pression normale 12/8, hypertension si $\geq 14/9$ cmHg.
\item \textbf{Décrire une courbe sans l'interpréter.} Après avoir lu les valeurs (pic, durée du retour), il faut obligatoirement relier la courbe au mécanisme : montée = absorption intestinale ; descente = action de l'insuline. Une description « photographique » sans explication biologique ne rapporte que la moitié des points.
\item \textbf{Diagnostiquer une hypertension sur une seule mesure.} La pression varie avec l'effort, le stress et la posture : le diagnostic d'hypertension exige des valeurs $\geq 14/9$ \emph{au repos}, confirmées par \emph{plusieurs mesures}. De même, une élévation à l'effort est physiologique, pas pathologique.
\item \textbf{Confondre diabète de type 1 et de type 2 dans le traitement.} Le diabète de type 1 (jeune, insulinémie effondrée) se traite par \emph{injections} d'insuline — jamais par voie orale, car l'insuline est une protéine digérée. Le diabète de type 2 (adulte, souvent obèse, insulinorésistance) se combat d'abord par le régime, l'exercice physique et les antidiabétiques oraux. Prescrire de l'insuline orale ou un simple régime à un diabétique de type 1 est une erreur grave.
\end{enumerate}
\end{attention}

\newpage

\coursec{Exercices}

\courssub{Je vérifie mes connaissances}

\begin{exercice}[QCM : la glycémie]
Pour chaque question, choisis la (ou les) bonne(s) réponse(s).
\begin{enumerate}
\item La glycémie normale à jeun est comprise entre :
\begin{enumerate}
\item 0,10 et 0,30 g/L ; \item 0,70 et 1,10 g/L ; \item 1,50 et 2,00 g/L ; \item 3,50 et 5,00 g/L.
\end{enumerate}
\item L'hormone hypoglycémiante est :
\begin{enumerate}
\item le glucagon ; \item l'adrénaline ; \item l'insuline ; \item le cortisol.
\end{enumerate}
\item L'insuline est sécrétée par :
\begin{enumerate}
\item les cellules $\alpha$ des îlots de Langerhans ; \item les cellules $\beta$ des îlots de Langerhans ; \item les hépatocytes ; \item les reins.
\end{enumerate}
\item Une glycémie à jeun supérieure à 1,26 g/L, vérifiée deux fois, signe :
\begin{enumerate}
\item une hypoglycémie ; \item un diabète ; \item une hypertension ; \item une glycémie normale.
\end{enumerate}
\end{enumerate}
\end{exercice}

\begin{exercice}[Vrai ou faux — à justifier]
Réponds par vrai ou faux, puis justifie ta réponse en une ou deux phrases.
\begin{enumerate}
\item Le foie peut stocker le glucose sous forme de glycogène après un repas.
\item Le glucagon provoque une baisse de la glycémie.
\item Le diabète de type 1 est lié à une destruction auto-immune des cellules $\beta$ du pancréas.
\item La pression artérielle systolique correspond à la phase de relaxation du cœur.
\item Une pression artérielle de 12/8 (en cm de Hg) est une valeur normale au repos.
\end{enumerate}
\end{exercice}

\begin{exercice}[Définitions]
Définis précisément, en une phrase chacune, les termes suivants : \cle{glycémie} ; \cle{diabète} ; \cle{pression artérielle} ; \cle{hypertension artérielle} ; \cle{hypoglycémie}.
\end{exercice}

\begin{exercice}[Calculs directs]
\begin{enumerate}
\item Un patient présente une glycémie à jeun de 1,40 g/L. Sachant que la masse molaire du glucose est $M = 180$ g/mol, exprime cette glycémie en mmol/L (on donnera le résultat avec deux chiffres significatifs).
\item Le sang d'un adulte représente environ 5 L. Calcule la masse totale de glucose circulant chez un sujet dont la glycémie est de 0,95 g/L.
\item Chez un patient, on mesure une pression artérielle maximale de 16 cm de Hg et une pression minimale de 9 cm de Hg. Calcule la pression différentielle (pression pulsée) et commente le résultat.
\end{enumerate}
\end{exercice}

\courssub{Je m'entraîne}

\begin{exercice}[Le rôle du foie dans la régulation de la glycémie]
On mesure la glycémie (en g/L) dans la veine porte hépatique (qui amène le sang digestif au foie) et dans la veine sus-hépatique (qui ramène le sang du foie vers le cœur) chez un chien, à jeun puis 2 h après un repas riche en glucides.
\begin{center}
\begin{tabularx}{\linewidth}{|l|X|X|}
\hline
\rowcolor{popblueL} \textbf{Moment du prélèvement} & \textbf{Veine porte (g/L)} & \textbf{Veine sus-hépatique (g/L)} \\
\hline
À jeun & 0,80 & 1,00 \\
\hline
2 h après le repas & 1,80 & 1,05 \\
\hline
\end{tabularx}
\end{center}
\begin{enumerate}
\item Compare les glycémies des deux veines à jeun, puis après le repas.
\item Que devient le glucose excédentaire au niveau du foie après le repas ? D'où vient le glucose libéré par le foie à jeun ?
\item Déduis-en le rôle du foie dans la régulation de la glycémie.
\end{enumerate}
\end{exercice}

\begin{exercice}[Schéma fonctionnel de la régulation de la glycémie]
Après un repas riche en glucides (attiéké, riz, igname), la glycémie s'élève transitoirement chez un sujet sain.
\begin{enumerate}
\item Réalise et légende le schéma fonctionnel de la régulation hormonale de la glycémie dans cette situation, en faisant apparaître : le pancréas (îlots de Langerhans), l'hormone mise en jeu, les organes cibles (foie, muscles, tissu adipeux) et leurs réponses.
\item Fais apparaître sur ton schéma la boucle de rétrocontrôle (rétroaction négative) qui ramène la glycémie à sa valeur normale.
\item Précise ce qui se passerait, sur le même schéma, en cas de jeûne prolongé.
\end{enumerate}
\end{exercice}

\begin{exercice}[La mesure de la pression artérielle]
Au centre de santé de quartier, l'infirmier mesure la pression artérielle de trois patients au repos, au niveau de l'artère humérale, à l'aide d'un sphygmomanomètre et d'un stéthoscope. Il obtient : patient A : 17/11 ; patient B : 12/8 ; patient C : 8/5 (valeurs en cm de Hg).
\begin{enumerate}
\item Décris, en les justifiant, les étapes de la mesure de la pression artérielle (gonflage du brassard, auscultation des bruits de Korotkoff, dégonflage progressif).
\item À quels événements cardiaques correspondent la pression maximale et la pression minimale ?
\item Interprète les valeurs obtenues pour chaque patient et propose pour chacun une conduite à tenir.
\end{enumerate}
\end{exercice}

\begin{exercice}[Pression artérielle et effort musculaire]
On enregistre la pression artérielle systolique d'un élève avant, pendant et après une course de 400 m.
\begin{center}
\begin{tabularx}{\linewidth}{|l|X|X|X|}
\hline
\rowcolor{popblueL} \textbf{Moment} & \textbf{Avant l'effort} & \textbf{Fin de l'effort} & \textbf{10 min après} \\
\hline
Pression systolique (cm de Hg) & 12 & 18 & 12,5 \\
\hline
Fréquence cardiaque (battements/min) & 70 & 150 & 85 \\
\hline
\end{tabularx}
\end{center}
\begin{enumerate}
\item Analyse les variations de la pression artérielle au cours de l'expérience.
\item Explique ces variations à l'aide des paramètres physiologiques étudiés (débit cardiaque, fréquence cardiaque, diamètre des vaisseaux).
\item Pourquoi la pratique régulière d'une activité physique modérée est-elle recommandée aux hypertendus ?
\end{enumerate}
\end{exercice}

\courssub{Je me prépare au Bac}

\begin{exercice}[Type Bac : les rôles du pancréas et de l'insuline]
Au cours d'une série d'expériences historiques réalisées sur des chiens, on a obtenu les résultats suivants :
\begin{center}
\begin{tabularx}{\linewidth}{|l|X|X|}
\hline
\rowcolor{popblueL} \textbf{Expérience} & \textbf{Protocole} & \textbf{Résultats} \\
\hline
1 & Ablation totale du pancréas & Glycémie élevée (2,5 g/L) ; présence de glucose dans les urines ; soif intense ; mort en quelques semaines \\
\hline
2 & Ligature du canal pancréatique (le suc digestif n'est plus évacué, les îlots de Langerhans restent intacts) & Glycémie normale (1 g/L) ; pas de glucose dans les urines \\
\hline
3 & Injection quotidienne d'extrait d'îlots de Langerhans (insuline) à un chien pancréatectomisé & Glycémie ramenée à 1 g/L ; disparition du glucose urinaire ; survie de l'animal \\
\hline
4 & Injection d'une forte dose d'insuline à un chien sain & Chute brutale de la glycémie (0,4 g/L) ; convulsions ; coma \\
\hline
\end{tabularx}
\end{center}
\begin{enumerate}
\item Analyse les résultats de l'expérience 1. Quel trouble métabolique présente le chien pancréatectomisé ? Justifie ta réponse à partir des symptômes.
\item En comparant les expériences 1 et 2, montre que la fonction digestive (exocrine) du pancréas n'est pas en cause dans ce trouble.
\item Que démontre l'expérience 3 quant à l'origine de la substance régulatrice ?
\item Interprète l'expérience 4 et précise le sens d'action de l'insuline sur la glycémie.
\item Conclus en dégageant, sous forme de schéma fonctionnel légendé, le rôle du pancréas et de l'insuline dans la régulation de la glycémie.
\end{enumerate}
\end{exercice}

\begin{exercice}[Type Bac : l'épreuve d'hyperglycémie provoquée]
À l'hôpital de district de Bafoussam, on réalise une épreuve d'hyperglycémie provoquée chez trois patients à jeun : chacun ingère 75 g de glucose dissous dans l'eau, et on mesure leur glycémie toutes les 30 min pendant 2 h. Les résultats sont consignés ci-dessous (glycémie en g/L).
\begin{center}
\begin{tabularx}{\linewidth}{|l|X|X|X|X|X|}
\hline
\rowcolor{popblueL} \textbf{Temps (min)} & \textbf{0} & \textbf{30} & \textbf{60} & \textbf{90} & \textbf{120} \\
\hline
Courbe A & 0,90 & 1,40 & 1,20 & 1,00 & 0,90 \\
\hline
Courbe B & 1,80 & 2,60 & 2,90 & 2,70 & 2,50 \\
\hline
Courbe C & 1,30 & 2,10 & 2,00 & 1,70 & 1,50 \\
\hline
\end{tabularx}
\end{center}
On sait que l'un des patients est sain, un autre est diabétique de type 1 non traité, et le troisième est diabétique de type 2 sous traitement (antidiabétiques oraux et régime).
\begin{enumerate}
\item Trace sur un même graphique les trois courbes de glycémie en fonction du temps (1 cm pour 15 min en abscisse ; 1 cm pour 0,2 g/L en ordonnée).
\item Analyse la courbe A : comment expliques-tu l'élévation transitoire de la glycémie puis son retour à la normale ?
\item Attribue chaque courbe à un patient, en justifiant rigoureusement ton choix à partir des valeurs à jeun et de l'évolution de la glycémie.
\item Explique, en t'appuyant sur les mécanismes de régulation hormonale, pourquoi la glycémie du diabétique de type 1 ne redescend pas.
\item Le médecin suspecte chez le patient B une glycosurie (présence de glucose dans les urines). Explique l'origine de ce symptôme sachant que le seuil rénal de réabsorption du glucose est d'environ 1,80 g/L.
\end{enumerate}
\end{exercice}

\begin{exercice}[Type Bac : étude de cas — diabète et hypertension]
Monsieur N., 52 ans, commerçant à Douala, consulte pour fatigue, soif permanente et mictions fréquentes. Le bilan médical donne les résultats suivants : glycémie à jeun : 1,65 g/L ; hémoglobine glyquée (HbA1c) : 8,2 \% (normale : inférieure à 6 \%) ; pression artérielle : 16/10 cm de Hg, confirmée à trois reprises ; masse corporelle : 95 kg pour 1,70 m (IMC = 32,9 kg/m\textsuperscript{2}) ; alimentation riche en graisses et en sel ; sédentarité ; antécédents familiaux de diabète.
\begin{enumerate}
\item Interprète les résultats des analyses sanguines de Monsieur N. Quel diagnostic le médecin pose-t-il ? Justifie.
\item Interprète les valeurs de la pression artérielle. De quel trouble supplémentaire souffre ce patient ?
\item À partir des données du dossier, dégage les facteurs de risque communs au diabète de type 2 et à l'hypertension artérielle présents chez ce patient.
\item Explique pourquoi l'association diabète + hypertension est particulièrement dangereuse pour les vaisseaux sanguins, le cœur et les reins.
\item Rédige, sous forme de programme structuré (alimentation, activité physique, suivi médical), un plan de prévention et de prise en charge adapté au contexte camerounais, que tu pourrais proposer lors d'une séance d'éducation à la santé dans un centre de santé.
\end{enumerate}
\end{exercice}

\newpage

\coursec{Corrigés des exercices}

\begin{corrige}[Exercice 1 — QCM : la glycémie]
\begin{enumerate}
\item \textbf{Réponse b.} La glycémie normale à jeun est comprise entre \SI{0,70}{g/L} et \SI{1,10}{g/L} (soit 3,9 à 6,1 mmol/L). En dessous de \SI{0,70}{g/L}, on parle d'hypoglycémie ; au-dessus de \SI{1,26}{g/L} vérifié deux fois, de diabète.
\item \textbf{Réponse c.} L'\cle{insuline} est la seule hormone \cle{hypoglycémiante} : elle fait baisser la glycémie. Le glucagon, l'adrénaline et le cortisol sont au contraire hyperglycémiants.
\item \textbf{Réponse b.} L'insuline est sécrétée par les \cle{cellules $\beta$} des îlots de Langerhans (pancréas endocrine). Les cellules $\alpha$ sécrètent le glucagon.
\item \textbf{Réponse b.} Une glycémie à jeun supérieure à \SI{1,26}{g/L}, confirmée par un second dosage, est le critère biologique du \cle{diabète}.
\end{enumerate}
\end{corrige}

\begin{corrige}[Exercice 2 — Vrai ou faux]
\begin{enumerate}
\item \textbf{Vrai.} Après un repas glucidique, le foie capte le glucose excédentaire et le stocke sous forme de \cle{glycogène} (glycogénogenèse), sous l'action de l'insuline.
\item \textbf{Faux.} Le glucagon est une hormone \cle{hyperglycémiante} : il déclenche la dégradation du glycogène hépatique (glycogénolyse) et la néoglucogenèse, ce qui \emph{élève} la glycémie.
\item \textbf{Vrai.} Le diabète de type 1 (insulino-dépendant) résulte de la destruction \cle{auto-immune} des cellules $\beta$ des îlots de Langerhans ; il en résulte une carence absolue en insuline.
\item \textbf{Faux.} La pression systolique (maximale) correspond à la \cle{contraction} des ventricules (systole) ; c'est la pression diastolique (minimale) qui correspond à la relaxation cardiaque (diastole).
\item \textbf{Vrai.} Au repos, une pression artérielle normale est voisine de 12/8 cm de Hg (soit 120/80 mm de Hg). On parle d'hypertension au-delà de 14/9 cm de Hg de façon permanente.
\end{enumerate}
\end{corrige}

\begin{corrige}[Exercice 3 — Définitions]
\begin{itemize}
\item \textbf{Glycémie} : taux de glucose dissous dans le sang, exprimé en g/L (ou en mmol/L), normalement compris entre 0,70 et \SI{1,10}{g/L} à jeun.
\item \textbf{Diabète} : maladie métabolique chronique caractérisée par une hyperglycémie permanente (glycémie à jeun $>$ \SI{1,26}{g/L}), due à un défaut de sécrétion ou d'action de l'insuline.
\item \textbf{Pression artérielle} : force exercée par le sang sur la paroi des artères ; elle s'exprime par deux valeurs, la pression maximale (systolique) et la pression minimale (diastolique), en cm de Hg.
\item \textbf{Hypertension artérielle} : élévation permanente de la pression artérielle au-delà de 14/9 cm de Hg au repos, confirmée par plusieurs mesures.
\item \textbf{Hypoglycémie} : baisse anormale de la glycémie en dessous de \SI{0,70}{g/L}, pouvant entraîner malaise, convulsions et coma.
\end{itemize}
\end{corrige}

\begin{corrige}[Exercice 4 — Calculs directs]
\begin{enumerate}
\item Conversion en mmol/L :
\[ c = \frac{1,40 \ \text{g/L}}{180 \ \text{g/mol}} = \num{0,00778} \ \text{mol/L} \approx \boxed{7,8 \ \text{mmol/L}} \]
(2 chiffres significatifs). Cette valeur, supérieure à 7 mmol/L ($\approx$ \SI{1,26}{g/L}), confirme un diabète si elle est retrouvée une seconde fois.
\item Masse totale de glucose circulant :
\[ m = 0,95 \ \text{g/L} \times 5 \ \text{L} = \boxed{4,75 \ \text{g} \approx 4,8 \ \text{g}} \]
\item Pression différentielle (pression pulsée) :
\[ \Delta P = P_{\text{max}} - P_{\text{min}} = 16 - 9 = \boxed{7 \ \text{cm de Hg}} \]
\textbf{Commentaire} : la pression différentielle normale est d'environ 4 cm de Hg. Une valeur de 7 cm de Hg est élevée ; elle traduit souvent une rigidité artérielle (artères peu élastiques, athérosclérose) et s'observe fréquemment chez l'hypertendu âgé.
\end{enumerate}
\end{corrige}

\begin{corrige}[Exercice 5 — Le rôle du foie dans la régulation de la glycémie]
\begin{enumerate}
\item \textbf{Comparaison.}
\begin{itemize}
\item À jeun : la glycémie de la veine sus-hépatique (1,00 g/L) est \emph{supérieure} à celle de la veine porte (0,80 g/L) : le sang qui sort du foie est plus riche en glucose que celui qui y entre. Le foie \textbf{libère} du glucose.
\item 2 h après le repas : la glycémie de la veine porte (1,80 g/L) est \emph{très supérieure} à celle de la veine sus-hépatique (1,05 g/L) : le foie \textbf{prélève} le glucose du sang digestif.
\end{itemize}
\item \textbf{Devenir du glucose.} Après le repas, le glucose excédentaire est capté par les hépatocytes et \textbf{stocké sous forme de glycogène} (glycogénogenèse), sous l'action de l'insuline. À jeun, le glucose libéré provient de la \textbf{dégradation du glycogène hépatique} (glycogénolyse, sous l'action du glucagon), complétée par la \textbf{néoglucogenèse} (fabrication de glucose à partir d'acides aminés et de glycérol) lors de jeûnes prolongés.
\item \textbf{Conclusion.} Le foie est un \textbf{réservoir et un régulateur de glucose} : il stocke le glucose en période d'abondance et le restitue au sang en période de jeûne, maintenant ainsi la glycémie constante autour de 1 g/L. C'est l'organe cible essentiel des hormones pancréatiques.
\end{enumerate}
\textbf{Point de vigilance} : ne pas écrire que le foie « fabrique » le glucose après le repas ; il le \emph{stocke}. À jeun, il le \emph{libère} puis le \emph{néo-synthétise}.
\end{corrige}

\begin{corrige}[Exercice 6 — Schéma fonctionnel de la régulation de la glycémie]
\begin{enumerate}
\item \textbf{Schéma fonctionnel (situation post-prandiale).}
\begin{popfigure}
\begin{tikzpicture}[font=\small, node distance=0.9cm,
box/.style={draw=popblue, fill=popblueL, rounded corners, align=center, minimum height=0.8cm, inner sep=4pt},
org/.style={draw=popgreen, fill=popgreenL, rounded corners, align=center, minimum height=0.8cm, inner sep=4pt},
fleche/.style={-{Stealth[length=2.5mm]}, thick, popdark}]
\node[box, align=center] (gly){Hyperglycémie\\ (glycémie $>$ 1,10 g/L)};
\node[box, below=of gly, align=center] (panc){Pancréas : cellules $\beta$\\ des îlots de Langerhans};
\node[box, below=of panc, align=center] (ins){Sécrétion d'\textbf{insuline}\\ (hormone hypoglycémiante)};
\node[org, below left=0.9cm and -0.5cm of ins, align=center] (foie){Foie : glycogénogenèse\\ (stockage du glucose)};
\node[org, below=of ins, align=center] (musc){Muscles : captation\\ et utilisation du glucose};
\node[org, below right=0.9cm and -0.5cm of ins, align=center] (adi){Tissu adipeux :\\ lipogenèse};
\node[box, below=2.2cm of musc, align=center] (norm){Retour de la glycémie\\ à la normale ($\approx$ 1 g/L)};
\draw[fleche] (gly) -- node[right, font=\scriptsize]{stimulation} (panc);
\draw[fleche] (panc) -- (ins);
\draw[fleche] (ins) -- (foie);
\draw[fleche] (ins) -- (musc);
\draw[fleche] (ins) -- (adi);
\draw[fleche] (foie.south) |- (norm.west);
\draw[fleche] (musc) -- (norm);
\draw[fleche] (adi.south) |- (norm.east);
\draw[fleche, poporange, dashed] (norm.east) .. controls +(3,-1) and +(3,1) .. node[right, font=\scriptsize, align=center]{rétroaction négative :\\ baisse de la sécrétion\\ d'insuline} (panc.east);
\end{tikzpicture}
\legende{Schéma fonctionnel de la régulation de la glycémie après un repas glucidique : l'hyperglycémie stimule les cellules $\beta$ du pancréas, qui sécrètent l'insuline ; celle-ci agit sur le foie, les muscles et le tissu adipeux, ce qui ramène la glycémie à sa valeur normale (rétrocontrôle négatif).}
\end{popfigure}
\item \textbf{Boucle de rétrocontrôle.} Le retour de la glycémie à la normale supprime le stimulus initial : les cellules $\beta$ \emph{cessent} de sécréter de l'insuline. C'est une \cle{rétroaction négative} : l'effet (normalisation de la glycémie) inhibe sa propre cause (la sécrétion d'insuline), ce qui stabilise le paramètre régulé.
\item \textbf{Cas du jeûne prolongé.} La glycémie s'abaisse ($<$ 0,80 g/L) : ce sont alors les \textbf{cellules $\alpha$} des îlots qui sont stimulées ; elles sécrètent le \textbf{glucagon} (hormone hyperglycémiante), qui agit sur le foie : glycogénolyse puis néoglucogenèse. Le glucose libéré dans le sang ramène la glycémie à la normale, ce qui arrête la sécrétion de glucagon (là encore, rétroaction négative). Sur le schéma, on remplace donc : cellules $\beta \to$ cellules $\alpha$ ; insuline $\to$ glucagon ; stockage $\to$ libération de glucose par le foie.
\end{enumerate}
\textbf{Point de vigilance} : au Bac, le schéma doit comporter les flèches orientées, le sens d'action de l'hormone (+ ou $-$) et la boucle de rétrocontrôle explicitement identifiée.
\end{corrige}

\begin{corrige}[Exercice 7 — La mesure de la pression artérielle]
\begin{enumerate}
\item \textbf{Étapes de la mesure (méthode auscultatoire).}
\begin{enumerate}
\item Le patient est au repos, assis, le bras à hauteur du cœur ; on place le brassard du sphygmomanomètre autour du bras, au niveau de l'artère humérale.
\item On \textbf{gonfle le brassard} au-delà de la pression systolique : l'artère est totalement comprimée, le sang ne passe plus, \emph{aucun bruit} n'est perçu au stéthoscope placé au pli du coude.
\item On \textbf{dégonfle progressivement} le brassard. Lorsque la pression du brassard devient égale à la pression systolique, le sang passe par à-coups à chaque systole : apparition des \textbf{bruits de Korotkoff} $\Rightarrow$ lecture de la \textbf{pression maximale}.
\item En poursuivant le dégonflage, les bruits s'atténuent puis \textbf{disparaissent} lorsque l'artère n'est plus comprimée et que l'écoulement redevient laminaire : lecture de la \textbf{pression minimale}.
\end{enumerate}
\item \textbf{Événements cardiaques.} La pression maximale correspond à la \textbf{systole} ventriculaire (contraction du cœur, éjection du sang) ; la pression minimale correspond à la \textbf{diastole} (relaxation cardiaque, remplissage des cavités).
\item \textbf{Interprétation.}
\begin{itemize}
\item \textbf{Patient A : 17/11 cm de Hg.} Hypertension artérielle nette ($>$ 14/9). Conduite : confirmer par des mesures répétées, rechercher les causes, prescrire un régime hyposodé, une activité physique, et un traitement antihypertenseur si l'élévation est confirmée.
\item \textbf{Patient B : 12/8 cm de Hg.} Valeurs \textbf{normales} au repos. Conduite : simple surveillance et conseils d'hygiène de vie.
\item \textbf{Patient C : 8/5 cm de Hg.} Hypotension artérielle. Conduite : rechercher une cause (déshydratation, saignement, trouble hormonal), conseiller une hydratation suffisante, se lever lentement ; avis médical si malaises.
\end{itemize}
\end{enumerate}
\end{corrige}

\begin{corrige}[Exercice 8 — Pression artérielle et effort musculaire]
\begin{enumerate}
\item \textbf{Analyse.} La pression systolique passe de 12 cm de Hg au repos à 18 cm de Hg en fin d'effort (augmentation de 50 \%), puis redescend à 12,5 cm de Hg 10 min après l'arrêt. La fréquence cardiaque suit la même évolution : 70 $\to$ 150 $\to$ 85 battements/min. Les deux paramètres varient donc dans le même sens, et l'effort provoque une élévation \emph{transitoire} et \emph{réversible} de la pression artérielle.
\item \textbf{Explication.} Pendant l'effort, les muscles ont des besoins accrus en O$_2$ et en glucose. La \textbf{fréquence cardiaque} et la force de contraction augmentent, donc le \textbf{débit cardiaque} augmente : plus de sang est éjecté par minute, ce qui élève la pression systolique. Par ailleurs, la \textbf{vasodilatation} des artérioles des muscles actifs (augmentation du diamètre des vaisseaux) facilite l'irrigation musculaire, tandis que la vasoconstriction des territoires non prioritaires (tube digestif) redistribue le débit. Après l'effort, fréquence cardiaque et débit reviennent aux valeurs de repos, et la pression se normalise.
\item \textbf{Intérêt de l'activité physique chez l'hypertendu.} L'exercice régulier et modéré entraîne le cœur (qui devient plus efficace à fréquence moindre), améliore l'élasticité des artères, favorise la vasodilatation périphérique, aide à réduire la masse grasse et le stress : tous ces effets contribuent à \emph{abaisser durablement} la pression artérielle de repos et à limiter les complications cardiovasculaires.
\end{enumerate}
\end{corrige}

\begin{corrige}[Exercice 9 — Type Bac : les rôles du pancréas et de l'insuline]
\begin{enumerate}
\item \textbf{Analyse de l'expérience 1.} Après ablation du pancréas, la glycémie s'élève à 2,5 g/L (normale $\approx$ 1 g/L) : c'est une \textbf{hyperglycémie permanente}. Le glucose dépasse le seuil rénal de réabsorption ($\approx$ 1,80 g/L) et apparaît dans les urines (\textbf{glycosurie}) ; l'élimination d'eau avec le glucose provoque une déshydratation, d'où la \textbf{soif intense} (polydipsie) et des urines abondantes (polyurie). L'ensemble de ces signes définit un \textbf{diabète} expérimental, mortel en quelques semaines. Le pancréas est donc indispensable au maintien d'une glycémie normale.
\item \textbf{Comparaison des expériences 1 et 2.} Dans l'expérience 2, la ligature du canal pancréatique supprime l'évacuation du suc digestif (fonction \textbf{exocrine} abolie) mais laisse les îlots de Langerhans intacts : la glycémie reste normale et il n'y a pas de glycosurie. Le diabète de l'expérience 1 n'est donc \emph{pas} dû à la perte de la fonction digestive : il est lié à la destruction de la partie \textbf{endocrine} du pancréas (les îlots de Langerhans).
\item \textbf{Expérience 3.} L'injection d'extrait d'îlots de Langerhans à un chien pancréatectomisé ramène la glycémie à 1 g/L, fait disparaître la glycosurie et assure la survie. Cela démontre que les îlots de Langerhans sécrètent une \textbf{substance hypoglycémiante}, l'\textbf{insuline}, capable de corriger le diabète : le pancréas agit sur la glycémie par voie \emph{sanguine} (sécrétion hormonale), et non par son canal excréteur.
\item \textbf{Expérience 4.} Une forte dose d'insuline chez un chien sain fait chuter la glycémie à 0,4 g/L (\textbf{hypoglycémie}) : le cerveau, privé de glucose, dysfonctionne, d'où convulsions puis coma. L'insuline est donc une hormone \textbf{hypoglycémiante} : elle fait baisser la glycémie en favorisant la captation, l'utilisation et le stockage du glucose par les cellules.
\item \textbf{Conclusion — schéma fonctionnel.}
\begin{popfigure}
\begin{tikzpicture}[font=\small, node distance=0.9cm,
box/.style={draw=popblue, fill=popblueL, rounded corners, align=center, minimum height=0.8cm, inner sep=4pt},
fleche/.style={-{Stealth[length=2.5mm]}, thick, popdark}]
\node[box, align=center] (gly){Glycémie $>$ normale\\ (après un repas)};
\node[box, below=of gly, align=center] (ilots){Îlots de Langerhans\\ (cellules $\beta$ du pancréas)};
\node[box, below=of ilots, align=center] (ins){Insuline déversée\\ dans le sang};
\node[box, below=of ins, align=center] (cibles){Cellules cibles (foie, muscles,\\ tissu adipeux) : captation,\\ utilisation et stockage du glucose};
\node[box, below=of cibles, align=center] (norm){Glycémie ramenée\\ à $\approx$ 1 g/L};
\draw[fleche] (gly) -- (ilots);
\draw[fleche] (ilots) -- (ins);
\draw[fleche] (ins) -- (cibles);
\draw[fleche] (cibles) -- (norm);
\draw[fleche, poporange, dashed] (norm.west) .. controls +(-3,-1) and +(-3,1) .. node[left, font=\scriptsize, align=center]{rétroaction\\ négative} (ilots.west);
\end{tikzpicture}
\legende{Rôle du pancréas et de l'insuline : l'hyperglycémie stimule les îlots de Langerhans ; l'insuline sécrétée provoque la baisse de la glycémie, qui coupe à son tour la sécrétion d'insuline (rétrocontrôle négatif).}
\end{popfigure}
Le pancréas endocrine, par l'insuline, est l'élément régulateur qui maintient la glycémie constante ; sa destruction entraîne le diabète, et l'insuline administrée corrige ce diabète.
\end{enumerate}
\textbf{Barème indicatif (sur 10)} : analyse de l'expérience 1 avec les trois symptômes justifiés (2 pts) ; comparaison 1/2 et conclusion exocrine/endocrine (2 pts) ; conclusion de l'expérience 3 (2 pts) ; interprétation de l'expérience 4 et sens d'action de l'insuline (2 pts) ; schéma fonctionnel légendé avec rétroaction (2 pts).

\textbf{Points de vigilance} : citer les \emph{valeurs chiffrées} dans chaque justification ; distinguer clairement pancréas \emph{exocrine} (digestif) et \emph{endocrine} (hormonal) ; ne pas écrire que l'insuline « détruit » le glucose : elle en favorise la captation, l'utilisation et le stockage.
\end{corrige}

\begin{corrige}[Exercice 10 — Type Bac : l'épreuve d'hyperglycémie provoquée]
\begin{enumerate}
\item \textbf{Graphique.}
\begin{popfigure}
\begin{tikzpicture}
\begin{axis}[width=0.95\linewidth, height=7.5cm,
xlabel={Temps (min)}, ylabel={Glycémie (g/L)},
xmin=0, xmax=125, ymin=0.6, ymax=3.1,
xtick={0,30,60,90,120}, ytick={0.8,1.0,1.26,1.4,1.8,2.0,2.5,3.0},
grid=both, grid style={popline!50}, legend pos=north east,
legend style={font=\scriptsize}]
\addplot[popcurve, color=popgreen, mark=*] coordinates {(0,0.9) (30,1.4) (60,1.2) (90,1.0) (120,0.9)};
\addplot[popcurve, color=popink, mark=square*] coordinates {(0,1.8) (30,2.6) (60,2.9) (90,2.7) (120,2.5)};
\addplot[popcurve, color=poporange, mark=triangle*] coordinates {(0,1.3) (30,2.1) (60,2.0) (90,1.7) (120,1.5)};
\addplot[domain=0:125, dashed, popmuted] {1.26};
\legend{Courbe A, Courbe B, Courbe C, Seuil du diabète (1,26 g/L)}
\end{axis}
\end{tikzpicture}
\legende{Épreuve d'hyperglycémie provoquée (ingestion de 75 g de glucose) chez trois patients : évolution de la glycémie en fonction du temps.}
\end{popfigure}
\item \textbf{Analyse de la courbe A.} La glycémie à jeun est normale (0,90 g/L). Après l'ingestion de glucose, l'absorption intestinale fait transitoirement monter la glycémie (pic de 1,40 g/L à 30 min). Cette hyperglycémie modérée stimule les cellules $\beta$ du pancréas, qui sécrètent de l'insuline ; celle-ci provoque la captation et le stockage du glucose par le foie, les muscles et le tissu adipeux : la glycémie redescend et revient à sa valeur initiale (0,90 g/L) en 2 h. C'est le profil d'une \textbf{régulation efficace}.
\item \textbf{Attribution des courbes.}
\begin{itemize}
\item \textbf{Courbe A $\to$ patient sain} : glycémie à jeun normale (0,90 g/L $<$ 1,10 g/L) et retour à la normale en moins de 2 h.
\item \textbf{Courbe B $\to$ diabétique de type 1 non traité} : glycémie à jeun très élevée (1,80 g/L $>$ 1,26 g/L), pic majeur (2,90 g/L) et \emph{aucune tendance} à la normalisation (encore 2,50 g/L à 120 min) : absence totale d'insuline.
\item \textbf{Courbe C $\to$ diabétique de type 2 sous traitement} : glycémie à jeun anormale mais modérée (1,30 g/L), élévation atténuée (pic 2,10 g/L) et descente progressive (1,50 g/L à 120 min) grâce au traitement (antidiabétiques oraux + régime), sans toutefois atteindre la normale.
\end{itemize}
\item \textbf{Pourquoi la glycémie du diabétique de type 1 ne redescend pas.} Dans le diabète de type 1, les cellules $\beta$ des îlots de Langerhans sont détruites par un processus auto-immun : il n'y a \textbf{plus de sécrétion d'insuline}. L'hyperglycémie provoquée ne déclenche donc aucune réponse hormonale hypoglycémiante : le glucose n'est ni capté ni stocké efficacement par le foie et les muscles, et le foie continue même de libérer du glucose. La boucle de régulation est rompue : la glycémie reste durablement élevée.
\item \textbf{Origine de la glycosurie chez le patient B.} Dans le rein, le glucose filtré par les glomérules est normalement \emph{totalement réabsorbé} par les tubules, tant que la glycémie reste inférieure au \textbf{seuil rénal de réabsorption} ($\approx$ 1,80 g/L). Chez le patient B, la glycémie dépasse largement ce seuil pendant toute l'épreuve (de 1,80 à 2,90 g/L) : les transporteurs tubulaires sont saturés, le glucose excédentaire n'est pas réabsorbé et passe dans les urines : c'est la \textbf{glycosurie}.
\end{enumerate}
\textbf{Barème indicatif (sur 10)} : graphique soigné (axes légendés, échelles respectées, courbes identifiées) (2 pts) ; analyse de la courbe A avec mécanisme hormonal (2 pts) ; attribution justifiée des trois courbes (3 pts) ; explication du défaut de régulation dans le type 1 (2 pts) ; explication de la glycosurie avec le seuil rénal (1 pt).

\textbf{Points de vigilance} : justifier chaque attribution par \emph{deux} arguments (valeur à jeun \emph{et} évolution) ; citer le seuil de 1,26 g/L pour le diagnostic et le seuil rénal de 1,80 g/L pour la glycosurie ; respecter les échelles demandées sur la copie.
\end{corrige}

\begin{corrige}[Exercice 11 — Type Bac : étude de cas — diabète et hypertension]
\begin{enumerate}
\item \textbf{Interprétation des analyses sanguines.} La glycémie à jeun de 1,65 g/L est supérieure au seuil de 1,26 g/L : critère de diabète. L'HbA1c de 8,2 \% (normale $<$ 6 \%) reflète une hyperglycémie \emph{chronique} installée depuis plusieurs mois (l'hémoglobine glyquée est le témoin de la glycémie moyenne sur environ 3 mois). \textbf{Diagnostic : diabète sucré}, très probablement de \textbf{type 2} (âge 52 ans, obésité, sédentarité, antécédents familiaux, début insidieux).
\item \textbf{Pression artérielle.} 16/10 cm de Hg, confirmée à trois reprises : les deux valeurs dépassent les seuils (14/9 cm de Hg). Le patient souffre en outre d'\textbf{hypertension artérielle} permanente.
\item \textbf{Facteurs de risque communs présents chez ce patient.}
\begin{itemize}
\item \textbf{Obésité} : IMC = 32,9 kg/m\textsuperscript{2} (obésité, norme 18,5--25) ;
\item \textbf{alimentation déséquilibrée} : riche en graisses (favorise le diabète de type 2 et l'athérosclérose) et en sel (favorise l'hypertension) ;
\item \textbf{sédentarité} (absence d'activité physique) ;
\item \textbf{antécédents familiaux} (terrain génétique) et âge $>$ 50 ans.
\end{itemize}
\item \textbf{Danger de l'association diabète + hypertension.} L'hyperglycémie chronique altère la paroi des vaisseaux (glycation des protéines, athérosclérose accélérée) et l'hypertension soumet ces parois fragilisées à une pression excessive. Conséquences :
\begin{itemize}
\item \textbf{vaisseaux} : athérosclérose, risque d'infarctus et d'accident vasculaire cérébral (AVC), atteinte des petits vaisseaux (rétine, membres inférieurs) ;
\item \textbf{cœur} : travail accru contre une pression élevée $\Rightarrow$ hypertrophie puis insuffisance cardiaque, coronaropathies ;
\item \textbf{reins} : destruction progressive des glomérules (néphropathie diabétique aggravée par l'HTA) pouvant conduire à l'insuffisance rénale chronique et à la dialyse.
\end{itemize}
L'association multiplie les risques de chaque maladie prise isolément.
\item \textbf{Programme de prévention et de prise en charge (contexte camerounais).}
\begin{itemize}
\item \textbf{Alimentation} : réduire le sel (limiter les cubes d'assaisonnement, les poissons fumés très salés), réduire les graisses (fritures, huile de palme en excès) et les boissons sucrées ; privilégier les légumes locaux (folong, ndolé, gombo), les fruits de saison, les céréales complètes et les légumineuses (niébé, arachides non salées) ; fractionner les repas et contrôler les portions de féculents (riz, igname, attiéké).
\item \textbf{Activité physique} : au moins 30 min de marche rapide par jour (marche au marché, travaux champêtres actifs, football modéré), 5 jours par semaine ; objectif de perte de poids progressive (viser un IMC $<$ 25 kg/m\textsuperscript{2}).
\item \textbf{Suivi médical} : consultation régulière au centre de santé de district ; contrôle mensuel de la pression artérielle et trimestriel de la glycémie et de l'HbA1c ; observance stricte des traitements prescrits (antidiabétiques oraux, antihypertenseurs) ; sevrage tabagique et limitation de l'alcool ; éducation du patient et de sa famille à reconnaître les signes d'alerte (malaise hypoglycémique, maux de tête violents, douleur thoracique).
\end{itemize}
\end{enumerate}
\textbf{Barème indicatif (sur 10)} : interprétation des analyses et diagnostic justifié (2,5 pts) ; interprétation de la PA (1,5 pts) ; facteurs de risque dégagés du dossier (2 pts) ; mécanismes de gravité de l'association (2 pts) ; programme structuré et contextualisé (2 pts).

\textbf{Points de vigilance} : citer les \emph{valeurs seuils} (1,26 g/L ; 6 \% d'HbA1c ; 14/9 cm de Hg ; IMC 25 et 30 kg/m\textsuperscript{2}) ; justifier le type 2 par les éléments du dossier et non par une affirmation gratuite ; le programme doit être \emph{concret, structuré et adapté au contexte local}.
\end{corrige}

\newpage

\coursec{Fiche bilan}

\begin{popfigure}
\begin{tikzpicture}[
  mindmap,
  grow cyclic,
  every node/.style={concept, execute at begin node=\hspace{0pt}},
  root concept/.append style={concept color=popblue, font=\bfseries\small, text=white, minimum size=3.2cm},
  level 1 concept/.append style={level distance=4.6cm, sibling angle=72, font=\scriptsize},
  level 2 concept/.append style={level distance=2.8cm, font=\tiny},
  concept connection/.append style={color=popmuted}
]
\node[root concept]{Santé nutritionnelle : glycémie et pression artérielle}
  child[concept color=poporange]{ node {Glycémie : définition et mesure}
    child { node {$\approx 1$ g/L à jeun} }
    child { node {Varie après un repas} }
  }
  child[concept color=popgreen]{ node {Régulation hormonale}
    child { node {Insuline : hypoglycémiante} }
    child { node {Glucagon : hyperglycémiant} }
    child { node {Foie : stockage / libération du glucose} }
  }
  child[concept color=poppink]{ node {Diabète et hypoglycémie}
    child { node {Diabète type 1 : insulinodépendant} }
    child { node {Diabète type 2 : insulinorésistant} }
    child { node {Hypoglycémie : malaise, coma} }
  }
  child[concept color=popgold]{ node {Pression artérielle : mesure}
    child { node {Sphygmomanomètre} }
    child { node {Maxi / mini : 12--13 / 8 cm Hg} }
  }
  child[concept color=poppurple]{ node {HTA et hypotension}
    child { node {HTA : $>$ 14/9 cm Hg} }
    child { node {Causes : sel, stress, hérédité} }
    child { node {Lutte : hygiène de vie} }
  };
\end{tikzpicture}
\legende{Carte mentale de la leçon : les cinq branches essentielles (glycémie, régulation hormonale, diabète et hypoglycémie, mesure de la pression artérielle, hyper/hypotension).}
\end{popfigure}

\begin{bilan}
\textbf{1. Définitions clés}
\begin{itemize}
  \item \cle{Glycémie} : taux de glucose dissous dans le sang ; valeur normale à jeun comprise entre \SI{0,70}{g/L} et \SI{1,10}{g/L} (environ \SI{1}{g/L}).
  \item \cle{Hyperglycémie} : glycémie durablement supérieure à \SI{1,26}{g/L} à jeun ; \cle{hypoglycémie} : glycémie inférieure à \SI{0,60}{g/L}.
  \item \cle{Pression artérielle} : pression exercée par le sang sur la paroi des artères ; elle comprend la pression \emph{systolique} (maxi, cœur en contraction) et la pression \emph{diastolique} (mini, cœur au repos). Valeur normale : environ $12$--$13 / 8$ cm de Hg.
  \item \cle{Hypertension artérielle (HTA)} : pression durablement supérieure à $14/9$ cm de Hg ; \cle{hypotension} : pression durablement inférieure à $10/6$ cm de Hg.
\end{itemize}

\textbf{2. Hormones et organes de la régulation de la glycémie (à connaître par cœur)}
\begin{center}
\begin{tabularx}{\linewidth}{|l|X|X|X|}
\hline
\rowcolor{popblueL} \textbf{Hormone} & \textbf{Origine} & \textbf{Effet sur la glycémie} & \textbf{Action sur le foie} \\
\hline
Insuline & Cellules $\beta$ des îlots de Langerhans (pancréas) & \textbf{Hypoglycémiante} : seule hormone qui fait baisser la glycémie & Favorise le stockage du glucose en glycogène (glycogénogenèse) et l'entrée du glucose dans les cellules \\
\hline
Glucagon & Cellules $\alpha$ des îlots de Langerhans (pancréas) & \textbf{Hyperglycémiante} & Provoque l'hydrolyse du glycogène hépatique (glycogénolyse) qui libère du glucose \\
\hline
Adrénaline & Glandes surrénales & Hyperglycémiante (stress, effort) & Stimule la glycogénolyse \\
\hline
\end{tabularx}
\end{center}
Le pancréas est la \emph{centrale de détection} (il mesure la glycémie) et le foie est l'\emph{organe effecteur} (il stocke ou libère le glucose) : c'est une régulation par \cle{rétroaction} (boucle négative) qui maintient la glycémie constante.

\textbf{3. Le diabète : typologie}
\begin{itemize}
  \item \textbf{Diabète de type 1} (insulinodépendant) : destruction des cellules $\beta$ $\Rightarrow$ absence d'insuline ; sujet jeune ; traitement par injections d'insuline.
  \item \textbf{Diabète de type 2} (insulinorésistant) : insuline produite mais inefficace ; sujet adulte, souvent associé à l'obésité et à la sédentarité ; lutte par régime alimentaire, exercice physique, antidiabétiques.
  \item Signes communs : hyperglycémie chronique, glucose dans les urines (\cle{glycosurie}), soif intense (polydipsie), urines abondantes (polyurie).
\end{itemize}

\textbf{4. Méthodes express}
\begin{itemize}
  \item \emph{Interpréter une courbe de glycémie} : repérer la valeur à jeun, le pic post-prandial (vers 1 h après le repas, jusqu'à \SI{1,4}{g/L} chez le sujet sain), puis le retour à la normale en 2 h grâce à l'insuline ; chez le diabétique, le pic est plus haut et le retour à la normale ne se fait pas.
  \item \emph{Schéma fonctionnel de régulation} : toujours faire apparaître la boucle : glycémie $\uparrow$ $\rightarrow$ pancréas sécrète insuline $\rightarrow$ foie stocke le glucose $\rightarrow$ glycémie $\downarrow$ (et inversement avec le glucagon).
  \item \emph{Mesure de la pression artérielle} : brassard gonflable (sphygmomanomètre) au bras, stéthoscope sur l'artère humérale ; le 1\ier{} bruit entendu en dégonflant = pression maxi (systolique), la disparition des bruits = pression mini (diastolique).
  \item \emph{Paramètres de la pression artérielle} : $P$ dépend du débit cardiaque, du volume sanguin et du diamètre des artérioles (résistances périphériques) ; elle varie avec l'effort, le stress, l'âge, la position du corps.
  \item \emph{Lutte contre l'HTA} : réduire le sel, pratiquer une activité physique, limiter alcool et tabac, suivre un traitement antihypertenseur ; contre l'hypotension : hydratation, apport sodé adapté, lever lent.
\end{itemize}
\end{bilan}

\begin{aretenir}[Checklist avant le Bac]
\begin{itemize}
  \item[$\square$] Je sais définir la glycémie et citer sa valeur normale à jeun avec son unité (\SI{1}{g/L} environ).
  \item[$\square$] Je sais interpréter une courbe de glycémie avant et après un repas, chez un sujet sain et chez un diabétique.
  \item[$\square$] Je connais les deux hormones pancréatiques (insuline, glucagon), leur origine cellulaire et leur effet sur la glycémie.
  \item[$\square$] Je sais réaliser le schéma fonctionnel de la régulation de la glycémie avec les flèches de rétroaction.
  \item[$\square$] Je sais distinguer le diabète de type 1 et le diabète de type 2 (causes, âge d'apparition, traitement).
  \item[$\square$] Je connais les signes du diabète (glycosurie, polyurie, polydipsie) et les moyens de lutte.
  \item[$\square$] Je sais expliquer les conséquences d'une hypoglycémie (malaise, convulsions, coma) et le geste d'urgence (apport de sucre).
  \item[$\square$] Je sais définir la pression artérielle et distinguer pression systolique (maxi) et diastolique (mini).
  \item[$\square$] Je sais décrire la technique de mesure de la pression artérielle au sphygmomanomètre.
  \item[$\square$] Je connais les paramètres physiologiques qui font varier la pression artérielle (débit cardiaque, volume sanguin, diamètre des vaisseaux).
  \item[$\square$] Je connais les seuils de l'hypertension ($> 14/9$ cm Hg) et de l'hypotension, leurs causes, symptômes et moyens de lutte.
  \item[$\square$] Je rédige mes interprétations en m'appuyant sur les données chiffrées du document, avec les unités.
\end{itemize}
\end{aretenir}

\findecours

\end{document}
